\documentclass[11pt,a4paper]{article}

% Modo artículo de Beamer: compila los mismos archivos de los temas que
% main.tex, pero como documento corrido.  Con `notheorems' Beamer no define
% los entornos matemáticos; se definen más abajo con amsthm, ya que el tema
% de beamerteaching (donde están para las transparencias) no se carga aquí.
\usepackage[notheorems]{beamerarticle}

% Paquetes, datos del documento y macros comunes con main.tex.
% !TEX root = main.tex
% Preámbulo común a las transparencias (main.tex) y a los apuntes
% (apuntes.tex).  Aquí solo van paquetes, datos del documento y macros que
% sirven igual en los dos modos; lo específico de cada uno está en su archivo.

% Compatibilidad con compilaciones mediante pdfLaTeX.
\usepackage[utf8]{inputenc}
\usepackage[T1]{fontenc}
\usepackage[spanish]{babel}

% Matemáticas y diagramas.
\usepackage{amsmath,amssymb,mathtools}
\usepackage{tikz}
\usetikzlibrary{matrix,positioning,decorations.pathmorphing}
\usepackage{tikz-cd}

\title{Álgebra Lineal y Geometría I}
\author{Fernando Muro}
\date{Curso 2026--2027}

% Comandos generales de la asignatura.
\newcommand{\N}{\mathbb{N}}
\newcommand{\Z}{\mathbb{Z}}
\newcommand{\Q}{\mathbb{Q}}
\newcommand{\R}{\mathbb{R}}
\newcommand{\C}{\mathbb{C}}
\newcommand{\K}{\mathbb{K}}
\newcommand{\Mat}{\operatorname{Mat}}
\newcommand{\rg}{\operatorname{rg}}
\newcommand{\im}{\operatorname{im}}
\newcommand{\id}{\operatorname{id}}
\renewcommand{\deg}{\operatorname{grado}}


\usepackage[margin=2.5cm]{geometry}
\usepackage{amsthm}
\usepackage{hyperref}

% En modo artículo las columnas se ponen una tras otra, pero sin cortar el
% párrafo: un dibujo de la segunda columna quedaría pegado al final del texto
% de la primera.  Cada columna pasa a ser un bloque de párrafos propio (y
% \centering dentro de ella solo le afecta a ella).
\renewenvironment{column}[2][]{\par}{\par}

% Colores de la Universidad de Sevilla, con los mismos nombres que en
% main.tex (se usan en los dibujos TikZ).
\definecolor{alertcolor}{RGB}{165,1,51}
\definecolor{secondarycolor}{RGB}{250,186,9}

% --------------------------------------------------------------------------
% Entornos matemáticos (equivalentes a los de beamerthemeTeaching.sty)
% --------------------------------------------------------------------------
\theoremstyle{plain}
\newtheorem{theorem}{Teorema}[section]
\newtheorem{proposition}[theorem]{Proposición}
\newtheorem{lemma}[theorem]{Lema}
\newtheorem{corollary}[theorem]{Corolario}

\theoremstyle{definition}
\newtheorem{definition}[theorem]{Definición}
\newtheorem{example}[theorem]{Ejemplo}
\newtheorem{exercise}[theorem]{Ejercicio}

\newenvironment{solution}[1][Solución]
  {\begin{proof}[#1]}
  {\end{proof}}

\newenvironment{proofpart}[1][\proofname]
  {\begingroup\renewcommand{\qedsymbol}{}\begin{proof}[#1]}
  {\end{proof}\endgroup}

\newenvironment{solutionpart}[1][Solución]
  {\begingroup\renewcommand{\qedsymbol}{}\begin{proof}[#1]}
  {\end{proof}\endgroup}

\newcommand{\proofcontinuationname}{\proofname~(continuación)}
\newcommand{\solutioncontinuationname}{Solución~(continuación)}

\begin{document}

\maketitle
\tableofcontents

\section{Sistemas de ecuaciones lineales, matrices y determinantes}

\begin{frame}{Cuerpos}

Un \alert{cuerpo} $k$ es un conjunto de ``números'' en el que podemos:

\begin{itemize}
    \item sumar,
    \item restar,
    \item multiplicar,
    \item dividir por cualquier número distinto de cero.
\end{itemize}

\medskip

Ejemplos de cuerpos:
\[
\mathbb{Q},
\qquad
\mathbb{R},
\qquad
\mathbb{C}.
\]

\medskip

En cambio, \(\mathbb{N}\) y \(\mathbb{Z}\) no son cuerpos:

\begin{itemize}
    \item en \(\mathbb{N}\) no siempre podemos restar:
    \[
    2-3\notin\mathbb{N};
    \]
    \item en \(\mathbb{Z}\) no siempre podemos dividir por un número no nulo:
    \[
    \frac{1}{2}\notin\mathbb{Z}.
    \]
\end{itemize}

\end{frame}
\begin{frame}{Matrices}

Sea \(k\) un cuerpo. Una \alert{matriz de tamaño \(m\times n\)}
con coeficientes en \(k\) es una disposición rectangular de números:

\[
A=(a_{ij})=
\left(
\begin{array}{ccccc}
a_{11} & a_{12} & \cdots & a_{1j} & \cdots \\[-1mm]
\vdots & \vdots &        & \vdots &        \\
a_{i1} & a_{i2} & \cdots & a_{ij} & \cdots \\
\vdots & \vdots &        & \vdots &        \\
a_{m1} & a_{m2} & \cdots & a_{mj} & \cdots
\end{array}
\right)
\in M_{m\times n}(k).
\]

\begin{center}
\begin{tikzpicture}[
    every node/.style={font=\small},
    >=stealth
]

% Matriz esquemática
\matrix (A) [
    matrix of math nodes,
    left delimiter={(},
    right delimiter={)},
    row sep=3mm,
    column sep=5mm
] {
 a_{11} & \cdots & a_{1j} & \cdots & a_{1n} \\
 \vdots &        & \vdots &        & \vdots \\
 a_{i1} & \cdots & a_{ij} & \cdots & a_{in} \\
 \vdots &        & \vdots &        & \vdots \\
 a_{m1} & \cdots & a_{mj} & \cdots & a_{mn} \\
};

% Resaltar fila i
\draw[rounded corners]
  ([xshift=-2mm,yshift=2mm]A-3-1.north west)
  rectangle
  ([xshift=2mm,yshift=-2mm]A-3-5.south east);

\node[left=8mm of A-3-1] (fila) {\alert{fila \(i\)}};
\draw[->] (fila) -- (A-3-1.west);

% Resaltar columna j
\draw[rounded corners]
  ([xshift=-2mm,yshift=2mm]A-1-3.north west)
  rectangle
  ([xshift=2mm,yshift=-2mm]A-5-3.south east);

\node[above=7mm of A-1-3] (columna) {\alert{columna \(j\)}};
\draw[->] (columna) -- (A-1-3.north);

% Entrada (i,j)
\node[right=13mm of A-3-5] (entrada)
  {\(\displaystyle a_{ij}\): entrada \((i,j)\)};
\draw[->] (entrada.west) -- (A-3-3.east);

% Tamaño
\node[below=7mm of A] {
  \(m\) filas \(\qquad\) \(n\) columnas
};

\end{tikzpicture}
\end{center}

\end{frame}
\begin{frame}{Operaciones elementales por filas: intercambio}

La \alert{operación elemental tipo 1}
\[
\alert{F_i\leftrightarrow F_j}
\]
intercambia las filas \(i\) y \(j\):

\medskip

\[
\begin{pmatrix}
\vdots&&\vdots\\
a_{i1} & \cdots & a_{in}\\
\vdots &&\vdots\\
a_{j1} & \cdots & a_{jn}\\
\vdots &&\vdots
\end{pmatrix}
\quad
\xrightarrow{\;F_i\leftrightarrow F_j\;}
\quad
\begin{pmatrix}
\vdots &&\vdots\\
a_{j1} & \cdots & a_{jn}\\
\vdots &&\vdots\\
a_{i1} & \cdots & a_{in}\\
\vdots &&\vdots
\end{pmatrix}.
\]

\end{frame}


\begin{frame}{Operaciones elementales por filas: producto por un escalar}

La \alert{operación elemental tipo 2}
\[
\alert{\lambda F_i},
\qquad \lambda\neq0,
\]
multiplica la fila \(i\) por el escalar no nulo \(\lambda\in k\):

\medskip

\[
\begin{pmatrix}
\vdots &&\vdots\\
a_{i1} & \cdots & a_{in}\\
\vdots &&\vdots
\end{pmatrix}
\quad
\xrightarrow{\;\lambda F_i\;}
\quad
\begin{pmatrix}
\vdots &&\vdots\\
\lambda a_{i1} & \cdots & \lambda a_{in}\\
\vdots &&\vdots
\end{pmatrix}.
\]

\end{frame}

\begin{frame}{Operaciones elementales por filas: suma de un múltiplo}

La \alert{operación elemental tipo 3}
\[
\alert{F_i+\lambda F_j},
\qquad i\neq j,
\]
sustituye la fila \(i\) por \(F_i+\lambda F_j\), donde $\lambda\in k$ es un escalar:

\medskip

\[
\begin{pmatrix}
\vdots &&\vdots\\
a_{i1} & \cdots & a_{in}\\
\vdots &&\vdots\\
a_{j1} & \cdots & a_{jn}\\
\vdots &&\vdots
\end{pmatrix}
\quad
\xrightarrow{\;F_i+\lambda F_j\;}
\quad
\begin{pmatrix}
\vdots &&\vdots\\
a_{i1}+\lambda a_{j1}
&
\cdots
&
a_{in}+\lambda a_{jn}\\
\vdots &&\vdots\\
a_{j1} & \cdots & a_{jn}\\
\vdots &&\vdots
\end{pmatrix}.
\]

\end{frame}

\begin{frame}{Equivalencia por filas}

\begin{definition}
Sean \(A,B\in M_{m\times n}(k)\). Diremos que \(A\) y \(B\) son
\alert{equivalentes por filas}, y escribiremos
\[
A\sim_f B,
\]
si \(B\) se obtiene a partir de \(A\) mediante una sucesión finita de
operaciones elementales por filas.
\end{definition}

\end{frame}

\begin{frame}{Matrices escalonadas}

\begin{definition}
Sea \(A\in M_{m\times n}(k)\). El \alert{pivote} de una fila no nula
es su primera entrada no nula.

La matriz \(A\) se dice \alert{escalonada por filas} si:

\begin{enumerate}
    \item todas las filas nulas, si las hay, están situadas después
    de las filas no nulas;

    \item el pivote de cada fila no nula está situado estrictamente
    a la derecha del pivote de la fila anterior.
\end{enumerate}

\[
\begin{pmatrix}
0    &\alert{\ast} & \ast & \ast & \ast & \ast & \ast & \cdots & \ast \\
0    & 0    & 0    & \alert{\ast} & \ast & \ast & \ast & \cdots & \ast \\
0    & 0    & 0    & 0    & \alert{\ast} & \ast & \ast & \cdots & \ast \\
\vdots &&&&&&&&\vdots \\
0    & 0    & 0    & 0    & 0    & 0    & 0    & \cdots & 0
\end{pmatrix}.
\]

\end{definition}

\end{frame}

\begin{frame}{Una matriz que no está escalonada}

Consideremos la matriz
\[
A=
\begin{pmatrix}
\alert{2} & 1 & -1 & 3 & 0\\
0 & 0 & \alert{1} & 2 & 4\\
0 & \alert{3} & 0 & 1 & 2\\
0 & \alert{5} & 2 & 0 & 1
\end{pmatrix}.
\]

Los pivotes están situados, respectivamente, en las columnas
\[
1,\qquad 3,\qquad 2,\qquad 2.
\]

\medskip

La matriz \(A\) \alert{no está escalonada por filas} porque:

\begin{itemize}
    \item el pivote de la tercera fila está situado a la izquierda
    del pivote de la segunda fila;

    \item el pivote de la cuarta fila está situado debajo del pivote
    de la tercera fila.
\end{itemize}

Cualquiera de estas dos condiciones es suficiente para que \(A\) no esté escalonada.

\end{frame}

\begin{frame}[t,allowframebreaks]{Eliminación gaussiana}

\begin{theorem}
Toda matriz $A_{m\times n}$ es equivalente por filas a una matriz escalonada.
\end{theorem}

\begin{proof}
Si \(A=0\), ya es escalonada. Supongamos \(A\neq0\), y sea \(j\)
la primera columna de \(A\) que contiene alguna entrada no nula.

\medskip

Mediante una operación elemental de tipo \(1\), si es necesario,
podemos situar una entrada no nula \(a_{1j}\) en la primera fila.

\medskip

Mediante operaciones elementales de tipo \(3\),
\[
F_i-\frac{a_{ij}}{a_{1j}}F_1,
\qquad i=2,\ldots,m,
\]
anulamos todas las entradas situadas debajo de \(a_{1j}\).

Obtenemos así una matriz por bloques de la forma
\[
\left(
\begin{array}{cccc|ccc}
0&\cdots&0&a_{1j}& * & \cdots & *\\
\hline
0&\cdots&0&0&\\
\vdots&&\vdots&\vdots&&B&\\
0&\cdots&0&0&
\end{array}
\right).
\]

Aplicamos el mismo procedimiento a la submatriz \(B\).

\medskip

Repitiendo sucesivamente:
\begin{enumerate}
    \item localizamos la primera columna no nula de la submatriz restante;
    \item mediante una operación de tipo \(1\), situamos una entrada
    no nula en su primera fila;
    \item mediante operaciones de tipo \(3\), anulamos todas las
    entradas situadas debajo de ella;
\end{enumerate}
y obtenemos, después de un número finito de pasos, una matriz escalonada.

\medskip

Cada paso consiste únicamente en operaciones elementales por filas;
por tanto, la matriz obtenida es equivalente por filas a \(A\).
\end{proof}

Este procedimiento se denomina \alert{método de eliminación gaussiana} o \alert{método de Gauss}.
\end{frame}

\begin{frame}[allowframebreaks]{Cálculo de una forma escalonada}

\[
\begin{pmatrix}
0&0&2&1&-1\\
2&4&2&0&8\\
0&0&0&2&-5\\
4&8&8&4&9
\end{pmatrix}
\xrightarrow{F_1\leftrightarrow F_2}
\begin{pmatrix}
2&4&2&0&8\\
0&0&2&1&-1\\
0&0&0&2&-5\\
4&8&8&4&9
\end{pmatrix}
\]

\[
\xrightarrow{F_4-2F_1}
\begin{pmatrix}
2&4&2&0&8\\
0&0&2&1&-1\\
0&0&0&2&-5\\
0&0&4&4&-7
\end{pmatrix}
\xrightarrow{F_4-2F_2}
\begin{pmatrix}
2&4&2&0&8\\
0&0&2&1&-1\\
0&0&0&2&-5\\
0&0&0&2&-5
\end{pmatrix}
\]

\[
\xrightarrow{F_4-F_3}
\begin{pmatrix}
2&4&2&0&8\\
0&0&2&1&-1\\
0&0&0&2&-5\\
0&0&0&0&0
\end{pmatrix}.
\]

La última matriz está escalonada por filas.

\end{frame}

\begin{frame}{Matrices escalonadas y reducidas}

\begin{definition}

Una matriz escalonada por filas se dice \alert{reducida} si, además:

\begin{enumerate}

    \item todos sus pivotes son iguales a \(1\);

    \item cada pivote es la única entrada no nula de su columna.
\end{enumerate}

\[
\begin{pmatrix}
0    & \alert{1}    & \ast & \alert{0}    & \alert{0}    & \ast & \ast & \cdots & \ast \\
0    & 0    & 0    & \alert{1}    & \alert{0}    & \ast & \ast & \cdots & \ast \\
0    & 0    & 0    & 0    & \alert{1}    & \ast & \ast & \cdots & \ast \\
\vdots &&&&&&&&\vdots \\
0    & 0    & 0    & 0    & 0    & 0    & 0    & \cdots & 0
\end{pmatrix}.
\]

\end{definition}

\end{frame}

\begin{frame}[allowframebreaks]{Forma escalonada reducida}

\begin{theorem}
Toda matriz es equivalente por filas a una matriz
escalonada y reducida.
\end{theorem}

\begin{proof}
Sea \(A\in M_{m\times n}(k)\). Sabemos que \(A\) es equivalente
por filas a una matriz escalonada \(E\).

Multiplicando cada fila no nula de \(E\) por el inverso de su pivote,
mediante operaciones elementales de tipo \(2\), podemos suponer que
todos los pivotes son iguales a \(1\).

A continuación, empezando por el último pivote y procediendo hacia
arriba, mediante operaciones elementales de tipo \(3\) anulamos
todas las entradas situadas por encima de cada pivote.

La matriz obtenida es escalonada y reducida y es equivalente por filas
a \(A\). Esto demuestra la existencia.
\end{proof}

\end{frame}

\begin{frame}[allowframebreaks]{De la forma escalonada a la reducida}

\[
\begin{pmatrix}
2&4&2&0&8\\
0&0&2&1&-1\\
0&0&0&2&-5\\
0&0&0&0&0
\end{pmatrix}
\xrightarrow{\substack{
\frac12F_1\\
\frac12F_2\\
\frac12F_3
}}
\begin{pmatrix}
1&2&1&0&4\\
0&0&1&\frac12&-\frac12\\
0&0&0&1&-\frac52\\
0&0&0&0&0
\end{pmatrix}
\]

\[
\xrightarrow{
F_2-\frac12F_3
}
\begin{pmatrix}
1&2&1&0&4\\
0&0&1&0&\frac34\\
0&0&0&1&-\frac52\\
0&0&0&0&0
\end{pmatrix}
\xrightarrow{
F_1-F_2
}
\begin{pmatrix}
1&2&0&0&\frac{13}{4}\\
0&0&1&0&\frac34\\
0&0&0&1&-\frac52\\
0&0&0&0&0
\end{pmatrix}.
\]

La última matriz está escalonada y reducida.

\end{frame}

\begin{frame}{Operaciones elementales inversas}

Toda operación elemental por filas puede \alert{deshacerse}
mediante otra operación elemental.

\medskip

\begin{enumerate}

\item Intercambiar dos filas:
\[
A
\xrightarrow{F_i\leftrightarrow F_j}
B
\xrightarrow{F_i\leftrightarrow F_j}
A.
\]

\item Multiplicar una fila por \(\lambda\neq0\):
\[
A
\xrightarrow{\lambda F_i}
B
\xrightarrow{\lambda^{-1}F_i}
A.
\]

\item Sumar a una fila un múltiplo de otra:
\[
A
\xrightarrow{F_i+\lambda F_j}
B
\xrightarrow{F_i-\lambda F_j}
A,
\qquad i\neq j.
\]

\end{enumerate}

\medskip

En cada caso, la segunda operación se denomina
\alert{operación elemental inversa} de la primera.

\end{frame}

\begin{frame}{Matrices elementales}

\begin{definition}
Una \alert{matriz elemental} es una matriz que se obtiene de la
matriz identidad \(I_m\) mediante una única operación elemental
por filas.
\end{definition}

\medskip

Hay, por tanto, tres tipos de matrices elementales:

\begin{enumerate}
    \item \(E_{ij}\): se obtiene intercambiando las filas \(i\) y \(j\)
    de \(I_m\).

    \item \(E_i(\lambda)\), \(\lambda\neq0\): se obtiene multiplicando
    la fila \(i\) de \(I_m\) por \(\lambda\).

    \item \(E_{ij}(\lambda)\), \(i\neq j\): se obtiene sumando a la
    fila \(i\) de \(I_m\) la fila \(j\) multiplicada por \(\lambda\).
\end{enumerate}

\end{frame}

\begin{frame}{Matrices elementales de tipo 1}

La matriz elemental correspondiente a
\(F_i\leftrightarrow F_j\) es

\[
E_{ij}=
\begin{pmatrix}
1      & 0      & \cdots & 0      & \cdots & 0      & \cdots & 0\\
0      & 1      & \ddots & \vdots &        & \vdots &        & \vdots\\
\vdots & \ddots & \ddots & 0      &        & 0      &        & \vdots\\
0      & \cdots & 0      & \alert{0} & \cdots & \alert{1} & \cdots & 0\\
\vdots &        &        & \vdots & \ddots & \vdots &        & \vdots\\
0      & \cdots & 0      & \alert{1} & \cdots & \alert{0} & \cdots & 0\\
\vdots &        &        & \vdots &        & \vdots & \ddots & 0\\
0      & \cdots & \cdots & 0      & \cdots & 0      & 0      & 1
\end{pmatrix}.
\]

\[
\hspace{34mm}
\underset{i}{\uparrow}
\hspace{19mm}
\underset{j}{\uparrow}
\]

Se obtiene de \(I_m\) intercambiando las filas \(i\) y \(j\).

\end{frame}

\begin{frame}{Matrices elementales de tipo 2}

La matriz elemental correspondiente a
\(\lambda F_i\), con \(\lambda\neq0\), es

\[
E_i(\lambda)=
\begin{pmatrix}
1      & 0      & \cdots & 0               & \cdots & 0\\
0      & 1      & \ddots & \vdots          &        & \vdots\\
\vdots & \ddots & \ddots & 0               &        & \vdots\\
0      & \cdots & 0      & \alert{\lambda} & \cdots & 0\\
\vdots &        &        & \vdots          & \ddots & 0\\
0      & \cdots & \cdots & 0               & 0      & 1
\end{pmatrix}.
\]

\[
\hspace{44mm}
\underset{i}{\uparrow}
\]

Se obtiene de \(I_m\) multiplicando la fila \(i\) por \(\lambda\).

\end{frame}

\begin{frame}{Matrices elementales de tipo 3}

La matriz elemental correspondiente a
\(F_i+\lambda F_j\), con \(i\neq j\), es

\[
E_{ij}(\lambda)=
\begin{pmatrix}
1      & 0      & \cdots & 0      & \cdots & 0               & \cdots & 0\\
0      & 1      & \ddots & \vdots &        & \vdots          &        & \vdots\\
\vdots & \ddots & \ddots & 0      &        & 0               &        & \vdots\\
0      & \cdots & 0      & 1      & \cdots & \alert{\lambda} & \cdots & 0\\
\vdots &        &        & \vdots & \ddots & \vdots          &        & \vdots\\
0      & \cdots & 0      & 0      & \cdots & 1               & \cdots & 0\\
\vdots &        &        & \vdots &        & \vdots          & \ddots & 0\\
0      & \cdots & \cdots & 0      & \cdots & 0               & 0      & 1
\end{pmatrix}.
\]

\[
\hspace{34mm}
\underset{i}{\uparrow}
\hspace{19mm}
\underset{j}{\uparrow}
\]

Se obtiene de \(I_m\) sumando a la fila \(i\) la fila \(j\)
multiplicada por \(\lambda\).

\end{frame}

\begin{frame}{Matrices elementales y operaciones por filas}

\begin{proposition}
Sea \(A\in M_{m\times n}(k)\). Entonces:

\begin{enumerate}

\item La operación de tipo \(1\) viene dada por
\[
A
\xrightarrow{F_i\leftrightarrow F_j}
E_{ij}A.
\]

\item La operación de tipo \(2\), con \(\lambda\neq0\), viene dada por
\[
A
\xrightarrow{\lambda F_i}
E_i(\lambda)A.
\]

\item La operación de tipo \(3\), con \(i\neq j\), viene dada por
\[
A
\xrightarrow{F_i+\lambda F_j}
E_{ij}(\lambda)A.
\]

\end{enumerate}
\end{proposition}


\end{frame}

\begin{frame}{Matrices invertibles}

\begin{definition}
Sea \(A\in M_{m\times n}(k)\). Diremos que \(A\) es
\alert{invertible} si existe una matriz \(B\in M_{n\times m}(k)\)
tal que
\[
AB=I_m
\qquad\text{y}\qquad
BA=I_n.
\]

En tal caso, \(B\) se denomina una \alert{inversa} de \(A\).
\end{definition}

\end{frame}

\begin{frame}{Unicidad de la inversa}

\begin{proposition}
Sea \(A\in M_{m\times n}(k)\) una matriz invertible. Entonces su
inversa es única.
\end{proposition}

\begin{proof}
Sean \(B,C\in M_{n\times m}(k)\) dos inversas de \(A\). Entonces
\[
B
=BI_m
=B(AC)
=(BA)C
=I_nC
=C.
\]
\end{proof}

Por tanto, podemos denotar la única \alert{inversa de \(A\)} por
\alert{\(A^{-1}.\)}

\end{frame}

\begin{frame}{Inversa de la inversa}

\begin{proposition}
Sea \(A\in M_{m\times n}(k)\) una matriz invertible. Entonces
\(A^{-1}\) es también invertible y
\[
\left(A^{-1}\right)^{-1}=A.
\]
\end{proposition}

\begin{proof}
Por definición de matriz inversa,
\[
AA^{-1}=I_m
\qquad\text{y}\qquad
A^{-1}A=I_n.
\]

Estas mismas igualdades muestran que \(A\) es una inversa de
\(A^{-1}\). Por tanto, \(A^{-1}\) es invertible.

Como la inversa es única,
\[
\left(A^{-1}\right)^{-1}=A.
\]
\end{proof}

\end{frame}


\begin{frame}{Inversa de una matriz elemental}

\begin{proposition}
Toda matriz elemental es invertible. Más concretamente:

\begin{enumerate}
    \item \(E_{ij}^{-1}=E_{ij}\).

    \item
    \(E_i(\lambda)^{-1}=E_i(\lambda^{-1})\) para \(\lambda\neq0\).

    \item 
    \(E_{ij}(\lambda)^{-1}=E_{ij}(-\lambda)\) para \(i\neq j\).
\end{enumerate}
\end{proposition}

\begin{proof}
En cada caso, la matriz indicada como inversa corresponde a la
operación elemental inversa. Por tanto, al multiplicar ambas matrices
en cualquiera de los dos órdenes se obtiene \(I_m\).
\end{proof}

\end{frame}

\begin{frame}{Producto de matrices invertibles}

\begin{proposition}
Sean \(A\in M_{m\times n}(k)\) y \(B\in M_{n\times p}(k)\)
matrices invertibles. Entonces \(AB\) es invertible y
\[
(AB)^{-1}=B^{-1}A^{-1}.
\]
\end{proposition}
\begin{proof}
Por asociatividad del producto,
\[
(AB)(B^{-1}A^{-1})
=A(BB^{-1})A^{-1}
=AI_nA^{-1}
=AA^{-1}
=I_m.
\]
Por otra parte,
\[
(B^{-1}A^{-1})(AB)
=B^{-1}(A^{-1}A)B
=B^{-1}I_nB
=B^{-1}B
=I_p.
\]
Por tanto, \(AB\) es invertible y
\[
(AB)^{-1}=B^{-1}A^{-1}.\qedhere
\]
\end{proof}

\end{frame}

\begin{frame}{Forma escalonada y matrices elementales}

\begin{corollary}
Sea \(A\in M_{m\times n}(k)\) y sea \(B\) una forma escalonada por filas
de \(A\). Entonces existe una matriz invertible \(P\in M_m(k)\),
producto de matrices elementales, tal que
\[
B=PA.
\]
\end{corollary}

\vspace{-20pt}

\begin{proof}
Como \(B\) se obtiene de \(A\) mediante una sucesión finita de
operaciones elementales por filas, existen matrices elementales
\[
E_1,\ldots,E_r\in M_m(k)
\]
tales que
\[
B=E_r\cdots E_1A.
\]

Como cada \(E_i\) es invertible y el producto de matrices invertibles
es invertible, obtenemos la ecuación deseada tomando
\[
P=E_r\cdots E_1,\qquad
B=PA.\qedhere
\]
\end{proof}

\end{frame}


\begin{frame}{La eliminación gaussiana mediante matrices elementales}

En el ejemplo anterior hemos obtenido
\[
A= \begin{pmatrix}
0&0&2&1&-1\\
2&4&2&0&8\\
0&0&0&2&-5\\
4&8&8&4&9
\end{pmatrix}
\xrightarrow{F_1\leftrightarrow F_2}
\xrightarrow{F_4-2F_1}
\xrightarrow{F_4-2F_2}
\xrightarrow{F_4-F_3}
\begin{pmatrix}
2&4&2&0&8\\
0&0&2&1&-1\\
0&0&0&2&-5\\
0&0&0&0&0
\end{pmatrix}=B.
\]

Las matrices elementales correspondientes a las operaciones elementales anteriores son, respectivamente
\[
E_{12},\qquad
E_{41}(-2),\qquad
E_{42}(-2),\qquad
E_{43}(-1).
\]

Por tanto, si
\[
P=
E_{43}(-1)E_{42}(-2)E_{41}(-2)E_{12},
\]
entonces
\[
B=PA.
\]

\end{frame}

\begin{frame}{Cálculo de la matriz \(P\)}

Para calcular
\[
P=E_{43}(-1)E_{42}(-2)E_{41}(-2)E_{12},
\]
realizamos sobre \(I_4\) las mismas operaciones elementales:

\[
I_4=
\begin{pmatrix}
1&0&0&0\\
0&1&0&0\\
0&0&1&0\\
0&0&0&1
\end{pmatrix}
\xrightarrow{F_1\leftrightarrow F_2}
\begin{pmatrix}
0&1&0&0\\
1&0&0&0\\
0&0&1&0\\
0&0&0&1
\end{pmatrix}
\xrightarrow{F_4-2F_1}
\begin{pmatrix}
0&1&0&0\\
1&0&0&0\\
0&0&1&0\\
0&-2&0&1
\end{pmatrix}\]

\[\xrightarrow{F_4-2F_2}
\begin{pmatrix}
0&1&0&0\\
1&0&0&0\\
0&0&1&0\\
-2&-2&0&1
\end{pmatrix}
\xrightarrow{F_4-F_3}
\begin{pmatrix}
0&1&0&0\\
1&0&0&0\\
0&0&1&0\\
-2&-2&-1&1
\end{pmatrix}=P.
\]


\end{frame}

\begin{frame}{Cálculo de la inversa de \(P\)}

La matriz
\[
P=
E_{43}(-1)E_{42}(-2)E_{41}(-2)E_{12}
\]
es invertible pues es producto de elementales, que son invertibles. Su inversa es
\begin{align*}
P^{-1}&=
(E_{43}(-1)E_{42}(-2)E_{41}(-2)E_{12})^{-1}\\
&= E_{12}^{-1} E_{41}(-2)^{-1} E_{42}(-2)^{-1} E_{43}(-1)^{-1}\\
&= E_{12} E_{41}(2) E_{42}(2) E_{43}(1)
\end{align*}

\end{frame}

\begin{frame}{Cálculo de la matriz \(P^{-1}\)}

Para calcular \(P^{-1}\), realizamos sobre \(I_4\) las operaciones
elementales inversas en orden contrario:

\[
I_4=
\begin{pmatrix}
1&0&0&0\\
0&1&0&0\\
0&0&1&0\\
0&0&0&1
\end{pmatrix}
\xrightarrow{F_4+F_3}
\begin{pmatrix}
1&0&0&0\\
0&1&0&0\\
0&0&1&0\\
0&0&1&1
\end{pmatrix}
\]

\[
\xrightarrow{F_4+2F_2}
\begin{pmatrix}
1&0&0&0\\
0&1&0&0\\
0&0&1&0\\
0&2&1&1
\end{pmatrix}
\xrightarrow{F_4+2F_1}
\begin{pmatrix}
1&0&0&0\\
0&1&0&0\\
0&0&1&0\\
2&2&1&1
\end{pmatrix}
\xrightarrow{F_1\leftrightarrow F_2}
\begin{pmatrix}
0&1&0&0\\
1&0&0&0\\
0&0&1&0\\
2&2&1&1
\end{pmatrix}
={P^{-1}}.
\]


\end{frame}

\begin{frame}{Vectores}

\begin{definition}
Un \alert{vector} de \(k^n\) es una matriz columna de \(n\) entradas:
\[
v=
\begin{pmatrix}
v_1\\
v_2\\
\vdots\\
v_n
\end{pmatrix},
\qquad v_1,\ldots,v_n\in k.
\]
Por tanto identificamos
\[
k^n=M_{n\times 1}(k).
\]
\end{definition}

\end{frame}

\begin{frame}{Sistemas de ecuaciones lineales}

\begin{definition}
Un \alert{sistema de \(m\) ecuaciones lineales con \(n\) incógnitas}
\(x_1,\ldots,x_n\) es una expresión de la forma
\[
\left\{
\begin{aligned}
a_{11}x_1+\cdots+a_{1n}x_n&=b_1,\\
a_{21}x_1+\cdots+a_{2n}x_n&=b_2,\\
&\ \vdots\\
a_{m1}x_1+\cdots+a_{mn}x_n&=b_m,
\end{aligned}
\right.
\]
donde \(a_{ij},b_i\in k\).

\medskip

Una \alert{solución} del sistema es una sucesión de escalares $v_1,\dots,v_n\in k$ tales que, sustituyendo $x_i=v_i$, se satisfacen simultáneamente todas las ecuaciones.
\end{definition}

\end{frame}

\begin{frame}{Forma matricial de un sistema}

Definimos
\[
\begin{array}{ccc}
A=
\begin{pmatrix}
a_{11}&\cdots&a_{1n}\\
\vdots&&\vdots\\
a_{m1}&\cdots&a_{mn}
\end{pmatrix},
&
x=
\begin{pmatrix}
x_1\\
\vdots\\
x_n
\end{pmatrix},
&
b=
\begin{pmatrix}
b_1\\
\vdots\\
b_m
\end{pmatrix}.\\[8mm]
\text{\alert{matriz de coeficientes}}&
\text{\alert{vector de incógnitas}}&
\text{\alert{vector de términos independientes}}
\end{array}
\]

Entonces el sistema puede escribirse en la forma
\[
\alert{Ax=b},
\]
y una \alert{solución} es un vector
\[v=
\begin{pmatrix}
v_1\\
\vdots\\
v_n
\end{pmatrix},\qquad 
Av=b.
\]

\end{frame}

\begin{frame}{Matriz ampliada}

La \alert{matriz ampliada} del sistema es
\[
(A\mid b)=
\left(
\begin{array}{ccc|c}
a_{11}&\cdots&a_{1n}&b_1\\
\vdots&&\vdots&\vdots\\
a_{m1}&\cdots&a_{mn}&b_m
\end{array}
\right).
\]


Cada fila de \((A\mid b)\) representa una de las ecuaciones del sistema.

\end{frame}

\begin{frame}{Carácter de un sistema}

\begin{definition}
Un sistema de ecuaciones lineales se denomina:
\begin{itemize}
    \item \alert{compatible} si tiene alguna solución;
    \item \alert{incompatible} si no tiene ninguna solución.
\end{itemize}

Un sistema compatible se denomina:
\begin{itemize}
    \item \alert{compatible determinado} si tiene una única solución;
    \item \alert{compatible indeterminado} si tiene más de una solución.
\end{itemize}
\end{definition}

\end{frame}

\begin{frame}[allowframebreaks, t]{Sistemas equivalentes}

\begin{definition}
Dos sistemas de ecuaciones lineales con las mismas incógnitas se dicen
\alert{equivalentes} si tienen las mismas soluciones.
\end{definition}

\begin{proposition}
Dos sistemas cuyas matrices ampliadas son equivalentes por filas
son equivalentes.
\end{proposition}

\begin{proof}
Sean
\[
Ax=b,
\qquad
A'x=b'
\]
los dos sistemas. Como sus matrices ampliadas son equivalentes por
filas, existe una matriz invertible \(P\), producto de matrices
elementales, tal que
\[
(A'\mid b')=P(A\mid b)=(PA\mid Pb).
\]

\framebreak

Por tanto,
\[
A'=PA,
\qquad
b'=Pb.
\]

Sea \(v\in k^n\). Entonces
\[
\begin{aligned}
v\text{ es solución de }Ax=b
&\iff Av=b\\
&\iff PAv=Pb\\
&\iff A'v=b'\\
&\iff v\text{ es solución de }A'x=b'.
\end{aligned}
\]
La segunda equivalencia se debe a que \(P\) es invertible.

\medskip

Así, un vector es solución del primer sistema si y solo si es
solución del segundo. Por definición, los dos sistemas son
equivalentes.
\end{proof}

\end{frame}

\begin{frame}{Sistemas en forma escalonada y reducida}

Supongamos que la matriz ampliada \((A\mid b)\) de un sistema
está escalonada y reducida.

\begin{itemize}
\item Si la última columna contiene un pivote, una de las filas es
\[
\left(
\begin{array}{cccc|c}
0&\cdots&0&0&1
\end{array}
\right),
\]
que representa la ecuación
\[
0=1.
\]
Por tanto, el sistema no tiene soluciones, es \alert{incompatible}.

\item Supongamos que la última columna no contiene ningún pivote.
Sean
\[
j_1<j_2<\cdots<j_r
\]
las columnas que contienen los pivotes. Las incógnitas
\[
x_{j_1},\ldots,x_{j_r}
\]
se denominan \alert{incógnitas pivote} y las restantes
\alert{incógnitas libres}.
\end{itemize}

\end{frame}

\begin{frame}{Obtención de las soluciones}

Supongamos que la última columna no contiene ningún pivote. Sean
\[
j_1<\cdots<j_r
\]
las columnas pivote y
\[
l_1<\cdots<l_{n-r}
\]
las restantes columnas.

Como la matriz está reducida, las ecuaciones correspondientes a sus
filas no nulas tienen la forma
\[
x_{j_i}
+a_{i l_1}x_{l_1}+\cdots+a_{i l_{n-r}}x_{l_{n-r}}
=b_i,
\qquad 1\leq i\leq r.
\]

Las incógnitas libres pueden tomar valores arbitrarios
\[
x_{l_1}=\lambda_1,\ldots,x_{l_{n-r}}=\lambda_{n-r},
\qquad
\lambda_1,\ldots,\lambda_{n-r}\in k.
\]

\end{frame}

\begin{frame}[t]{Obtención de las soluciones}

Para cada elección de estos valores, las incógnitas pivote quedan
determinadas de manera única:
\[
x_{j_i}
=
b_i-\sum_{s=1}^{n-r}a_{i l_s}\lambda_s,
\qquad 1\leq i\leq r.
\]

De este modo se obtienen todas las soluciones del sistema, que es \alert{compatible}. El sistema es \alert{determinado} cuando no hay incógnitas libres, es decir, cuando todas las columnas contienen un pivote, excepto la última.

\medskip

Por tanto, para resolver un sistema de ecuaciones lineales, podemos escalonar y reducir su matriz ampliada mediante el procedimiento de eliminación gaussiana y resolver el sistema equivalente cuya matriz ampliada es esta forma escalonada y reducida.

\end{frame}


\begin{frame}{Teorema de Rouché--Frobenius}

\begin{theorem}[Rouché--Frobenius]
Consideremos un sistema de \(m\) ecuaciones lineales con \(n\)
incógnitas
\[
Ax=b.
\]
Sea $R$ una forma escalonada y reducida de su matriz ampliada\((A\mid b).
\)

\medskip

El carácter del sistema viene determinado por las posiciones de los
pivotes de \(R\):
\begin{enumerate}
\item Si la última columna de \(R\) contiene un pivote, el sistema es
\alert{incompatible}.

\item Si la última columna no contiene ningún pivote pero todas las demás sí, el sistema es
\alert{compatible determinado}.

\item Si la última columna no contiene ningún pivote y alguna de las
otras tampoco, el sistema es
\alert{compatible indeterminado}.
\end{enumerate}
\end{theorem}

\end{frame}

\begin{frame}[allowframebreaks, t]{Ejemplo: sistema compatible indeterminado}

Consideremos el sistema sobre el cuerpo $k=\Q$ cuya matriz ampliada es
\[
\left(
\begin{array}{ccccc|c}
1&2&0&-1&0&3\\
0&0&1&4&0&2\\
0&0&0&0&1&-1
\end{array}
\right).
\]

La matriz de coeficientes está escalonada y reducida. No hay pivotes en la última columna, así que el sistema es compatible. Es más, dos de las columnas anteriores tampoco contienen pivotes, con lo cual es indeterminado.

\medskip

Las incógnitas
pivote son
\[
x_1,\quad x_3,\quad x_5,
\]
y las incógnitas libres son \(x_2,x_4\).

Asignamos a estas últimas valores arbitrarios
\[
x_2=\lambda,\qquad x_4=\mu,
\qquad \lambda,\mu\in \Q.
\]

Las ecuaciones permiten obtener las incógnitas pivote:
\[
x_1=3-2\lambda+\mu,\qquad
x_3=2-4\mu,\qquad
x_5=-1.
\]

Por tanto, todas las soluciones son
\[
v=
\begin{pmatrix}
3-2\lambda+\mu\\
\lambda\\
2-4\mu\\
\mu\\
-1
\end{pmatrix}
,
\qquad \lambda,\mu\in \Q.
\]

\end{frame}


\begin{frame}{Resolución de un sistema escalonado}

Supongamos que la matriz ampliada de un sistema está
\alert{escalonada por filas}.

\begin{itemize}
\item Si aparece una fila de la forma
\[
\left(
\begin{array}{cccc|c}
0&\cdots&0&0&c
\end{array}
\right),
\qquad c\neq0,
\]
el sistema es incompatible.

\item En caso contrario, las incógnitas correspondientes a columnas
sin pivote son \alert{incógnitas libres} y pueden tomar valores
arbitrarios.

\item Asignamos a las incógnitas libres valores
\[
\lambda_1,\ldots,\lambda_s\in k.
\]

\item Comenzando por la última fila no nula y procediendo hacia arriba,
cada ecuación determina la incógnita correspondiente a su pivote
en función de las incógnitas ya conocidas.
\end{itemize}

\end{frame}

\begin{frame}{Sistemas homogéneos}

\begin{definition}
Un sistema de ecuaciones lineales
\[
Ax=b
\]
se denomina \alert{homogéneo} si \(b=0\), es decir, si es de la forma
\[
Ax=0.
\]
\end{definition}

Todo sistema homogéneo tiene al menos una solución:
\[
v=0.
\]

\end{frame}

\begin{frame}{Sistemas homogéneos}

Además, sus soluciones son cerradas para la suma y el producto por
escalares. En efecto, si \(v,w\in k^n\) son soluciones y
\(\lambda\in k\), entonces
\[
A(v+w)=Av+Aw=0+0=0,
\]
por lo que \(v+w\) es también una solución, y
\[
A(\lambda v)=\lambda Av=\lambda0=0,
\]
por lo que \(\lambda v\) es también una solución.

\end{frame}

\begin{frame}{Soluciones de un sistema homogéneo}

Consideremos un sistema homogéneo \(Ax=0\) cuyas incógnitas libres,
tras reducirlo, son
\[
x_{l_1},\ldots,x_{l_s}.
\]
Toda solución se obtiene para unos \alert{únicos} valores
\[
x_{l_1}=\lambda_1,\ldots,x_{l_s}=\lambda_s,
\qquad
\lambda_1,\ldots,\lambda_s\in k,
\]
pues estos son precisamente los valores de las incógnitas libres.

Para cada \(i=1,\ldots,s\), sea \(v_i\) la solución obtenida tomando
\[
\lambda_i=1,
\qquad
\lambda_j=0\quad\text{si }j\neq i.
\]

\begin{proposition}
Para cada solución \(v\) existen unos únicos
\(\lambda_1,\ldots,\lambda_s\in k\) tales que
\[
v=\lambda_1v_1+\cdots+\lambda_sv_s.
\]
\end{proposition}

\end{frame}

\begin{frame}{Soluciones de un sistema homogéneo}

\begin{proof}
Sean \(\lambda_1,\ldots,\lambda_s\) los valores de las incógnitas
libres de \(v\). La combinación lineal
\[
w=\lambda_1v_1+\cdots+\lambda_sv_s
\]
es una solución y sus incógnitas libres toman esos mismos valores.
Como los valores de las incógnitas libres determinan de manera única
las incógnitas pivote, \(v=w\).

Además, si
\[
v=\mu_1v_1+\cdots+\mu_sv_s,
\]
el valor de la incógnita libre \(x_{l_i}\) del segundo miembro es
\(\mu_i\). Por tanto, necesariamente
\[
\mu_i=\lambda_i,\qquad i=1,\ldots,s.
\]
\end{proof}

\end{frame}

\begin{frame}[allowframebreaks, t]{Ejemplo: soluciones de un sistema homogéneo}

Consideremos el sistema homogéneo sobre el cuerpo $k=\Q$ cuya matriz de coeficientes
escalonada y reducida es
\[
\begin{pmatrix}
1&2&0&-1&0\\
0&0&1&4&0\\
0&0&0&0&1
\end{pmatrix}.
\]
Las incógnitas libres son \(x_2,x_4\). Tomando
\[
x_2=\lambda,\qquad x_4=\mu,
\]
obtenemos
\[
x_1=-2\lambda+\mu,\qquad
x_3=-4\mu,\qquad
x_5=0.
\]
Por tanto, la solución correspondiente a \(\lambda,\mu\in \Q\) es
\[
v=
\begin{pmatrix}
-2\lambda+\mu\\
\lambda\\
-4\mu\\
\mu\\
0
\end{pmatrix}.
\]

Tomando respectivamente \((\lambda,\mu)=(1,0)\) y \((0,1)\),
obtenemos
\[
v_1=
\begin{pmatrix}-2\\1\\0\\0\\0\end{pmatrix},
\qquad
v_2=
\begin{pmatrix}1\\0\\-4\\1\\0\end{pmatrix}.
\]

Así,
\[
v=\lambda v_1+\mu v_2,
\qquad \lambda,\mu\in \Q,
\]
y cada solución se expresa de esta forma para unos
\alert{únicos} valores de \(\lambda\) y \(\mu\).

\end{frame}

\begin{frame}{Sistema homogéneo asociado}

\begin{definition}
Dado un sistema de ecuaciones lineales
\[
Ax=b,
\]
su \alert{sistema homogéneo asociado} u \alert{homogeneizado} es
\[
Ax=0.
\]
\end{definition}

\end{frame}

\begin{frame}{Soluciones de un sistema en función de las del homogeneizado}

\begin{proposition}
Supongamos que \(v_0\) es una solución particular de \(Ax=b\).
Entonces \(v\in k^n\) es solución de \(Ax=b\) si y solo si
\[
v=v_0+w,
\]
para alguna solución \(w\) del sistema homogéneo asociado \(Ax=0\).
\end{proposition}

\begin{proof}
Si \(v\) es una solución de \(Ax=b\), entonces
\[
A(v-v_0)=Av-Av_0=b-b=0.
\]
Por tanto, \(w=v-v_0\) es solución del sistema homogéneo asociado y
\(v=v_0+w\).

Recíprocamente, si \(Aw=0\), entonces
\[
A(v_0+w)=Av_0+Aw=b+0=b,
\]
por lo que \(v_0+w\) es solución de \(Ax=b\).
\end{proof}

\end{frame}


\begin{frame}[allowframebreaks, t]{Ejemplo: sistema compatible indeterminado}

Anteriormente, habíamos obtenido las siguientes soluciones de un sistema no homogéneo
\[
v=
\begin{pmatrix}
3-2\lambda+\mu\\
\lambda\\
2-4\mu\\
\mu\\
-1
\end{pmatrix}
,
\qquad \lambda,\mu\in k.
\]

\end{frame}

\begin{frame}[allowframebreaks, t]{Ejemplo: sistema compatible indeterminado}

Anteriormente, habíamos obtenido las siguientes soluciones de un sistema no homogéneo sobre el cuerpo \(k=\Q\):
\[
v=
\begin{pmatrix}
3-2\lambda+\mu\\
\lambda\\
2-4\mu\\
\mu\\
-1
\end{pmatrix}
,
\qquad \lambda,\mu\in \Q.
\]

Tomando \(\lambda=\mu=0\), obtenemos la solución particular
\[
v_0=
\begin{pmatrix}
3\\
0\\
2\\
0\\
-1
\end{pmatrix}.
\]

Las soluciones del sistema homogéneo asociado se obtienen eliminando
los términos independientes:
\[
w=
\begin{pmatrix}
-2\lambda+\mu\\
\lambda\\
-4\mu\\
\mu\\
0
\end{pmatrix}.
\]

Tomando respectivamente
\(
(\lambda,\mu)=(1,0)
\)
y
\(
(\lambda,\mu)=(0,1),
\)
obtenemos las siguientes soluciones del sistema homogéneo asociado:
\[
w_1=
\begin{pmatrix}
-2\\
1\\
0\\
0\\
0
\end{pmatrix},
\qquad
w_2=
\begin{pmatrix}
1\\
0\\
-4\\
1\\
0
\end{pmatrix}
\]

Por tanto, la solución general del sistema original puede escribirse
como
\[
v=v_0+\lambda w_1+\mu w_2,
\qquad \lambda,\mu\in \Q,
\]
es decir,
\[
v=
\begin{pmatrix}
3\\
0\\
2\\
0\\
-1
\end{pmatrix}
+
\lambda
\begin{pmatrix}
-2\\
1\\
0\\
0\\
0
\end{pmatrix}
+
\mu
\begin{pmatrix}
1\\
0\\
-4\\
1\\
0
\end{pmatrix},
\qquad \lambda,\mu\in \Q.
\]

Además, cada solución se obtiene de esta forma para unos
\alert{únicos} valores de \(\lambda\) y \(\mu\).

\end{frame}

\begin{frame}[allowframebreaks, t]
{Sistemas homogéneos escalonados y reducidos}

\begin{proposition}
Sean \(Ax=0\) y \(Bx=0\) dos sistemas homogéneos cuyas matrices de
coeficientes \(A\) y \(B\) están escalonadas y reducidas.
Si ambos sistemas tienen las mismas soluciones, entonces
\[
A=B.
\]
\end{proposition}

\begin{proof}
Observemos primero que, para una matriz escalonada y reducida, la
columna \(j\) contiene un pivote si y solo si toda solución que
satisface
\[
x_{j+1}=\cdots=x_n=0
\]
satisface también \(x_j=0\).

En efecto, si la columna \(j\) contiene un pivote, la ecuación
correspondiente tiene la forma
\[
x_j+a_{j+1}x_{j+1}+\cdots+a_nx_n=0,
\]
por lo que \(x_{j+1}=\cdots=x_n=0\) implica \(x_j=0\).

Recíprocamente, si la columna \(j\) no contiene un pivote, \(x_j\)
es una incógnita libre. Podemos tomar
\[
x_j=1
\]
y hacer iguales a \(0\) todas las demás incógnitas libres.
Las incógnitas pivote quedan entonces determinadas y obtenemos una
solución que satisface
\[
x_j=1,
\qquad
x_{j+1}=\cdots=x_n=0.
\]

Por tanto, las columnas que contienen pivotes están completamente
determinadas por el conjunto de soluciones.

Como \(Ax=0\) y \(Bx=0\) tienen las mismas soluciones, \(A\) y \(B\)
tienen las mismas columnas pivote y, por tanto, las mismas incógnitas
pivote y las mismas incógnitas libres.

Sean
\[
j_1<\cdots<j_r
\]
las columnas pivote y
\[
l_1<\cdots<l_{n-r}
\]
las columnas libres. La ecuación correspondiente al pivote \(j_i\)
en el primer sistema es
\[
x_{j_i}
+\sum_{s=1}^{n-r}a_{i l_s}x_{l_s}=0.
\]

Fijemos \(l_s\). Tomando
\[
x_{l_s}=1
\]
y haciendo iguales a \(0\) las demás incógnitas libres, obtenemos
una única solución, para la cual
\[
x_{j_i}=-a_{i l_s}.
\]

Haciendo la misma elección de las incógnitas libres en el segundo
sistema obtenemos
\[
x_{j_i}=-b_{i l_s}.
\]
Como los dos sistemas tienen las mismas soluciones y las incógnitas
pivote quedan determinadas de manera única por las libres,
\[
a_{i l_s}=b_{i l_s}.
\]

Así, \(A\) y \(B\) tienen los mismos pivotes y los mismos coeficientes
en las columnas libres. Al estar ambas escalonadas y reducidas,
todas sus entradas coinciden. Por tanto,
\[
A=B.
\]
\end{proof}

\end{frame}

\begin{frame}[allowframebreaks, t]{Forma escalonada y reducida}

\begin{theorem}
Toda matriz \(A\in M_{m\times n}(k)\) es equivalente por filas a una
única matriz escalonada y reducida, su \alert{forma escalonada y reducida}.
\end{theorem}

\begin{proof}
Ya sabemos que, mediante eliminación de Gauss y operaciones
elementales por filas, toda matriz es equivalente por filas a una
matriz escalonada y reducida. Por tanto, basta probar la unicidad.

\medskip

Supongamos que \(R,S\in M_{m\times n}(k)\) son matrices escalonadas
y reducidas tales que
\[
A\sim_f R
\qquad\text{y}\qquad
A\sim_f S.
\]

Consideremos los sistemas homogéneos
\[
Ax=0,\qquad Rx=0,\qquad Sx=0.
\]

Como \(A\sim_f R\), los sistemas
\[
Ax=0
\qquad\text{y}\qquad
Rx=0
\]
tienen las mismas soluciones. Del mismo modo, como \(A\sim_f S\),
los sistemas
\[
Ax=0
\qquad\text{y}\qquad
Sx=0
\]
tienen las mismas soluciones.

\medskip

Por tanto, los sistemas homogéneos
\[
Rx=0
\qquad\text{y}\qquad
Sx=0
\]
tienen las mismas soluciones.

\medskip

Como \(R\) y \(S\) están escalonadas y reducidas, el resultado
anterior implica que
\[
R=S.
\]

Por tanto, \(A\) es equivalente por filas a una única matriz
escalonada y reducida.
\end{proof}

\end{frame}

\begin{frame}{Rango de una matriz}

\begin{definition}
Sea \(A\in M_{m\times n}(k)\). El \alert{rango} de \(A\), que
denotaremos por
\[
\operatorname{rg}(A),
\]
es el número de pivotes de la forma escalonada y reducida de \(A\).
\end{definition}

\medskip

Equivalentemente, \(\operatorname{rg}(A)\) es el número de filas
no nulas de la forma escalonada y reducida de \(A\), o simplemente de una forma escalonada de $A$.

\medskip

En particular,
\[
0\leq \operatorname{rg}(A)\leq \min\{m,n\}.
\]

\end{frame}

\begin{frame}{Teorema de Rouché--Frobenius}

\begin{theorem}[Rouché--Frobenius]
Consideremos un sistema de \(m\) ecuaciones lineales con \(n\)
incógnitas
\(
Ax=b.
\)
Entonces:
\begin{enumerate}
\item El sistema es \alert{compatible} si y solo si
\[
\operatorname{rg}(A)=\operatorname{rg}(A\mid b).
\]

\item Si el sistema es compatible, es
\alert{compatible determinado} si y solo si
\[
\operatorname{rg}(A)=\operatorname{rg}(A\mid b)=n.
\]

\item Si el sistema es compatible, es
\alert{compatible indeterminado} si y solo si
\[
\operatorname{rg}(A)=\operatorname{rg}(A\mid b)<n.
\]
\end{enumerate}
\end{theorem}

\end{frame}


\begin{frame}[allowframebreaks, t]{Demostración del teorema de Rouché--Frobenius}

\begin{proof}
Mediante eliminación gaussiana llevamos la matriz ampliada a su
forma escalonada y reducida:
\[
(A\mid b)\sim_f(A'\mid b').
\]

Las mismas operaciones elementales transforman \(A\) en \(A'\).
Además, como \((A'\mid b')\) está escalonada y reducida, también
lo está \(A'\), así que
\(A'\) es la forma escalonada y reducida de \(A\).

Por tanto,
\[
\operatorname{rg}(A)
=
\text{número de pivotes de }A',
\]
mientras que
\[
\operatorname{rg}(A\mid b)
=
\text{número de pivotes de }(A'\mid b').
\]

Como \((A'\mid b')\) solo tiene una columna más que \(A'\), estos
rangos son distintos precisamente cuando la última columna contiene
un pivote, lo que equivale a la incompatibilidad del sistema.

Por tanto, el sistema es compatible si y solo si
\[
\operatorname{rg}(A)=\operatorname{rg}(A\mid b).
\]

Supongamos ahora que el sistema es compatible. Entonces la última
columna no contiene ningún pivote y
\[
\operatorname{rg}(A)
\]
es el número de columnas pivote entre las \(n\) columnas
correspondientes a las incógnitas.

Por tanto, todas las incógnitas son pivote si y solo si
\[
\operatorname{rg}(A)=n.
\]
En este caso no hay incógnitas libres y el sistema es compatible
determinado.

Por el contrario, existe alguna incógnita libre si y solo si
\[
\operatorname{rg}(A)<n.
\]
En este caso el sistema es compatible indeterminado.
\end{proof}

\end{frame}

\begin{frame}{Rango y sistemas homogéneos}

\begin{corollary}
Sea \(A\in M_{m\times n}(k)\). Entonces
\[
\operatorname{rg}(A)=n
\]
si y solo si el sistema homogéneo
\[
Ax=0
\]
tiene como única solución \(v=0\).
\end{corollary}

\end{frame}

\begin{frame}[allowframebreaks, t]{Rango de una matriz invertible}

\begin{proposition}
Si \(A\in M_{m\times n}(k)\) es invertible, entonces \(m=n\) y
\[
\operatorname{rg}(A)=n.
\]
\end{proposition}

\vspace{-5mm}

\begin{proof}
Si \(v\) es una solución de
\(
Ax=0,
\)
entonces
\[
v=I_nv=A^{-1}Av=A^{-1}0=0.
\]
Por tanto, el sistema \(Ax=0\) tiene una única solución y, por el
resultado anterior,
\[
\operatorname{rg}(A)=n.
\]
Como \(\operatorname{rg}(A)\leq m\), se tiene \(n\leq m\).

La inversa de una matriz invertible también es invertible. Como $A^{-1}\in M_{n\times m}(k)$, si aplicamos lo anterior a $A^{-1}$, deducimos que $\operatorname{rg}(A^{-1})=m$ y \(m\leq n\).
Por tanto, por doble desigualdad,
\[
m=n,
\]
y \(A\) es cuadrada de rango \(n\).
\end{proof}

\end{frame}

\begin{frame}[allowframebreaks, t]{Matrices cuadradas de rango máximo}

\begin{lemma}
	Una matriz cuadrada \(A\in M_n(k)\) tiene
	\(
	\operatorname{rg}(A)=n
	\)
	si y solo si su forma escalonada y reducida es la matriz identidad $I_n$.
\end{lemma}

\begin{proof}
Sea $R$ la forma escalonada y reducida de $A$.
Como
\(
\operatorname{rg}(A)=n,
\)
la matriz \(R\) tiene \(n\) pivotes. Por tanto, todas sus columnas
contienen un pivote y, al ser \(R\) escalonada y reducida,
\(
R=I_n.
\)
\end{proof}


\begin{corollary}
Sea \(A\in M_n(k)\). Si
\(
\operatorname{rg}(A)=n,
\)
entonces \(A\) es producto de matrices elementales, en particular es  invertible.
\end{corollary}

\begin{proof}
La forma escalonada y reducida de $A$ es $I_n$, así que existe una matriz invertible $$P=E_1\cdots E_r,$$ producto de elementales $E_i$, $1\leq i\leq r$, tal que $$PA=I_n.$$ Por tanto 
\begin{align*}
A&=P^{-1}\\
&=(E_1\cdots E_r)^{-1}\\
&=E_r^{-1}\cdots E_1^{-1}
\end{align*}
es producto de matrices elementales, puesto que la inversa de una matriz elemental es elemental.
\end{proof}

\end{frame}

\begin{frame}{Caracterización de las matrices invertibles}

\begin{theorem}
Sea \(A\in M_{m\times n}(k)\). Entonces,
\begin{align*}
A\text{ es invertible}
&\Longleftrightarrow
m=n=\operatorname{rg}(A)\\
&\Longleftrightarrow
A\text{ es producto de matrices elementales}.
\end{align*}
\end{theorem}

Estas caracterizaciones de las matrices invertibles son consecuencia de los resultados anteriores.

\end{frame}

\begin{frame}{Caracterización de la equivalencia por filas}

\begin{corollary}
	Dos matrices $A,B \in M_{m\times n}(k)$ son equivalentes por filas si y solo si existe una matriz invertible $P\in M_{n\times n}(k)$ tal que $PA=B$.
\end{corollary}

\begin{proof}
	Si $A$ y $B$ son equivalentes por filas, podemos pasar de $A$ a $B$ realizando operaciones elementales por filas. Entonces, si $P$ es el producto de matrices elementales correspondiente a dichas operaciones, $PA=B$. La matriz $P$ es invertible por ser producto de elementales, que son invertibles.
	
	\medskip
	
	Si $PA=B$ con $P$ invertible, entonces $P$ es producto de elementales, por tanto podemos pasar de $A$ a $B$ realizando las correspondientes operaciones elementales, así que $A$ y $B$ son equivalentes por filas.
\end{proof}

\end{frame}

\begin{frame}[allowframebreaks, t]{Cálculo de la inversa por eliminación gaussiana}

\begin{proposition}
    Dada una matriz invertible \(A\in M_n(k)\), la forma escalonada y reducida de la matriz por bloques
    \[
    (A\mid I_n)
    \]
    es
    \[
    (I_n\mid A^{-1}).
    \]
\end{proposition}

\begin{proof}
Como \(A\) es invertible, su rango es \(n\). Por tanto, la forma escalonada y reducida de \(A\) es
\(
I_n.
\)

\medskip

Aplicamos a la matriz por bloques
\[
(A\mid I_n)
\]
las operaciones elementales por filas que transforman \(A\) en
\(I_n\).

\medskip

Si \(P\) es el producto de las matrices elementales correspondientes
a estas operaciones, obtenemos
\[
P(A\mid I_n)
=
(PA\mid P)
=
(I_n\mid P).
\]
Esta última matriz está escalonada y reducida.

\medskip

Como
\[
PA=I_n,
\]
se tiene que
\[
A^{-1}=P.
\]
\end{proof}
\end{frame}

\begin{frame}[allowframebreaks, t]{Ejemplo: cálculo de una inversa}

Consideremos
\[
A=
\begin{pmatrix}
1&1&0\\
0&1&1\\
1&1&1
\end{pmatrix}.
\]

Realizamos el siguiente cálculo de forma escalonada y reducida por eliminación gaussiana:
\[
(A\mid I_3)=
\left(
\begin{array}{ccc|ccc}
1&1&0&1&0&0\\
0&1&1&0&1&0\\
1&1&1&0&0&1
\end{array}
\right)
\xrightarrow{F_3-F_1}
\left(
\begin{array}{ccc|ccc}
1&1&0&1&0&0\\
0&1&1&0&1&0\\
0&0&1&-1&0&1
\end{array}
\right)
\]

\[
\xrightarrow{F_2-F_3}
\left(
\begin{array}{ccc|ccc}
1&1&0&1&0&0\\
0&1&0&1&1&-1\\
0&0&1&-1&0&1
\end{array}
\right)
\xrightarrow{F_1-F_2}
\left(
\begin{array}{ccc|ccc}
1&0&0&0&-1&1\\
0&1&0&1&1&-1\\
0&0&1&-1&0&1
\end{array}
\right).
\]

Por tanto,
\[
A^{-1}
=
\begin{pmatrix}
0&-1&1\\
1&1&-1\\
-1&0&1
\end{pmatrix}.
\]

\end{frame}
\begin{frame}{Definición inductiva del determinante}

Sea \(A=(a_{ij})\in M_n(k)\). Denotaremos por \(A_{ij}\) la matriz
de orden \(n-1\) que se obtiene de \(A\) suprimiendo la fila \(i\)
y la columna \(j\).

\medskip

\begin{center}
\begin{tikzpicture}[baseline=(M.center)]
\matrix (M) [
    matrix of math nodes,
    ampersand replacement=\&,
    left delimiter=(,
    right delimiter=),
    row sep=2mm,
    column sep=3mm
] {
a_{11} \& \cdots \& a_{1j} \& \cdots \& a_{1n}\\
\vdots \&        \& \vdots \&        \& \vdots\\
a_{i1} \& \cdots \& a_{ij} \& \cdots \& a_{in}\\
\vdots \&        \& \vdots \&        \& \vdots\\
a_{n1} \& \cdots \& a_{nj} \& \cdots \& a_{nn}\\
};

\draw[ultra thick, opacity =.5]
    ([xshift=-2mm]M-3-1.west)
    --
    ([xshift=2mm]M-3-5.east);

\draw[ultra thick, opacity =.5]
    ([yshift=2mm]M-1-3.north)
    --
    ([yshift=-2mm]M-5-3.south);

\node[right=7mm of M] {\(= A_{ij}\)};
\end{tikzpicture}
\end{center}
\end{frame}

\begin{frame}{Definición inductiva del determinante}

\begin{definition}
Definimos inductivamente el \alert{determinante} de una matriz
cuadrada.

Si \(A=(a)\in M_1(k)\), definimos
\[
\det(A)=a.
\]

Si \(n>1\), suponiendo definido el determinante para matrices de
orden \(n-1\), para \(A=(a_{ij})\in M_n(k)\) definimos
\[
\det(A)
=
\sum_{j=1}^n(-1)^{1+j}a_{1j}\det(A_{1j}).
\]
Esta fórmula se conoce como el \alert{desarrollo del determinante por la primera fila}.

El determinante de una matriz $A$ también se denotará por 
\[|A|=\det(A).\]
\end{definition}

\end{frame}

\begin{frame}{Determinantes de orden \(2\)}

Sea
\[
A=
\begin{pmatrix}
a&b\\
c&d
\end{pmatrix}.
\]

Por la definición inductiva,
\[
\begin{aligned}
\det(A)
&=
a\det(d)-b\det(c)\\
&=ad-bc.
\end{aligned}
\]

Por tanto,
\[
\det
\begin{pmatrix}
a&b\\
c&d
\end{pmatrix}
=
ad-bc.
\]

\end{frame}

\begin{frame}[allowframebreaks]{Determinantes de orden \(3\)}

Sea
\[
A=
\begin{pmatrix}
a_{11}&a_{12}&a_{13}\\
a_{21}&a_{22}&a_{23}\\
a_{31}&a_{32}&a_{33}
\end{pmatrix}.
\]

Desarrollando por la primera fila,
\[
\begin{aligned}
\det(A)
={}&
a_{11}
\begin{vmatrix}
a_{22}&a_{23}\\
a_{32}&a_{33}
\end{vmatrix}
-
a_{12}
\begin{vmatrix}
a_{21}&a_{23}\\
a_{31}&a_{33}
\end{vmatrix}\\
&+
a_{13}
\begin{vmatrix}
a_{21}&a_{22}\\
a_{31}&a_{32}
\end{vmatrix}.
\end{aligned}
\]

Usando la fórmula para determinantes de orden \(2\),
\[
\begin{aligned}
\det(A)
={}&
a_{11}a_{22}a_{33}
-a_{11}a_{23}a_{32}\\
&-a_{12}a_{21}a_{33}
+a_{12}a_{23}a_{31}\\
&+a_{13}a_{21}a_{32}
-a_{13}a_{22}a_{31}.
\end{aligned}
\]
\end{frame}

\begin{frame}[allowframebreaks]{Determinantes de orden \(3\)}

\begin{columns}
\begin{column}{0.5\textwidth}
Reordenando los términos obtenemos la \alert{regla de Sarrus}:
\[
\begin{aligned}
\det(A)
={}&
a_{11}a_{22}a_{33}
+a_{12}a_{23}a_{31}
+a_{13}a_{21}a_{32}\\
&-
a_{13}a_{22}a_{31}
-a_{11}a_{23}a_{32}
-a_{12}a_{21}a_{33}.
\end{aligned}
\]
Esta regla solo es válida para matrices de orden \(3\).
\end{column}

\begin{column}{0.5\textwidth}
% Source - https://tex.stackexchange.com/a/603570
% Posted by SebGlav
% Retrieved 2026-09-24, License - CC BY-SA 4.0
\begin{center}
\begin{tikzpicture}[baseline=(A.center)]
  \tikzset{BarreStyle/.style =   {opacity=.4,line width=4 mm,line cap=round,color=#1}}
    \tikzset{SignePlus/.style =   {above left,inner sep=1.5pt,opacity=.4,circle,fill=#1}}
    \tikzset{SigneMoins/.style =   {below left,inner sep=1.5pt,opacity=.4,circle,fill=#1}}
% the matrices
\matrix (A) [matrix of math nodes, nodes = {node style ge},column sep=0 mm, ampersand replacement=\&,] 
{ a_{11} \& a_{12} \& a_{13}  \\
  a_{21} \& a_{22} \& a_{23}  \\
  a_{31} \& a_{32} \& a_{33}  \\
  a_{11} \& a_{12} \& a_{13} \\
  a_{21} \& a_{22} \& a_{13}\\
};

 \draw [BarreStyle=blue] (A-1-1.north west) node[SignePlus=blue] {$+$} to (A-3-3.south east) ;
 \draw [BarreStyle=blue] (A-2-1.north west) node[SignePlus=blue] {$+$} to (A-4-3.south east) ;
 \draw [BarreStyle=blue] (A-3-1.north west) node[SignePlus=blue] {$+$} to (A-5-3.south east) ;
 \draw [BarreStyle=red]  (A-3-1.south west) node[SigneMoins=red] {$-$} to (A-1-3.north east);
 \draw [BarreStyle=red]  (A-4-1.south west) node[SigneMoins=red] {$-$} to (A-2-3.north east);
 \draw [BarreStyle=red]  (A-5-1.south west) node[SigneMoins=red] {$-$} to (A-3-3.north east);
\end{tikzpicture}
\end{center}
\end{column}
\end{columns}

\end{frame}

\begin{frame}[allowframebreaks]{Linealidad respecto de las filas}

\begin{proposition}
El determinante es lineal respecto de cada fila. Más concretamente,
si \(u,v\in M_{1\times n}(k)\) y \(\lambda\in k\), entonces
\[
\det
\left(
\begin{array}{c}
\vdots\\
\hline
u+v\\
\hline
\vdots
\end{array}
\right)
=
\det
\left(
\begin{array}{c}
\vdots\\
\hline
u\\
\hline
\vdots
\end{array}
\right)
+
\det
\left(
\begin{array}{c}
\vdots\\
\hline
v\\
\hline
\vdots
\end{array}
\right),
\]
y
\[
\det
\left(
\begin{array}{c}
\vdots\\
\hline
\lambda u\\
\hline
\vdots
\end{array}
\right)
=
\lambda
\det
\left(
\begin{array}{c}
\vdots\\
\hline
u\\
\hline
\vdots
\end{array}
\right),
\]
donde las restantes filas son las mismas en todas las matrices.
\end{proposition}

\end{frame}

\begin{frame}[allowframebreaks, t]{Linealidad respecto de la primera fila}

Sean
\[
u=(u_1,\ldots,u_n),\qquad
v=(v_1,\ldots,v_n).
\]
Mantengamos fijas las demás filas de la matriz.

\medskip

Por la definición inductiva del determinante,
\[
\begin{aligned}
\det
\left(
\begin{array}{c}
u+v\\
\hline
\vdots
\end{array}
\right)
&=
\sum_{j=1}^n
(-1)^{1+j}(u_j+v_j)\det(A_{1j})\\
&=
\sum_{j=1}^n
(-1)^{1+j}u_j\det(A_{1j})
+
\sum_{j=1}^n
(-1)^{1+j}v_j\det(A_{1j})\\
&=
\det
\left(
\begin{array}{c}
u\\
\hline
\vdots
\end{array}
\right)
+
\det
\left(
\begin{array}{c}
v\\
\hline
\vdots
\end{array}
\right).
\end{aligned}
\]

Del mismo modo, para \(\lambda\in k\),
\[
\begin{aligned}
\det
\left(
\begin{array}{c}
\lambda u\\
\hline
\vdots
\end{array}
\right)
&=
\sum_{j=1}^n
(-1)^{1+j}(\lambda u_j)\det(A_{1j})\\
&=
\lambda\sum_{j=1}^n
(-1)^{1+j}u_j\det(A_{1j})\\
&=
\lambda\det
\left(
\begin{array}{c}
u\\
\hline
\vdots
\end{array}
\right)
\end{aligned}
\]

Por tanto, el determinante es lineal respecto de la primera fila.

\end{frame}

\begin{frame}[allowframebreaks, t]{Linealidad respecto otra fila}

El resultado es obvio para matrices de orden \(1\). Supongamos que \(n>1\) y que el resultado es cierto para matrices de orden \(n-1\).

\medskip

Supongamos también que la suma de vectores o el producto por un escalar se realiza en otra fila, la fila \(i>1\). 

\medskip

Dado un vector $w=(w_1,\ldots,w_n)$, denotamos 
\[w^j=(w_1,\ldots,w_{j-1},w_{j+1},\ldots,w_n)\]
el vector obtenido suprimiendo su componente \(j\). 

\framebreak

La definición inductiva del determinante y la hipótesis de inducción nos dan:
\begin{multline*}
\det
\left(
\begin{array}{c}
\vdots\\
\hline
u+v\\
\hline
\vdots
\end{array}
\right)
=
\sum_{j=1}^n
(-1)^{1+j}a_{1j}\det
\left(
\begin{array}{c}
\vdots\\
\hline\\[-12pt]
u^j+v^j\\
\hline
\vdots
\end{array}
\right)\\
=
\sum_{j=1}^n
(-1)^{1+j}a_{1j}\left(\det
\left(
\begin{array}{c}
\vdots\\
\hline\\[-12pt]
u^j\\
\hline
\vdots
\end{array}
\right)
+
\det
\left(
\begin{array}{c}
\vdots\\
\hline\\[-12pt]
v^j\\
\hline
\vdots
\end{array}
\right)\right)
\\
=
\sum_{j=1}^n
(-1)^{1+j}a_{1j}\det
\left(
\begin{array}{c}
\vdots\\
\hline\\[-12pt]
u^j\\
\hline
\vdots
\end{array}
\right)
+
\sum_{j=1}^n
(-1)^{1+j}a_{1j}\det
\left(
\begin{array}{c}
\vdots\\
\hline\\[-12pt]
v^j\\
\hline
\vdots
\end{array}
\right)
=
\det
\left(
\begin{array}{c}
\vdots\\
\hline
u\\
\hline
\vdots
\end{array}
\right)
+
\det
\left(
\begin{array}{c}
\vdots\\
\hline
v\\
\hline
\vdots
\end{array}
\right).
\end{multline*}

\framebreak

Del mismo modo, para \(\lambda\in k\),
\[
\begin{aligned}
\det
\left(
\begin{array}{c}
\vdots\\
\hline
\lambda u\\
\hline
\vdots
\end{array}
\right)
&=
\sum_{j=1}^n
(-1)^{1+j}a_{1j}\det
\left(
\begin{array}{c}
\vdots\\
\hline\\[-12pt]
\lambda u^j\\
\hline
\vdots
\end{array}
\right)=
\sum_{j=1}^n
(-1)^{1+j}a_{1j}\lambda\det
\left(
\begin{array}{c}
\vdots\\
\hline\\[-12pt]
u^j\\
\hline
\vdots
\end{array}
\right)\\
&=
\lambda\sum_{j=1}^n
(-1)^{1+j}a_{1j}\det
\left(
\begin{array}{c}
\vdots\\
\hline\\[-12pt]
u^j\\
\hline
\vdots
\end{array}
\right)
=
\lambda\det
\left(
\begin{array}{c}
\vdots\\
\hline
u\\
\hline
\vdots
\end{array}
\right).
\end{aligned}
\]

Por tanto, el determinante es lineal respecto de cualquier fila.

\end{frame}

\begin{frame}[allowframebreaks, t]{Dos filas consecutivas iguales}

\begin{proposition}
Si una matriz cuadrada tiene dos filas consecutivas iguales, entonces
su determinante es \(0\).
\end{proposition}

\begin{proof}
Procedemos por inducción sobre el orden \(n\).

Para \(n=2\),
\[
\det
\begin{pmatrix}
a&b\\
a&b
\end{pmatrix}
=ab-ba=0.
\]

Supongamos ahora \(n>2\) y que el resultado es cierto para matrices
de orden \(n-1\).

Supongamos primero que las filas iguales son las filas \(i\) e
\(i+1\), con \(i>1\). Al desarrollar por la primera fila,
\[
\det(A)
=
\sum_{j=1}^n
(-1)^{1+j}a_{1j}\det(A_{1j}).
\]
En cada \(A_{1j}\), las filas \(i-1\) e \(i\) son iguales y
consecutivas. Por la hipótesis de inducción,
\[
\det(A_{1j})=0
\]
para todo \(j\). Por tanto,
\[
\det(A)=0.
\]

Queda el caso en que las dos primeras filas son iguales. Escribamos
\[
A=
\begin{pmatrix}
a_1&a_2&\cdots&a_n\\
a_1&a_2&\cdots&a_n\\
\multicolumn{4}{c}{\raisebox{1ex}{\(\vdots\)}}
\end{pmatrix}.
\]

Para \(j<k\), denotemos por \(B_{jk}\) la matriz obtenida suprimiendo
las dos primeras filas y las columnas \(j\) y \(k\).

Desarrollamos primero por la primera fila y, en cada una de las
matrices obtenidas, nuevamente por la primera fila.

Para cada par \(j<k\) aparecen exactamente dos términos que contienen
el producto
\[
a_ja_k\det(B_{jk}).
\]

El que se obtiene suprimiendo primero la columna \(j\) y después la
columna \(k\) lleva el signo
\[
(-1)^{1+j}(-1)^k
=
(-1)^{j+k+1},
\]
mientras que el que se obtiene suprimiendo primero la columna \(k\)
y después la columna \(j\) lleva el signo
\[
(-1)^{1+k}(-1)^{1+j}
=
(-1)^{j+k+2}.
\]

Estos signos son opuestos. Por tanto, los términos se cancelan dos a
dos y
\[
\det(A)=0.
\]
\end{proof}

\end{frame}

\begin{frame}[allowframebreaks, t]{Intercambio de dos filas consecutivas}

\begin{proposition}
Si \(A'\) se obtiene de \(A\) intercambiando dos filas consecutivas,
entonces
\[
\det(A')=-\det(A).
\]
\end{proposition}

\begin{proof}
Sean \(u,v\in M_{1\times n}(k)\) las dos filas consecutivas. Como
la matriz
\[
\left(
\begin{array}{c}
\vdots\\
\hline
u+v\\
u+v\\
\hline
\vdots
\end{array}
\right)
\]
tiene dos filas consecutivas iguales, su determinante es \(0\).

\framebreak

Por la linealidad respecto de cada una de esas filas,
\[
\begin{aligned}
0 ={}&
\det
\left(
\begin{array}{c}
\vdots\\
\hline
u+v\\
u+v\\
\hline
\vdots
\end{array}
\right)
={}
\det
\left(
\begin{array}{c}
\vdots\\
\hline
u\\
u+v\\
\hline
\vdots
\end{array}
\right)
+
\det
\left(
\begin{array}{c}
\vdots\\
\hline
v\\
u+v\\
\hline
\vdots
\end{array}
\right)\\
={}&
\det
\left(
\begin{array}{c}
\vdots\\
\hline
u\\
u\\
\hline
\vdots
\end{array}
\right)
+
\det
\left(
\begin{array}{c}
\vdots\\
\hline
u\\
v\\
\hline
\vdots
\end{array}
\right)
+
\det
\left(
\begin{array}{c}
\vdots\\
\hline
v\\
u\\
\hline
\vdots
\end{array}
\right)
+
\det
\left(
\begin{array}{c}
\vdots\\
\hline
v\\
v\\
\hline
\vdots
\end{array}
\right).
\end{aligned}
\]

El primer y el último determinante de la última línea son \(0\). Por tanto,
\[
\det
\left(
\begin{array}{c}
\vdots\\
\hline
v\\
u\\
\hline
\vdots
\end{array}
\right)
=
-
\det
\left(
\begin{array}{c}
\vdots\\
\hline
u\\
v\\
\hline
\vdots
\end{array}
\right).
\]
\end{proof}

\end{frame}

\begin{frame}{Intercambio de dos filas}

\begin{corollary}
Si \(A'\) se obtiene de \(A\) intercambiando dos filas distintas,
entonces
\[
\det(A')=-\det(A).
\]
\end{corollary}

\vspace{-15pt}

\begin{proof}
Supongamos que se intercambian las filas \(i\) y \(j\), con \(i<j\).

Podemos llevar la fila \(i\) hasta la posición \(j\) mediante
\(j-i\) intercambios de filas consecutivas, y llevar después la fila
que inicialmente ocupaba la posición \(j\) hasta la posición \(i\)
mediante \(j-i-1\) intercambios más.

En total realizamos
\[
(j-i)+(j-i-1)=2(j-i)-1
\]
intercambios de filas consecutivas.

Cada uno cambia el signo del determinante y \(2(j-i)-1\) es impar.
Por tanto,
\[
\det(A')=-\det(A).\qedhere
\]
\end{proof}

\end{frame}

\begin{frame}{Dos filas iguales}

\begin{corollary}
Si una matriz cuadrada $A$ tiene dos filas iguales, entonces
\[
\det(A)=0.
\]
\end{corollary}

\vspace{-15pt}

\begin{proof}
Supongamos que las filas \(i\) y \(j\), con \(i<j\), son iguales. 
Si son consecutivas, el resultado ya está demostrado. Si no, sea
\(A'\) es la matriz que se obtiene de \(A\) intercambiando las filas \(i\) y \(j-1\). Entonces
\[\det(A)=-\det(A'),\]
además \(A'\) tiene dos filas consecutivas iguales, de modo que
\[\det(A')=0.
\]
En consecuencia,
\[\det(A)=0.\qedhere
\]
\end{proof}

\end{frame}

\begin{frame}[allowframebreaks]{Desarrollo por una fila cualquiera}

\begin{proposition}
Sea \(A=(a_{ij})\in M_n(k)\). Para cualquier fila \(i\),
\[
\det(A)
=
\sum_{j=1}^n
(-1)^{i+j}a_{ij}\det(A_{ij}),
\]
donde \(A_{ij}\) es la matriz que se obtiene de \(A\) suprimiendo
la fila \(i\) y la columna \(j\).
\end{proposition}

\begin{proof}
Llevamos la fila \(i\) de \(A\) a la primera posición mediante
\(i-1\) intercambios de filas consecutivas, y denotamos por \(B\)
la matriz obtenida. Entonces
\[
\det(B)=(-1)^{i-1}\det(A).
\]

La primera fila de \(B\) es
\[
(a_{i1},\ldots,a_{in}).
\]
Además, al suprimir de \(B\) su primera fila y la columna \(j\)
obtenemos precisamente \(A_{ij}\).

Por la definición del determinante,
\[
\det(B)
=
\sum_{j=1}^n
(-1)^{1+j}a_{ij}\det(A_{ij}).
\]

Por tanto,
\[
\begin{aligned}
\det(A)
&=
(-1)^{i-1}\det(B)\\
&=
\sum_{j=1}^n
(-1)^{i+j}a_{ij}\det(A_{ij}).
\end{aligned}
\]
\end{proof}

\end{frame}

\begin{frame}[allowframebreaks, t]{Determinantes y operaciones elementales}

\begin{theorem}
Sea \(A\in M_n(k)\).

\begin{enumerate}
\item Si
\(
A\xrightarrow{F_i\leftrightarrow F_j}B,
\)
entonces
\[
\det(B)=-\det(A).
\]

\item Si
\(
A\xrightarrow{\lambda F_i}B\),
\(
\lambda\neq0,
\)
entonces
\[
\det(B)=\lambda\det(A).
\]

\item Si
\(
A\xrightarrow{F_i+\lambda F_j}B\),
\(
i\neq j,
\)
entonces
\[
\det(B)=\det(A).
\]
\end{enumerate}
\end{theorem}

\framebreak 

\begin{proof}
La primera afirmación ya está demostrada.

\medskip

La segunda se sigue de la linealidad del determinante respecto de
cada fila.

\medskip

Para la tercera, si denotamos $u$ y $v$ a as filas $i$ y $j$ de la matriz $A$,
\[
A=
\left(
\begin{array}{c}
\vdots\\
\hline
u\\
\hline
\vdots\\
\hline
v\\
\hline
\vdots
\end{array}
\right),
\qquad
B=
\left(
\begin{array}{c}
\vdots\\
\hline
u+\lambda v\\
\hline
\vdots\\
\hline
v\\
\hline
\vdots
\end{array}
\right).
\]
Por linealidad respecto de la fila \(i\),
\[
\det(B)
=
\det
\left(
\begin{array}{c}
\vdots\\
\hline
u+\lambda v\\
\hline
\vdots\\
\hline
v\\
\hline
\vdots
\end{array}
\right)
=
\det
\left(
\begin{array}{c}
\vdots\\
\hline
u\\
\hline
\vdots\\
\hline
v\\
\hline
\vdots
\end{array}
\right)
+
\lambda
\det
\left(
\begin{array}{c}
\vdots\\
\hline
v\\
\hline
\vdots\\
\hline
v\\
\hline
\vdots
\end{array}
\right)
=
\det(A)+ \lambda 0
=
\det(A).
\]
\end{proof}

\end{frame}


\begin{frame}{Diagonal principal}

\begin{definition}
Sea \(A=(a_{ij})\in M_n(k)\). La \alert{diagonal principal} de \(A\)
es la formada por las entradas
\[
a_{11},a_{22},\ldots,a_{nn}.
\]
\end{definition}

\[
A=
\begin{pmatrix}
\alert{a_{11}} & a_{12} & \cdots & a_{1n}\\
a_{21} & \alert{a_{22}} & \cdots & a_{2n}\\
\vdots & \vdots & \ddots & \vdots\\
a_{n1} & a_{n2} & \cdots & \alert{a_{nn}}
\end{pmatrix}
\]

\end{frame}

\begin{frame}{Matrices triangulares superiores}

\begin{definition}
Una matriz \(A=(a_{ij})\in M_n(k)\) se denomina
\alert{triangular superior} si
\[
a_{ij}=0
\qquad\text{si }i>j.
\]

Es decir, tiene la forma
\[
A=
\begin{pmatrix}
a_{11} & a_{12} & \cdots & a_{1n}\\
0      & a_{22} & \cdots & a_{2n}\\
\vdots & \ddots & \ddots & \vdots\\
0      & \cdots & 0      & a_{nn}
\end{pmatrix}.
\]
Dicho de otro modo, todos los elementos por debajo de la diagonal principal son \(0\).
\end{definition}

\end{frame}
\begin{frame}[allowframebreaks]{Determinante de una matriz triangular superior}

\begin{proposition}
Sea
\[
A=
\begin{pmatrix}
a_{11} & a_{12} & \cdots & a_{1n}\\
0      & a_{22} & \cdots & a_{2n}\\
\vdots & \ddots & \ddots & \vdots\\
0      & \cdots & 0      & a_{nn}
\end{pmatrix}
\]
una matriz triangular superior. Entonces
\[
\det(A)=a_{11}a_{22}\cdots a_{nn}.
\]
\end{proposition}

\begin{proof}
Procedemos por inducción sobre \(n\).

Para \(n=1\), el resultado es inmediato.

Supongamos que el resultado es cierto para matrices triangulares
superiores de orden \(n-1\).

Desarrollamos el determinante de \(A\) por la última fila. Como todos
sus elementos salvo el último son \(0\),
\[
\det(A)
=
(-1)^{n+n}a_{nn}\det(A_{nn})
=
a_{nn}\det(A_{nn}).
\]

La matriz \(A_{nn}\) es triangular superior de orden \(n-1\), con
diagonal principal
\[
a_{11},a_{22},\ldots,a_{n-1,n-1}.
\]
Por la hipótesis de inducción,
\[
\det(A_{nn})
=
a_{11}a_{22}\cdots a_{n-1,n-1}.
\]

Por tanto,
\[
\det(A)
=
a_{11}a_{22}\cdots a_{nn}.
\]
\end{proof}

\end{frame}

\begin{frame}{Cálculo del determinante mediante eliminación gaussiana}

Sea \(A\in M_n(k)\). Mediante operaciones elementales de tipos \(1\)
y \(3\), obtenemos una forma escalonada \(E\) de \(A\).

Como \(A\) es cuadrada, \(E\) es triangular superior y, por tanto,
\[
\det(E)=e_{11}\cdots e_{nn}.
\]

Las operaciones de tipo \(3\) no modifican el determinante y cada
intercambio de filas cambia su signo. Por tanto, si se han realizado
\(s\) intercambios,
\[
\det(A)=(-1)^s e_{11}\cdots e_{nn}.
\]

\end{frame}

\begin{frame}{Ejemplo: cálculo del determinante por Gauss}

Sea
\[
A=
\begin{pmatrix}
0&2&1&3\\
1&2&-1&0\\
0&4&1&8\\
0&0&-1&5
\end{pmatrix}.
\]

Mediante eliminación gaussiana,
\[
A
\xrightarrow{F_1\leftrightarrow F_2}
\begin{pmatrix}
1&2&-1&0\\
0&2&1&3\\
0&4&1&8\\
0&0&-1&5
\end{pmatrix}
\xrightarrow{F_3-2F_2}
\begin{pmatrix}
1&2&-1&0\\
0&2&1&3\\
0&0&-1&2\\
0&0&-1&5
\end{pmatrix}
\xrightarrow{F_4-F_3}
\begin{pmatrix}
1&2&-1&0\\
0&2&1&3\\
0&0&-1&2\\
0&0&0&3
\end{pmatrix}.
\]

Se ha realizado un único intercambio de filas. Por tanto,
\[
\det(A)
=
-\,\bigl(1\cdot2\cdot(-1)\cdot3\bigr)
=
6.
\]

\end{frame}

\begin{frame}{Determinante de las matrices elementales}

La matriz identidad es triangular superior.

\begin{corollary}
	El determinante de la matriz identidad es $|I|=1$.
\end{corollary}

Como las matrices elementales se obtienen a partir de la identidad $I$ mediante la realización de la correspondiente operación elemental, tenemos el siguiente cálculo.

\begin{corollary}
	Los determinantes de los tres tipos de matrices elementales son:
	\begin{enumerate}
		\item $| E_{ij} |=-1$.
		\item $| E_i(\lambda) |=\lambda$.
		\item $| E_{ij}(\lambda) |=1$.
	\end{enumerate}
\end{corollary}

\end{frame}

\begin{frame}[allowframebreaks, t]{Determinante e invertibilidad}

\begin{theorem}
Sea \(A\in M_n(k)\). Entonces
\(A\) es invertible si y solo si
\(
\det(A)\neq0.
\)
\end{theorem}

\begin{proof}
Sea $E$ una forma escalonada de \(A\) obtenida mediante operaciones elementales por filas tipo $1$ y $3$. Según hemos visto $\det(A)=\pm\det(E)$, con lo cual uno de los dos es nulo si y solo si el otro lo es. Al ser $E$ cuadrada y escalonada, es triangular superior.

\medskip

Si \(A\) no es invertible, entonces \(\operatorname{rg}(A)<n\) y, por tanto $E$ tiene la última fila nula. En consecuencia $e_{nn}=0$ y al ser \(E\) triangular superior, \(\det(E)=e_{11}\cdots e_{nn}=0\).

\medskip

Si \(A\) es invertible, entonces \(\operatorname{rg}(A)=n\) y, por tanto, \(E\) no tiene filas nulas. Como cada pivote de $E$ está más a la derecha que el pivote de la fila anterior, los pivotes de $E$ han  de ser las entradas de la diagonal principal. Luego $e_{ii}\neq0$ para todo \(i=1,\ldots,n\) y $\det(E)\neq 0$.
\end{proof}

\end{frame}

\begin{frame}[allowframebreaks]{Determinante de un producto}

\begin{theorem}
Si \(A,B\in M_n(k)\), entonces
\[
\det(AB)=\det(A)\det(B).
\]
\end{theorem}

\begin{proof}
Supongamos primero que \(A\) no es invertible. 
gaussiana obtenemos una forma escalonada \(E\) de \(A\) con la última
fila nula. Aplicando a \(AB\) las mismas operaciones elementales obtenemos
\(EB\), que también tiene última fila nula. Por tanto,
\[
\det(AB)=0.
\]
Como \(A\) no es invertible, también \(\det(A)=0\), y la fórmula se
cumple.

Supongamos ahora que \(A\) es invertible. Entonces $A=E_1\cdots E_r$ es producto de matrices elementales $E_i$, $1\leq i\leq r$. En este caso, demostraremos la fórmula por inducción en $r$.

\medskip

Si $r=0$ entonces $A=I$ y $AB=B$, así que
\[
\det(AB)=\det(B)=\det(I)\det(B)=\det(A)\det(B).
\]

Si $r=1$ entonces $A=E_1$ es una matriz elemental y $AB=E_1B$ se obtiene a partir de $B$ realizando la correspondiente operación elemental por filas, por tanto la fórmula es equivalente a la variación del determinante por operaciones elementales por filas, ya demostrada.

\medskip

Si $r>1$ y suponemos el resultado cierto para productos de $r-1$ matrices elementales, entonces:
\begin{align*}
	\det(AB)&=\det((E_1E_2\cdots E_r)B)\\
	&=\det(E_1(E_2\cdots E_rB))\\
	&=\det(E_1)\det(E_2\cdots E_rB)\\
	&=\det(E_1)\det(E_2\cdots E_r)\det(B)\\
	&=\det(E_1E_2\cdots E_r)\det(B)\\
	&=\det(A)\det(B).
\end{align*}
Aquí hemos aplicado dos veces el caso $r=1$, ya probado, y una vez el caso $r-1$, que es la hipótesis de inducción.
\end{proof}

\end{frame}


\begin{frame}{Menores de una matriz}

\begin{definition}
Sea \(A\in M_{m\times n}(k)\). Elegimos \(r\) filas $i_1<\cdots<i_r$ y \(r\) columnas $j_1<\cdots<j_r$
de \(A\). La matriz formada por las entradas situadas en sus
intersecciones es una matriz cuadrada de orden \(r\).
\[
\left(
\begin{array}{ccc}
a_{i_1j_1} & \cdots & a_{i_1j_r}\\
\vdots & \ddots & \vdots\\
a_{i_rj_1} & \cdots & a_{i_rj_r}
\end{array}
\right).
\]
El determinante de esta matriz se denomina un
\alert{menor de orden \(r\)} de \(A\).
\end{definition}

\medskip

Así, los menores de orden \(1\) son las propias entradas de \(A\).

\end{frame}

\begin{frame}[allowframebreaks]{Menores y operaciones elementales}

\begin{proposition}
Sea \(1\leq r\leq\min\{m,n\}\). Si \(A,B\in M_{m\times n}(k)\)
son equivalentes por filas, entonces los siguientes enunciados son equivalentes
\begin{enumerate}
	\item Todos los menores de orden \(r\) de \(A\) son nulos.
	\item Todos los menores de orden \(r\) de \(B\) son nulos.
\end{enumerate}
\end{proposition}

\begin{proof}
Basta comprobar que la propiedad se conserva mediante una operación
elemental por filas.

Si intercambiamos dos filas, cada menor de la nueva matriz coincide,
salvo posiblemente el signo, con un menor de la matriz anterior.

Si multiplicamos una fila por \(\lambda\neq0\), cada menor que
contiene esa fila queda multiplicado por \(\lambda\), y los demás
no cambian.

Finalmente, supongamos que sustituimos
\[
F_i\quad\text{por}\quad F_i+\lambda F_j.
\]
Un menor que no contiene la fila \(i\) no cambia.

Si contiene las filas \(i\) y \(j\), por la linealidad del
determinante respecto de una fila queda igual, pues el término
adicional tiene dos filas iguales.

Si contiene la fila \(i\), pero no la fila \(j\), por linealidad se
convierte en la suma del menor original y \(\lambda\) veces, salvo
posiblemente el signo, otro menor de orden \(r\) de la matriz
original.

Por tanto, si todos los menores de orden \(r\) eran nulos, continúan
siéndolo después de cualquier operación elemental.
\end{proof}

\end{frame}

\begin{frame}[allowframebreaks]{Caracterización del rango por menores}

\begin{theorem}
Sea \(A\in M_{m\times n}(k)\) y sea
\(1\leq r\leq\min\{m,n\}\). Entonces
\[
\operatorname{rg}(A)\geq r
\]
si y solo si \(A\) posee algún menor de orden \(r\) distinto de \(0\).
\end{theorem}

\begin{proof}
Sea \(E\) una forma escalonada de \(A\). Por el resultado anterior,
\(A\) posee un menor no nulo de orden \(r\) si y solo si \(E\) lo
posee.

Sea
\[
q=\operatorname{rg}(A).
\]
Entonces \(E\) tiene exactamente \(q\) filas no nulas y \(q\)
pivotes.

Si \(q\geq r\), tomamos las primeras \(r\) filas no nulas y las
\(r\) columnas que contienen sus pivotes. La matriz cuadrada así
obtenida es triangular superior y los elementos de su diagonal
principal son precisamente esos \(r\) pivotes.

Por tanto, su determinante es distinto de \(0\), y \(E\) posee un
menor no nulo de orden \(r\).

Recíprocamente, si \(r>q\), cualquier elección de \(r\) filas de
\(E\) contiene alguna fila nula, pues \(E\) solo tiene \(q\) filas
no nulas. Por tanto, todos sus menores de orden \(r\) son nulos.

Así,
\[
\operatorname{rg}(A)\geq r
\quad\Longleftrightarrow\quad
A\text{ posee un menor no nulo de orden }r.
\]
\end{proof}

\end{frame}

\begin{frame}{Cálculo del rango mediante menores}

\begin{corollary}
Sea \(A\in M_{m\times n}(k)\). Entonces
\[
\operatorname{rg}(A)=r
\]
si y solo si:
\begin{enumerate}
\item existe algún menor de orden \(r\) distinto de \(0\);
\item todos los menores de orden \(r+1\) son \(0\).
\end{enumerate}
\end{corollary}

\medskip

Por tanto, el rango de una matriz es el \alert{mayor orden de un
menor no nulo} de la matriz.

\end{frame}


\begin{frame}{Matriz traspuesta}

\begin{definition}
Sea
\[
A=(a_{ij})=\begin{pmatrix}
a_{11}&a_{12}&\cdots&a_{1n}\\
a_{21}&a_{22}&\cdots&a_{2n}\\
\vdots&\vdots&&\vdots\\
a_{m1}&a_{m2}&\cdots&a_{mn}
\end{pmatrix}\in M_{m\times n}(k).
\]
La \alert{matriz traspuesta} de \(A\) es la matriz
\[
A^t=(a_{ij}^t)=(a_{ji})=\begin{pmatrix}
a_{11}&a_{21}&\cdots&a_{m1}\\
a_{12}&a_{22}&\cdots&a_{m2}\\
\vdots&\vdots&&\vdots\\
a_{1n}&a_{2n}&\cdots&a_{mn}
\end{pmatrix}\in M_{n\times m}(k).
\]
Es decir,
\[a^t_{ij}=a_{ji}.\]
\end{definition}

\end{frame}

\begin{frame}{Traspuesta de un producto}

\begin{proposition}
Sean
\(
A\in M_{m\times n}(k),
\) y \(
B\in M_{n\times p}(k).
\). 
Entonces
\[
(AB)^t=B^tA^t.
\]
\end{proposition}

\begin{proof}
Para \(1\leq i\leq p\) y \(1\leq j\leq m\),
\[
\begin{aligned}
\bigl((AB)^t\bigr)_{ij}
&=(AB)_{ji}=\sum_{r=1}^n a_{jr}b_{ri}\\
&=\sum_{r=1}^n b_{ri}a_{jr}=\sum_{r=1}^n b^t_{ir}a^t_{rj}=(B^tA^t)_{ij}.
\end{aligned}
\]

Por tanto, como ambos lados son matrices de orden \(p\times m\) y tienen las mismas entradas, coinciden.
\end{proof}

\end{frame}

\begin{frame}{Traspuesta de una matriz invertible}

\begin{corollary}
Sea \(A\in M_n(k)\) una matriz invertible. Entonces \(A^t\) es
invertible y
\[
(A^t)^{-1}=(A^{-1})^t.
\]
\end{corollary}

\begin{proof}
Como \(A\) es invertible,
\[
AA^{-1}=I_n
\qquad\text{y}\qquad
A^{-1}A=I_n.
\]

Trasponiendo y usando la fórmula para la traspuesta de un producto,
obtenemos
\[
(AA^{-1})^t=(A^{-1})^tA^t =I_n^t=I_n
\qquad\text{y}\qquad
(A^{-1}A)^t=A^t(A^{-1})^t =I_n^t =I_n.
\]

Por tanto, \((A^{-1})^t\) es una inversa de \(A^t\). Así, \(A^t\)
es invertible y
\[
(A^t)^{-1}=(A^{-1})^t.\qedhere
\]
\end{proof}

\end{frame}

\begin{frame}{Traspuestas de matrices elementales}

\begin{proposition}
Las traspuestas de los tres tipos de matrices elementales son:
\begin{enumerate}
	\item $E_{ij}^t=E_{ij}$,
	\item $E_i(\lambda)^t=E_i(\lambda)$,
	\item $E_{ij}(\lambda)^t=E_{ji}(\lambda)$.
\end{enumerate}
En particular, si $E$ es una matriz elemental entonces
\[
\det(E^t)=\det(E).
\]
\end{proposition}

\end{frame}

\begin{frame}[allowframebreaks]{Determinante de la matriz traspuesta}

\begin{theorem}
Para toda \(A\in M_n(k)\),
\[
\det(A^t)=\det(A).
\]
\end{theorem}

\begin{proof}
Supongamos primero que \(A\) es invertible. Como su forma escalonada
y reducida es \(I_n\), \(A\) es producto de matrices elementales:
\[
A=E_1\cdots E_r.
\]

Entonces
\[
A^t=E_r^t\cdots E_1^t.
\]
Por la fórmula del determinante de un producto y el resultado
anterior,
\[
\begin{aligned}
\det(A^t)
&=\det(E_r^t)\cdots\det(E_1^t)\\
&=\det(E_r)\cdots\det(E_1)\\
&=\det(E_1)\cdots\det(E_r)\\
&=\det(A).
\end{aligned}
\]

Supongamos ahora que \(A\) no es invertible. Entonces \(A^t\)
tampoco es invertible pues si \(A^t\) fuese invertible,
también lo sería
\[
(A^t)^t=A.
\]
En este caso \(\det(A)=0=\det(A^t)\).
\end{proof}

\end{frame}

\begin{frame}{Rango de una matriz traspuesta}

\begin{corollary}
Para toda matriz \(A\),
\[
\operatorname{rg}(A^t)=\operatorname{rg}(A).
\]
\end{corollary}

\begin{proof}
Como el determinante de una matriz coincide con el de su traspuesta, los menores de \(A\) de cada tamaño coinciden con los de \(A^t\). Por tanto, la caracterización del rango en términos de menores demuestra que \(A\) y \(A^t\) tienen el mismo rango.
\end{proof}

\end{frame}


\begin{frame}{Propiedades del determinante por columnas}

\begin{corollary}
El determinante es lineal respecto de cada columna. Más concretamente, si $u,v\in M_{n\times 1}(k)$ y $\lambda\in k$, entonces
\begin{align*}
	\det\left(\cdots\mid u+v\mid\cdots\right)
	&=
	\det\left(\cdots\mid u\mid\cdots\right)
	+ \det\left(\cdots\mid v\mid\cdots\right),\\
	\det\left(\cdots\mid \lambda u\mid\cdots\right)
	&=
	\lambda\det\left(\cdots\mid u\mid\cdots\right).
\end{align*}
Además, si una matriz cuadrada tiene dos columnas iguales, entonces
su determinante es \(0\).
\end{corollary}

\end{frame}



\begin{frame}{Operaciones elementales por columnas: intercambio}

La \alert{operación elemental tipo 1}
\[
\alert{C_i\leftrightarrow C_j}
\]
intercambia las columnas \(i\) y \(j\):

\medskip

\[
\begin{pmatrix}
\cdots & a_{1i} & \cdots & a_{1j} & \cdots\\
       & \vdots &        & \vdots &       \\
\cdots & a_{mi} & \cdots & a_{mj} & \cdots
\end{pmatrix}
\quad
\xrightarrow{\;C_i\leftrightarrow C_j\;}
\quad
\begin{pmatrix}
\cdots & a_{1j} & \cdots & a_{1i} & \cdots\\
       & \vdots &        & \vdots &       \\
\cdots & a_{mj} & \cdots & a_{mi} & \cdots
\end{pmatrix}.
\]

\end{frame}

\begin{frame}{Operaciones elementales por columnas: producto por un escalar}

La \alert{operación elemental tipo 2}
\[
\alert{\lambda C_i},
\qquad \lambda\neq0,
\]
multiplica la columna \(i\) por el escalar no nulo
\(\lambda\in k\):

\medskip

\[
\begin{pmatrix}
\cdots & a_{1i} & \cdots\\
       & \vdots &       \\
\cdots & a_{mi} & \cdots
\end{pmatrix}
\quad
\xrightarrow{\;\lambda C_i\;}
\quad
\begin{pmatrix}
\cdots & \lambda a_{1i} & \cdots\\
       & \vdots          &       \\
\cdots & \lambda a_{mi} & \cdots
\end{pmatrix}.
\]

\end{frame}

\begin{frame}{Operaciones elementales por columnas: suma de un múltiplo}

La \alert{operación elemental tipo 3}
\[
\alert{C_i+\lambda C_j},
\qquad i\neq j,
\]
sustituye la columna \(i\) por \(C_i+\lambda C_j\), donde
\(\lambda\in k\) es un escalar:

\medskip

\[
\begin{pmatrix}
\cdots & a_{1i} & \cdots & a_{1j} & \cdots\\
       & \vdots &        & \vdots &       \\
\cdots & a_{mi} & \cdots & a_{mj} & \cdots
\end{pmatrix}
\quad
\xrightarrow{\;C_i+\lambda C_j\;}
\quad
\begin{pmatrix}
\cdots & a_{1i}+\lambda a_{1j} & \cdots & a_{1j} & \cdots\\
       & \vdots                  &        & \vdots &       \\
\cdots & a_{mi}+\lambda a_{mj} & \cdots & a_{mj} & \cdots
\end{pmatrix}.
\]

\end{frame}

\begin{frame}{Operaciones elementales y trasposición}

\begin{theorem}
Sea \(A\) una matriz.

\begin{enumerate}
\item Si \(A\xrightarrow{F_i\leftrightarrow F_j}B\), entonces
\(A^t\xrightarrow{C_i\leftrightarrow C_j}B^t\).

\item Si \(A\xrightarrow{\lambda F_i}B\), con \(\lambda\neq0\), entonces
\(A^t\xrightarrow{\lambda C_i}B^t\).

\item Si \(A\xrightarrow{F_i+\lambda F_j}B\), con \(i\neq j\), entonces
\(A^t\xrightarrow{C_i+\lambda C_j}B^t\).
\end{enumerate}
\end{theorem}

\end{frame}

\begin{frame}{Operaciones por columnas y matrices elementales}

\begin{theorem}
Sea \(A\in M_{m\times n}(k)\). Entonces:

\begin{enumerate}
\item Si \(A\xrightarrow{C_i\leftrightarrow C_j}B\), entonces
\[
B=AE_{ij}.
\]

\item Si \(A\xrightarrow{\lambda C_i}B\), con \(\lambda\neq0\), entonces
\[
B=AE_i(\lambda).
\]

\item Si \(A\xrightarrow{C_i+\lambda C_j}B\), con \(i\neq j\), entonces
\[
B=AE_{ji}(\lambda).
\]
\end{enumerate}
\end{theorem}

\vspace{-10pt}

Este resultado es consecuencia de la relación entre las operaciones elementales por filas y por columnas mediante la trasposición y de las propiedades de la trasposición.

\end{frame}

\begin{frame}{Determinantes y operaciones elementales por columnas}

\begin{theorem}
Sea \(A\in M_n(k)\).

\begin{enumerate}
\item Si
\(
A\xrightarrow{C_i\leftrightarrow C_j}B,
\)
entonces
\[
\det(B)=-\det(A).
\]

\item Si
\(
A\xrightarrow{\lambda C_i}B\),
\(\lambda\neq0\),
entonces
\[
\det(B)=\lambda\det(A).
\]

\item Si
\(
A\xrightarrow{C_i+\lambda C_j}B\),
\(i\neq j\),
entonces
\[
\det(B)=\det(A).
\]
\end{enumerate}
\end{theorem}

\vspace{-10pt}

Se deduce de las propiedades análogas para operaciones elementales
por filas, aplicándolas a las matrices traspuestas y usando
\(
\det(A^t)=\det(A).
\)

\end{frame}


\begin{frame}{Equivalencia por columnas}

\begin{definition}
Sean \(A,B\in M_{m\times n}(k)\). Diremos que \(A\) y \(B\) son
\alert{equivalentes por columnas}, y escribiremos
\[
A\sim_c B,
\]
si \(B\) se obtiene a partir de \(A\) mediante una sucesión finita
de operaciones elementales por columnas.
\end{definition}

\medskip

Como trasponer transforma operaciones elementales por columnas en
las correspondientes operaciones elementales por filas,
\[
A\sim_c B
\quad\Longleftrightarrow\quad
A^t\sim_f B^t,
\]
donde \(\sim_f\) denota la equivalencia por filas.

\end{frame}


\begin{frame}{Equivalencia por columnas}

\begin{proposition}
Dos matrices $A,B\in M_{m\times n}(k)$ son equivalentes por columnas si y solo si existe una matriz invertible $Q\in M_n(k)$ tal que $AQ=B$.\end{proposition}

\begin{proof}
	Si $A$ y $B$ son equivalentes por columnas entonces $A^t$ y $B^t$ son equivalentes por filas, luego existe una matriz invertible $P\in M_n(k)$ tal que $PA^t=B^t$, por tanto
	\[(PA^t)^t=(A^t)^tP^t=AP^t=(B^t)^t=B.\]
	La matriz $Q=P^t$ es invertible por serlo $P$.
	
	\medskip
	
	Recíprocamente, si $AQ=B$, al ser $Q$ producto de matrices elementales, $B$ se obtiene a partir de $A$ realizando operaciones elementales por columnas, luego $A$ Y $B$ son equivalentes por columnas.
\end{proof}

\end{frame}

\begin{frame}{Equivalencia por filas y columnas}

\begin{definition}
Sean \(A,B\in M_{m\times n}(k)\). Diremos que \(A\) y \(B\) son
\alert{equivalentes por filas y columnas}, y escribiremos
\[
A\sim_{fc}B,
\]
si \(B\) se obtiene a partir de \(A\) mediante una sucesión finita
de operaciones elementales por filas y por columnas.
\end{definition}

\begin{proposition}
	Dos matrices $A,B\in M_{m\times n}(k)$ son equivalentes por filas y columnas si y solo si existen matrices invertibles $P\in M_m(k)$ y $Q\in M_n(k)$ tal que $PAQ=B$.
\end{proposition}


\end{frame}

\begin{frame}[allowframebreaks, t]
{Equivalencia por filas y columnas y rango}

\begin{theorem}
Dos matrices \(A,B\in M_{m\times n}(k)\) son equivalentes por filas
y columnas si y solo si
\(
\operatorname{rg}(A)=\operatorname{rg}(B).
\)
\end{theorem}


\begin{proof}
Las operaciones elementales por filas y por columnas no modifican
el rango. Por tanto, is $A$ y $B$ son equivalentes por filas y columnas
\(
\operatorname{rg}(A)=\operatorname{rg}(B).
\)

Recíprocamente, sea
\(
\operatorname{rg}(A)=r.
\)
Mediante operaciones elementales por filas llevamos \(A\) a su forma
escalonada y reducida.

Esta matriz tiene \(r\) columnas pivote. Mediante intercambios de
columnas podemos llevarlas a las primeras \(r\) posiciones. Obtenemos
así una matriz de la forma
\[
\left(
\begin{array}{c|c}
I_r&C\\
\hline
0&0
\end{array}
\right).
\]

\framebreak

Ahora anulamos las entradas de \(C\) mediante operaciones elementales
por columnas:
\[
\left(
\begin{array}{cccc|ccc}
1 & 0 & \cdots & 0
  & \cdots
  & \alert{c_{1j}}
  & \cdots\\
0 & 1 & \cdots & 0
  & \cdots
  & \alert{c_{2j}}
  & \cdots\\
\vdots & \vdots & \ddots & \vdots
  & \vdots
  & \alert{\vdots}
  & \vdots\\
0 & 0 & \cdots & 1
  & \cdots
  & \alert{c_{rj}}
  & \cdots\\
\hline
0 & 0 & \cdots & 0
  & \cdots
  & \alert{0}
  & \cdots\\
\vdots & \vdots &        & \vdots
  & \vdots
  & \alert{\vdots}
  & \vdots\\
0 & 0 & \cdots & 0
  & \cdots
  & \alert{0}
  & \cdots
\end{array}
\right)
\xrightarrow{\substack{C_j-c_{1j}C_1\\ C_j-c_{2j}C_2\\\vdots\\ C_j-c_{rj}C_r}}
\left(
\begin{array}{cccc|ccc}
1 & 0 & \cdots & 0
  & \cdots
  & \alert{0}
  & \cdots\\
0 & 1 & \cdots & 0
  & \cdots
  & \alert{0}
  & \cdots\\
\vdots & \vdots & \ddots & \vdots
  & \vdots
  & \alert{\vdots}
  & \vdots\\
0 & 0 & \cdots & 1
  & \cdots
  & \alert{0}
  & \cdots\\
\hline
0 & 0 & \cdots & 0
  & \cdots
  & \alert{0}
  & \cdots\\
\vdots & \vdots &        & \vdots
  & \vdots
  & \alert{\vdots}
  & \vdots\\
0 & 0 & \cdots & 0
  & \cdots
  & \alert{0}
  & \cdots
\end{array}
\right)
\]


De este modo obtenemos
\[
A\sim_{fc}
J_r:=
\left(
\begin{array}{c|c}
I_r&0\\
\hline
0&0
\end{array}
\right).
\]

Por tanto, toda matriz de rango \(r\) es equivalente por filas y
columnas a \(J_r\), luego dos de ellas son equivalentes entre sí.
\end{proof}

\end{frame}

\begin{frame}{Menores orlados}

\begin{definition}
Sea \(A\in M_{m\times n}(k)\) y sea
\[
\Delta = \det
\begin{pmatrix}
a_{i_1j_1} & \cdots & a_{i_1j_r}\\
\vdots      & \ddots & \vdots\\
a_{i_rj_1} & \cdots & a_{i_rj_r}
\end{pmatrix}
\]
un menor de orden \(r\) de \(A\), determinado por las filas
\(i_1<\cdots<i_r\) y las columnas \(j_1<\cdots<j_r\).

\medskip

Si \(i\notin\{i_1,\ldots,i_r\}\) y
\(j\notin\{j_1,\ldots,j_r\}\), el menor de orden \(r+1\)
determinado por las filas
\(i_1<\cdots<i <\cdots <i_r\)
y las columnas
\(j_1<\cdots<j <\cdots <j_r\)
se denomina un \alert{menor orlado} de $\Delta$.
\end{definition}

\end{frame}

\begin{frame}{Menores orlados}

Gráficamente, el menor orlado de $\Delta$ es
\[
\det\left(
\begin{array}{ccccccc}
a_{i_1j_1} & \cdots & \alert{a_{i_1j}} & \cdots & \cdots & a_{i_1j_r}\\
\vdots     &        & \alert{\vdots}   &        &        & \vdots\\
\alert{a_{ij_1}} & \alert{\cdots} & \alert{a_{ij}} & \alert{\cdots} & \alert{\cdots} & \alert{a_{ij_r}}\\
\vdots     &        & \alert{\vdots}   &        &        & \vdots\\
\vdots     &        & \alert{\vdots}   &        &        & \vdots\\
a_{i_rj_1} & \cdots & \alert{a_{i_rj}} & \cdots & \cdots & a_{i_rj_r}
\end{array}
\right)
\]

\end{frame}

\begin{frame}[allowframebreaks, t]{Método del orlado}

\begin{theorem}
Sea \(A\in M_{m\times n}(k)\) y sea \(\Delta\neq0\) un menor de
orden \(r\) de \(A\). Entonces
\[
\operatorname{rg}(A)=r
\]
si y solo si todos los menores de orden \(r+1\) que orlan a
\(\Delta\) son nulos.
\end{theorem}

\begin{proof}
Sean
\[
i_1<\cdots<i_r,
\qquad
j_1<\cdots<j_r
\]
las filas y columnas que determinan el menor
\(\Delta\neq0\).

Mediante intercambios de filas llevamos las filas
\(i_1,\ldots,i_r\) a las primeras \(r\) posiciones, y mediante
intercambios de columnas llevamos las columnas
\(j_1,\ldots,j_r\) a las primeras \(r\) posiciones.

Obtenemos así una matriz \(A'\) de la forma
\[
A'=
\left(
\begin{array}{c|c}
B & *\\
\hline
* & *
\end{array}
\right),
\]
donde \(B\in M_r(k)\) y
\[
\det(B)=\Delta\neq0.
\]

Esto no modifica el rango. Además, los menores que orlan a \(\Delta\) se corresponden con los
menores que orlan a \(\det(B)\), salvo posibles cambios de
signo debido al reordenamiento de filas y columnas. Así que podemos suponer que el menor \(\Delta\) está determinado
por las primeras \(r\) filas y las primeras \(r\) columnas.

Como \(\det(B)\neq0\), la matriz \(B\) es invertible. Mediante
operaciones elementales sobre las primeras \(r\) filas podemos
transformar \(B\) en \(I_r\). Aplicando esas mismas operaciones a
las filas completas de \(A\), obtenemos una matriz de la forma
\[
\left(
\begin{array}{c|c}
I_r & C'\\
\hline
D' & E'
\end{array}
\right).
\]

Estas operaciones no modifican el rango. Además, aplicadas a cada
menor orlado, multiplican su determinante por un escalar no nulo,
por lo que preservan su nulidad.

A continuación, utilizando las primeras \(r\) filas, anulamos las
entradas de \(D'\) mediante operaciones elementales de tipo \(3\).
Obtenemos así
\[
\left(
\begin{array}{c|c}
I_r & C'\\
\hline
0 & E''
\end{array}
\right).
\]
Estas operaciones tampoco modifican el rango ni la nulidad de los
menores orlados.

Para \(i>r\) y \(j>r\), el menor orlado correspondiente a la fila
\(i\) y la columna \(j\) es ahora
\[
\left|
\begin{array}{c|c}
I_r & *\\
\hline
0\ \cdots\ 0 & e''_{ij}
\end{array}
\right|
=e''_{ij}.
\]

Por tanto, todos los menores que orlan a \(\Delta\) son nulos si y
solo si
\(
E''=0.
\)
En tal caso, la matriz obtenida es
\[
\left(
\begin{array}{c|c}
I_r & C'\\
\hline
0 & 0
\end{array}
\right),
\]
que está escalonada y tiene exactamente \(r\) pivotes. Por tanto,
\[
\operatorname{rg}(A)=r.
\]

Recíprocamente, si \(\operatorname{rg}(A)=r\), todos los menores de
orden \(r+1\) son nulos y, en particular, todos los menores que orlan
a \(\Delta\) son nulos.

\end{proof}

\end{frame}

\begin{frame}{Cálculo del rango por el método del orlado}

Sea \(A\in M_{m\times n}(k)\).

\begin{enumerate}
\item Si \(A=0\), entonces
\(
\operatorname{rg}(A)=0.
\)

\item En caso contrario, elegimos inductivamente un menor no nulo
\[
\Delta_r \neq0
\]
de orden \(r\), empezando por $r=1$, que es simplemente elegir una entrada no nula.
\begin{enumerate}

\item Si $r=\min\{m,n\}$ entonces $\operatorname{rg}(A)=r$.

\item En caso contrario, comenzamos a calcular menores de orden \(r+1\) que orlen a
\(\Delta_r\).
\begin{enumerate}
\item Si alguno es no nulo, lo tomamos como un nuevo menor
\[
\Delta_{r+1}\neq0
\]
y repetimos el procedimiento desde \alert{2.}

\item Si todos son nulos, entonces
\[
\operatorname{rg}(A)=r.
\]
\end{enumerate}
\end{enumerate}
\end{enumerate}

\end{frame}

\begin{frame}{Adjuntos de una matriz}

\begin{definition}
Sea \(A=(a_{ij})\in M_n(k)\). El \alert{adjunto} de la entrada
\(a_{ij}\) es el escalar
\[
(-1)^{i+j}\det(A_{ij}),
\]
donde \(A_{ij}\) se obtiene suprimiendo la fila \(i\) y la columna
\(j\) de \(A\).

La \alert{matriz adjunta} de \(A\) es
\[
\operatorname{adj}(A)=((-1)^{i+j}\det(A_{ij})),
\]
es decir, aquella cuya entrada \((i,j)\) es el adjunto de \(a_{ij}\).
\end{definition}

\end{frame}



\begin{frame}{La matriz adjunta}

\begin{proposition}
Para toda \(A\in M_n(k)\),
\[
A\operatorname{adj}(A)^t=\det(A)I_n.
\]
\end{proposition}

\vspace{-15pt}

\begin{proof}
La entrada \((i,j)\) de \(A\operatorname{adj}(A)^t\) es
\[
\sum_{r=1}^n a_{ir}(-1)^{j+r}\det(A_{jr}).
\]

Si \(i=j\), esta suma es precisamente el desarrollo de
\(\det(A)\) por la fila \(i\). 

Si \(i\neq j\), la fórmula es el desarrollo por la fila \(j\) del determinante de la matriz \(B\) la matriz obtenida de
\(A\) sustituyendo la fila \(j\) por la fila \(i\). Como \(B\) tiene dos filas iguales (la fila \(i\) y la fila \(j\)),
\[
\det(B)=0.
\]
Así, las entradas diagonales de \(A\operatorname{adj}(A)^t\) son
\(\det(A)\) y las restantes son \(0\). 
\end{proof}

\end{frame}

\begin{frame}{Fórmula de la inversa por adjuntos}

\begin{corollary}
Si \(A\in M_n(k)\) es invertible, entonces
\[
A^{-1}
=
\frac{1}{\det(A)}\operatorname{adj}(A)^t.
\]
\end{corollary}

\end{frame}

\begin{frame}{Regla de Cramer}

\begin{theorem}[Regla de Cramer]
Consideremos un sistema de \(n\) ecuaciones lineales con \(n\)
incógnitas
\[
Ax=b,
\]
donde \(A\in M_n(k)\) es invertible. La única solución de este sistema es
\[x_i=\frac{
\det\left(
\begin{array}{ccccccc}
a_{11} & \cdots & a_{1,i-1} & \alert{b_1} & a_{1,i+1} & \cdots & a_{1n}\\
a_{21} & \cdots & a_{2,i-1} & \alert{b_2} & a_{2,i+1} & \cdots & a_{2n}\\
\vdots &        & \vdots    & \alert{\vdots} & \vdots    &        & \vdots\\
a_{n1} & \cdots & a_{n,i-1} & \alert{b_n} & a_{n,i+1} & \cdots & a_{nn}
\end{array}
\right)
}{\det(A)},\qquad 1\leq i\leq n.\]
\end{theorem}

\end{frame}

\begin{frame}[allowframebreaks, t]{Demostración de la regla de Cramer}

\begin{proof}
Escribamos las columnas de \(A\) como
\[
A=(C_1\mid \cdots\mid C_n),
\]
y sea
\[
v=
\begin{pmatrix}
v_1\\
\vdots\\
v_n
\end{pmatrix}
\]
la única solución del sistema \(Ax=b\).

\medskip

La igualdad \(Ax=b\) puede escribirse como
\[
v_1C_1+\cdots+v_nC_n=b.
\]

Fijemos \(i\), con \(1\leq i\leq n\). Entonces, el determinante de la matriz que aparece en el enunciado como numerador es:
\begin{multline*}
	\det
(C_1\mid \cdots\mid C_{i-1}\mid b\mid C_{i+1}\mid \cdots\mid C_n)
\\
=
\det (C_1\mid \cdots\mid C_{i-1}\mid v_1C_1+\cdots+v_nC_n\mid C_{i+1}\mid \cdots\ C_n)\\
=
\sum_{j=1}^n
v_j\det(C_1\mid \ldots\mid C_{i-1}\mid C_j\mid C_{i+1}\mid \ldots\mid C_n)\\
=
v_i\det(C_1\mid \ldots\mid C_{i-1}\mid C_i\mid C_{i+1}\mid \ldots\mid C_n)=v_i\det(A).
\end{multline*}
Aquí hemos aplicado la linealidad del determinante por columnas y que el determinante de una matriz con dos filas iguales es \(0\). 
Despejando $\det(A)\neq 0$ obtenemos la fórmula del enunciado.
\end{proof}

\end{frame}

% \input{temas/01-sistemas-matrices-determinantes/ejercicios/...}

\section{Espacios vectoriales}



\begin{frame}
    \frametitle{Espacios vectoriales}

    \[
        k^n,
        \qquad
        M_{m\times n}(k),
        \qquad
        k[x],
        \qquad
        k^X
    \]

    \begin{center}
        \alert{¿Qué tienen en común estos conjuntos?}
    \end{center}

    \medskip

    \begin{alertblock}{}
        En los primeros ejemplos puede pensarse en \(k=\mathbb{R}\).
    \end{alertblock}
\end{frame}


\begin{frame}
    \frametitle{El ejemplo prototípico: \(k^n\)}

    Dados
    \[
        u=(u_1,\ldots,u_n),
        \qquad
        v=(v_1,\ldots,v_n)\in k^n,
    \]
    y \(\lambda\in k\), definimos
    \[
        u+v
        =
        (u_1+v_1,\ldots,u_n+v_n)
    \]
    y
    \[
        \lambda u
        =
        (\lambda u_1,\ldots,\lambda u_n).
    \]

    Estas operaciones satisfacen, entre otras, las propiedades
    \[
        u+v=v+u,
        \qquad
        (u+v)+w=u+(v+w),
    \]
    \[
        \lambda(u+v)=\lambda u+\lambda v,
        \qquad
        (\lambda+\mu)u=\lambda u+\mu u.
    \]
\end{frame}


\begin{frame}
    \frametitle{Otros ejemplos}

    Las mismas operaciones pueden definirse en conjuntos de naturaleza muy distinta:

    \begin{itemize}
        \item las matrices
        \[
            M_{m\times n}(k),
        \]

        \item los polinomios
        \[
            k[x],
        \]

        \item las funciones de un conjunto \(X\) en \(k\)
        \[
            k^X=\{f\colon X\longrightarrow k\}.
        \]
    \end{itemize}

    En todos estos casos tenemos una suma y un producto por escalares
    que satisfacen las mismas propiedades formales que en \(k^n\).

    \medskip

    \begin{center}
        \alert{La naturaleza de los elementos cambia, pero la estructura es la misma.}
    \end{center}
\end{frame}

\begin{frame}
    \frametitle{Espacios vectoriales}

    \begin{definition}
        Sea \(k\) un cuerpo. Un \alert{espacio vectorial sobre \(k\)} es un conjunto \(V\)
        dotado de dos operaciones
        \begin{alignat*}{2}
            +\colon V\times V &\longrightarrow V, \qquad &(u,v) &\longmapsto u+v,\\
            \cdot\;\colon k\times V &\longrightarrow V, \qquad &(\lambda,v) &\longmapsto \lambda v,
        \end{alignat*}
        llamadas, respectivamente, \alert{suma} y \alert{producto por escalares},
        que satisfacen ciertos axiomas.
    \end{definition}

    \medskip

    Los elementos de \(V\) se llaman \alert{vectores}.
\end{frame}


\begin{frame}
    \frametitle{Axiomas de un espacio vectorial}

    Para cualesquiera \(u,v,w\in V\),

    \begin{align*}
        (u+v)+w &= u+(v+w),\\
        u+v &= v+u,
    \end{align*}

    existe un elemento \(0\in V\) llamado \alert{vector cero}, \alert{nulo} o \alert{trivial} tal que
    \[
        v+0=v,
    \]

    y para cada \(v\in V\) existe un elemento \(-v\in V\) tal que
    \[
        v+(-v)=0.
    \]

    \medskip

    Es decir, \((V,+)\) es un \alert{grupo abeliano}.
\end{frame}


\begin{frame}
    \frametitle{Axiomas de un espacio vectorial}

    Para cualesquiera \(\lambda,\mu\in k\) y \(u,v\in V\),

    \begin{align*}
        \lambda(u+v) &= \lambda u+\lambda v,\\
        (\lambda+\mu)v &= \lambda v+\mu v,\\
        (\lambda\mu)v &= \lambda(\mu v),\\
        1v &= v.
    \end{align*}

    \medskip

    Obsérvese que utilizamos el mismo símbolo \(0\) para denotar
    el elemento neutro de \(k\) y el vector cero de \(V\):
    \[
        0\in k,
        \qquad
        0\in V.
    \]
\end{frame}


\begin{frame}
    \frametitle{El espacio vectorial \(k^n\)}

    \begin{example}
        El conjunto
        \[
            k^n=\{(u_1,\ldots,u_n):u_i\in k\}
        \]
        es un espacio vectorial sobre \(k\) con las operaciones
        \[
            (u_1,\ldots,u_n)+(v_1,\ldots,v_n)
            =
            (u_1+v_1,\ldots,u_n+v_n),
        \]
        \[
            \lambda(u_1,\ldots,u_n)
            =
            (\lambda u_1,\ldots,\lambda u_n).
        \]
    \end{example}

    \medskip

    Su vector cero es
    \[
        (0,\ldots,0),
    \]
    y el opuesto de \((u_1,\ldots,u_n)\) es
    \[
        (-u_1,\ldots,-u_n).
    \]
\end{frame}


\begin{frame}
    \frametitle{El espacio vectorial de las matrices}

    \begin{example}
        El conjunto
        \[
            M_{m\times n}(k)
        \]
        de matrices \(m\times n\) con coeficientes en \(k\) es un espacio vectorial
        sobre \(k\), con la suma de matrices y el producto por escalares habituales.
    \end{example}

    \medskip

    El vector cero es la matriz nula
    \[
        0=
        \begin{pmatrix}
            0 & \cdots & 0\\
            \vdots & & \vdots\\
            0 & \cdots & 0
        \end{pmatrix},
    \]
    y el opuesto de una matriz \(A\) es \(-A\).
\end{frame}


\begin{frame}
    \frametitle{El espacio vectorial de los polinomios}

    \begin{example}
        El conjunto
        \[
            k[x]
        \]
        de polinomios con coeficientes en \(k\) es un espacio vectorial sobre \(k\),
        con la suma y el producto por escalares habituales.
    \end{example}

    \medskip

    También lo es el conjunto
    \[
        k_n[x]
        =
        \{p(x)\in k[x]:\deg p\leq n\}\cup\{0\}.
    \]

    \medskip

    Por ejemplo,
    \[
        1,\quad x,\quad x^2,\quad \ldots,\quad x^n
    \]
    son vectores de \(k_n[x]\).
\end{frame}


\begin{frame}
    \frametitle{Espacios vectoriales de funciones}

    \begin{example}
        Sea \(X\) un conjunto. El conjunto
        \[
            k^X=\{f\colon X\longrightarrow k\}
        \]
        es un espacio vectorial sobre \(k\), definiendo
        \[
            (f+g)(x)=f(x)+g(x),
        \]
        \[
            (\lambda f)(x)=\lambda f(x),
        \]
        para todo \(x\in X\).
    \end{example}

    \medskip

    El vector cero es la función constante
    \[
        0\colon X\longrightarrow k,
        \qquad
        x\longmapsto 0.
    \]
\end{frame}

\begin{frame}
    \frametitle{Un ejemplo menos evidente}

    \begin{example}
        Consideremos
        \[
            V=\mathbb{R}_{>0}
        \]
        y definamos la \alert{suma} y el \alert{producto por escalares} del siguiente modo inusual:
        \[
            x\oplus y=xy,
            \qquad
            \lambda\odot x=x^\lambda,
        \]
        para \(x,y\in V\) y \(\lambda\in\mathbb{R}\). Entonces \(V\) es un espacio vectorial sobre \(\mathbb{R}\). En este espacio vectorial, el vector cero es
    \[
        1,
    \]
    pues
    \[
        x\oplus 1=x.
    \]

    Además, el opuesto de \(x\) es
    \[
        x^{-1}.
    \]
        \end{example}
\end{frame}


\begin{frame}
    \frametitle{Más ejemplos}

    Observa que los siguientes conjuntos, con las operaciones habituales,
    son espacios vectoriales sobre \(k\):

    \begin{enumerate}
        \item
        \(
            \{(x_1,\ldots,x_n)\in k^n:
            x_1+\cdots+x_n=0\}.
        \)

        \item
        \(
            \{p(x)\in k[x]:p(0)=0\}.
        \)
    \end{enumerate}

    \medskip

	Es un buen ejercicio verificar que los axiomas se cumplen en cada caso.
\end{frame}

\begin{frame}
    \frametitle{Resta de vectores}

    Sean $u,v\in V$. Definimos la \alert{resta} de $v$ a $u$ por
    \[
        u-v:=u+(-v).
    \]

    Es decir, restar un vector es, por definición, sumar su opuesto.

    \begin{proposition}
        Para todo $u,v,w\in V$,
        \[
            u-v=w \quad\Longleftrightarrow\quad u=w+v.
        \]
    \end{proposition}

    \begin{proof}
        ($\Rightarrow$) $w+v=(u-v)+v=u+((-v)+v)=u$.

        ($\Leftarrow$) $u-v=u+(-v)=(w+v)+(-v)=w+(v+(-v))=w$.
    \end{proof}
\end{frame}


\begin{frame}
    \frametitle{Unicidad del vector cero}

    \begin{proposition}
        Sea \(V\) un espacio vectorial sobre \(k\). El vector cero de \(V\) es único.
    \end{proposition}

    \begin{proof}
        Supongamos que \(0,0'\in V\) son dos vectores cero. Entonces
        \[
            0=0+0'=0'.
        \]
        Por tanto, el vector cero es único.
    \end{proof}
\end{frame}

\begin{frame}
    \frametitle{Unicidad del opuesto}

    \begin{proposition}
        Sea \(V\) un espacio vectorial sobre \(k\). Para cada \(v\in V\),
        existe un único vector \(-v\in V\) tal que
        \[
            v+(-v)=0.
        \]
    \end{proposition}

    \begin{proof}
        La existencia forma parte de los axiomas de espacio vectorial.

        Supongamos que \(w,w'\in V\) satisfacen
        \[
            v+w=0,
            \qquad
            v+w'=0.
        \]
        Entonces
        \[
            w
            =
            w+0
            =
            w+(v+w')
            =
            (w+v)+w'
            =
            (v+w)+w'
            =
            0+w'
            =
            w'+0
            =
            w'.
        \]

        Por tanto, el opuesto de \(v\) es único.
    \end{proof}
\end{frame}

\begin{frame}[allowframebreaks, t]
    \frametitle{Primeras consecuencias de los axiomas}

    \begin{proposition}
        Sea \(V\) un espacio vectorial sobre \(k\). Para todo \(v\in V\)
        y todo \(\lambda\in k\),
        \[
            0v=0,
            \qquad
            \lambda 0=0,
            \qquad
            (-1)v=-v.
        \]
    \end{proposition}

    \begin{proof}
        Como
        \[
            0v=(0+0)v=0v+0v,
        \]
        despejando obtenemos
        \[
            0v=0.
        \]

        Análogamente,
        \[
            \lambda 0
            =
            \lambda(0+0)
            =
            \lambda 0+\lambda 0,
        \]
        y por tanto
        \[
            \lambda 0=0.
        \]

        Por otra parte,
        \[
            v+(-1)v
            =
            1v+(-1)v
            =
            (1-1)v
            =
            0v
            =
            0.
        \]
        Por unicidad del opuesto,
        \[
            (-1)v=-v.
        \]
    \end{proof}
\end{frame}


\begin{frame}
    \frametitle{Operaciones que no definen un espacio vectorial}

    En \(\mathbb{R}^2\), consideremos la suma habitual y definamos
    \[
        \lambda\cdot(x,y)=(\lambda x,y).
    \]

    \medskip

    ¿Definen estas operaciones una estructura de espacio vectorial sobre \(\mathbb{R}\)?

    \medskip

    No. Por ejemplo,
    \[
        0\cdot(x,y)=(0,y),
    \]
    que en general no es el vector cero.

    \medskip

    Por tanto, falla una consecuencia necesaria de los axiomas:
    \[
        0v=0.
    \]
\end{frame}

\begin{frame}
    \frametitle{Despejar un escalar no nulo}

    \begin{proposition}
        Sean $\lambda\in k$ no nulo y $u,v\in V$. Entonces
        \[
            \lambda u=v \quad\Longleftrightarrow\quad u=\lambda^{-1}v.
        \]
    \end{proposition}

    \begin{proof}
        ($\Rightarrow$) $u=1\cdot u=(\lambda^{-1}\lambda)u=\lambda^{-1}(\lambda u)=\lambda^{-1}v$.

        ($\Leftarrow$) $\lambda u=\lambda(\lambda^{-1}v)=(\lambda\lambda^{-1})v=1\cdot v=v$.
    \end{proof}
\end{frame}


\begin{frame}
    \frametitle{Cuando el producto por escalares es cero}

    \begin{corollary}
        Sea \(V\) un espacio vectorial sobre \(k\). Si \(\lambda\in k\) y \(v\in V\),
        entonces
        \[
            \lambda v=0
            \quad\Longrightarrow\quad
            \lambda=0
            \quad\text{o}\quad
            v=0.
        \]
    \end{corollary}

    \begin{proof}
        Si \(\lambda\neq 0\), entonces 
        \[v=\lambda^{-1}0=0.\]
    \end{proof}
\end{frame}


\subsection{Dependencia lineal}

\begin{frame}
    \frametitle{Combinaciones lineales}

    \begin{definition}
        Sean \(v_1,\ldots,v_r\) vectores de un espacio vectorial \(V\) sobre \(k\).
        Una \alert{combinación lineal} de \(v_1,\ldots,v_r\) es un vector de la forma
        \[
            \lambda_1v_1+\cdots+\lambda_rv_r,
        \]
        donde
        \[
            \lambda_1,\ldots,\lambda_r\in k.
        \]

        Los escalares \(\lambda_1,\ldots,\lambda_r\) se llaman
        \alert{coeficientes} de la combinación lineal.
    \end{definition}

    En particular, el vector cero siempre es combinación lineal de
    \(v_1,\ldots,v_r\):
    \[
        0=0v_1+\cdots+0v_r.
    \]
\end{frame}


\begin{frame}
    \frametitle{Combinaciones lineales en \(k^n\)}

    \begin{example}
        En \(\mathbb{R}^3\), sean
        \[
            v_1=(1,0,2),
            \qquad
            v_2=(0,1,-1).
        \]
        Entonces
        \[
            2v_1-3v_2
            =
            2(1,0,2)-3(0,1,-1)
            =
            (2,-3,7).
        \]

        Por tanto, \((2,-3,7)\) es combinación lineal de \(v_1\) y \(v_2\).
    \end{example}

    \medskip

    En general, determinar si \(v\) es combinación lineal de
    \(v_1,\ldots,v_r\) equivale a estudiar si existen escalares
    \(\lambda_1,\ldots,\lambda_r\) tales que
    \[
        \lambda_1v_1+\cdots+\lambda_rv_r=v.
    \]
\end{frame}

\begin{frame}
    \frametitle{¿Es combinación lineal?}

        Consideremos en \(\mathbb{R}^3\)
        \[
            v_1=(1,1,0),
            \qquad
            v_2=(0,1,1),
            \qquad
            w=(2,3,1).
        \]

        Para decidir si \(w\) es combinación lineal de \(v_1,v_2\),
        buscamos \(\lambda_1,\lambda_2\in\mathbb{R}\) tales que
        \[
            \lambda_1v_1+\lambda_2v_2=(\lambda_1, \lambda_1+ \lambda_2, \lambda_2)=w.
        \]
        Esto equivale al sistema
        \[
            \begin{cases}
                \lambda_1=2,\\
                \lambda_1+\lambda_2=3,\\
                \lambda_2=1.
            \end{cases}
        \]
        Este sistema es compatible. Su única solución es $(\lambda_1,\lambda_2)=(2,1)$, por tanto,
        \[
            v=2v_1+v_2.
        \]

\end{frame}



\begin{frame}
    \frametitle{Combinaciones lineales y sistemas de ecuaciones}

    \begin{proposition}
        Sean $v_1,\ldots,v_r,w\in k^n$ y \[A=(v_1\mid\cdots\mid v_r).\]
        El vector $w$ es combinación lineal de $v_1,\ldots,v_r$ si y solo si
        el sistema \[Ax=w\] es compatible.
        En tal caso, toda solución $(\lambda_1,\ldots,\lambda_r)$ 
        verifica \[w=\lambda_1v_1+\cdots+\lambda_rv_r.\]
    \end{proposition}
    
    \vspace{-10pt}

    \begin{proof}
        Por definición, $w$ es combinación lineal de $v_1,\ldots,v_r$
        si y solo si existen $x_1,\ldots,x_r\in k$ tales que
        $x_1v_1+\cdots+x_rv_r=w$.
        Esta ecuación es equivalente a $Ax=w$,
        pues $Ax=x_1v_1+\cdots+x_rv_r$.
    \end{proof}
\end{frame}


\begin{frame}
    \frametitle{Combinaciones lineales de polinomios}

    \begin{example}
        En \(\mathbb{R}_2[x]\), consideremos
        \[
            p_1(x)=1+x,
            \qquad
            p_2(x)=x+x^2.
        \]
        Entonces
        \[
            2p_1(x)-p_2(x)
            =
            2+x-x^2.
        \]

        Por tanto,
        \[
            2+x-x^2
        \]
        es combinación lineal de \(p_1(x)\) y \(p_2(x)\).
    \end{example}

\end{frame}


\begin{frame}
    \frametitle{Dependencia lineal}

    \begin{definition}
        Un conjunto de vectores \(\{v_1,\ldots,v_r\}\) de un espacio vectorial \(V\)
        es \alert{linealmente dependiente} si existen
        \(\lambda_1,\ldots,\lambda_r\in k\), 
        \begin{center}  
			\textbf{NO TODOS NULOS}, 
        \end{center}
        tales que
        \[
            \lambda_1v_1+\cdots+\lambda_rv_r=0.
        \]

        Una igualdad de este tipo con coeficientes no todos nulos se llama
        \alert{relación de dependencia lineal}.
    \end{definition}

    \medskip

    La condición ``no todos nulos'' es esencial, pues siempre tenemos
    \[
        0v_1+\cdots+0v_r=0.
    \]
\end{frame}


\begin{frame}
    \frametitle{Un ejemplo de dependencia lineal}

    \begin{example}
        Los vectores
        \[
            v_1=(1,0,1),
            \qquad
            v_2=(0,1,1),
            \qquad
            v_3=(1,1,2)
        \]
        de \(\mathbb{R}^3\) son linealmente dependientes, pues
        \[
            v_1+v_2-v_3=0.
        \]

        Hemos encontrado una relación de dependencia lineal con coeficientes
        \[
            1,\quad 1,\quad -1.
        \]
    \end{example}

    Obsérvese que, equivalentemente,
    \[
        v_3=v_1+v_2.
    \]
\end{frame}


\begin{frame}
    \frametitle{Independencia lineal}
    
    La independencia lineal se define como lo contrario de la dependencia lineal y se puede caracterizar del siguiente modo.

    \begin{definition}
        Un conjunto de vectores \(\{v_1,\ldots,v_r\}\) de un espacio vectorial \(V\)
        es \alert{linealmente independientes} si
        \[
            \lambda_1v_1+\cdots+\lambda_rv_r=0
        \]
        implica
        \[
            \lambda_1=\cdots=\lambda_r=0.
        \]
    \end{definition}

    \medskip

    Es decir, \(v_1,\ldots,v_r\) son linealmente independientes
    si la única combinación lineal de ellos que da el vector cero
    es la que tiene todos sus coeficientes nulos.
\end{frame}


\begin{frame}
    \frametitle{Un ejemplo de independencia lineal}

        Los vectores
        \[
            v_1=(1,1),
            \qquad
            v_2=(1,-1)
        \]
        de \(\mathbb{R}^2\) son linealmente independientes.

        En efecto, si
        \[
            \lambda_1v_1+\lambda_2v_2=0,
        \]
        entonces
        \[
            (\lambda_1+\lambda_2,\lambda_1-\lambda_2)=(0,0),
        \]
        de donde
        \[
        \begin{cases}
        	\lambda_1+\lambda_2=0,\\
            \lambda_1-\lambda_2=0.
        \end{cases}
        \]
		Este sistema homogéneo es compatible determinado, por lo que necesariamente,
        \[
            \lambda_1=\lambda_2=0.
        \]
\end{frame}

\begin{frame}[allowframebreaks]
    \frametitle{Un ejemplo de dependencia lineal}


        Los vectores
        \[
            v_1=(1,2,3),\qquad v_2=(1,1,1),\qquad v_3=(3,4,5)
        \]
        de $\mathbb{R}^3$ son linealmente dependientes.

        Buscamos $\lambda_1,\lambda_2,\lambda_3\in\mathbb{R}$, no todos nulos,
        tales que $\lambda_1v_1+\lambda_2v_2+\lambda_3v_3=0$. Operando,
        \[
            \lambda_1(1,2,3)+\lambda_2(1,1,1)+\lambda_3(3,4,5)
            =
            (\lambda_1+\lambda_2+3\lambda_3,\;
             2\lambda_1+\lambda_2+4\lambda_3,\;
             3\lambda_1+\lambda_2+5\lambda_3),
        \]
        de modo que la ecuación equivale al sistema homogéneo
        \[
            \begin{cases}
                \lambda_1+\lambda_2+3\lambda_3=0,\\
                2\lambda_1+\lambda_2+4\lambda_3=0,\\
                3\lambda_1+\lambda_2+5\lambda_3=0.
            \end{cases}
        \]
        Escalonamos la matriz de coeficientes:
        \[
            \begin{pmatrix}
                1&1&3\\
                2&1&4\\
                3&1&5
            \end{pmatrix}
            \xrightarrow{\;F_2-2F_1\;}
            \begin{pmatrix}
                1&1&3\\
                0&-1&-2\\
                3&1&5
            \end{pmatrix}
            \xrightarrow{\;F_3-3F_1\;}
            \begin{pmatrix}
                1&1&3\\
                0&-1&-2\\
                0&-2&-4
            \end{pmatrix}
        \]
        \[
            \xrightarrow{\;F_3-2F_2\;}
            \begin{pmatrix}
                1&1&3\\
                0&-1&-2\\
                0&0&0
            \end{pmatrix}.
        \]
        Como $\rg=2<3$, el sistema tiene soluciones no triviales.
        La única variable libre es $\lambda_3$. Tomando $\lambda_3=-1$:
        de la segunda fila de la forma escalonada obtenermos $\lambda_2=2$;
        in entonces de la primera, $\lambda_1=1$.

        Por tanto,
        \[
            v_1+2v_2-v_3=(1,2,3)+2(1,1,1)-(3,4,5)=(0,0,0)
        \]
        es una relación de dependencia lineal.

\end{frame}



\begin{frame}
    \frametitle{Independencia lineal y sistemas homogéneos}

    \begin{proposition}
        Sean $v_1,\ldots,v_r\in k^n$ y $A=(v_1\mid\cdots\mid v_r)$.
        Los vectores $v_1,\ldots,v_r$ son linealmente independientes
        si y solo si el sistema homogéneo $Ax=0$ es determinado,
        lo que equivale a $\rg(A)=r$.
    \end{proposition}

    \begin{proof}
        Por definición, $v_1,\ldots,v_r$ son linealmente independientes
        si y solo si la única solución de $x_1v_1+\cdots+x_rv_r=0$
        es $x_1=\cdots=x_r=0$.
        Esta ecuación coincide con $Ax=0$, pues $Ax=x_1v_1+\cdots+x_rv_r$.
        El sistema $Ax=0$ es determinado si y solo $\rg(A)=r$.
    \end{proof}

    Si $\rg(A)<r$, los vectores son linealmente dependientes
    y $Ax=0$ tiene soluciones no triviales.
    Toda solución no trivial $(\lambda_1,\ldots,\lambda_r)$ da lugar a una 
    relación de dependencia lineal \[\lambda_1v_1+\cdots+\lambda_rv_r=0.\]
\end{frame}



\begin{frame}
    \frametitle{Independencia lineal mediante eliminación gaussiana}

    \begin{example}
        Estudiemos la independencia lineal de
        \[
            v_1=(1,1,0),\qquad
            v_2=(1,0,1),\qquad
            v_3=(0,1,1).
        \]

        Por la proposición anterior, basta calcular el rango de la matriz
        \(A=(v_1\mid v_2\mid v_3)\), que tiene los vectores por columnas.
        Escalonando,
        \[
            A=
            \begin{pmatrix}
                1&1&0\\
                1&0&1\\
                0&1&1
            \end{pmatrix}
            \xrightarrow{\;F_2-F_1\;}
            \begin{pmatrix}
                1&1&0\\
                0&-1&1\\
                0&1&1
            \end{pmatrix}
            \xrightarrow{\;F_3+F_2\;}
            \begin{pmatrix}
                1&1&0\\
                0&-1&1\\
                0&0&2
            \end{pmatrix}.
        \]

        Hay tres pivotes, luego \(\rg(A)=3\), que es el número de vectores.
        Por tanto, \(v_1,v_2,v_3\) son linealmente independientes.
    \end{example}
\end{frame}





\begin{frame}
    \frametitle{Un solo vector}

    \begin{proposition}
        Un conjunto unitario \(\{v\}\subseteq V\) es linealmente independiente si y solo si
        \[
            v\neq 0.
        \]
    \end{proposition}

    \begin{proof}
        Si \(v=0\), entonces
        \[
            1v=0,
        \]
        es una relación de dependencia lineal.

        Supongamos ahora que \(v\neq0\). Si
        \[
            \lambda v=0,
        \]
        entonces \(\lambda=0\).
    \end{proof}
\end{frame}


\begin{frame}
    \frametitle{Familias que contienen el vector cero}

    \begin{proposition}
        Si uno de los vectores \(v_1,\ldots,v_r\) es el vector cero,
        entonces \(v_1,\ldots,v_r\) son linealmente dependientes.
    \end{proposition}

    \begin{proof}
        Supongamos, por ejemplo, que \(v_i=0\). Entonces
        \[
            0v_1+\cdots+0v_{i-1}
            +1v_i
            +0v_{i+1}+\cdots+0v_r
            =0
        \]
        es una relación de dependencia lineal.
    \end{proof}
\end{frame}


\begin{frame}
    \frametitle{Dos vectores}

    \begin{proposition}
        Un conjunto formado por dos vectores \(\{u,v\}\in V\) es linealmente dependiente
        si y solo si uno es múltiplo escalar del otro.
    \end{proposition}

    \begin{proof}
        Si \(v=\lambda u\), entonces
        \[
            \lambda u-v=0
        \]
        es una relación de dependencia lineal.

        Recíprocamente, si \(u,v\) son linealmente dependientes, existen
        \(\lambda,\mu\in k\), no ambos nulos, tales que
        \[
            \lambda u+\mu v=0.
        \]
        Si, por ejemplo, \(\mu\neq0\), entonces
        \[
            v=-\mu^{-1}\lambda u.
        \]
        El caso \(\lambda\neq0\) es análogo.
    \end{proof}
\end{frame}


\begin{frame}
    \frametitle{Dependencia lineal de polinomios}

    \begin{example}
        Consideremos en \(\mathbb{R}_2[x]\)
        \[
            p_1(x)=1+x,
            \qquad
            p_2(x)=x+x^2,
            \qquad
            p_3(x)=1+2x+x^2.
        \]

        Como
        \[
            p_1(x)+p_2(x)-p_3(x)=0,
        \]
        los polinomios \(p_1,p_2,p_3\) son linealmente dependientes.
    \end{example}

    \medskip

    De nuevo, la dependencia lineal no es una propiedad exclusiva de
    vectores de \(k^n\).
\end{frame}


\begin{frame}[allowframebreaks]
    \frametitle{Dependencia y combinaciones lineales}

    \begin{theorem}
        Sean \(v_1,\ldots,v_r\) vectores de \(V\), con \(r\geq2\).
        Son linealmente dependientes si y solo si alguno de ellos
        es combinación lineal de los que le preceden.
    \end{theorem}

    \begin{proof}
        Supongamos primero que, para algún \(i\),
        \[
            v_i
            =
            \lambda_1v_1+\cdots+\lambda_{i-1}v_{i-1}.
        \]
        Entonces
        \[
            \lambda_1v_1+\cdots+\lambda_{i-1}v_{i-1}-v_i
            +0v_{i+1}+\cdots+0v_r=0
        \]
        es una relación de dependencia lineal.

        Recíprocamente, supongamos que \(v_1,\ldots,v_r\) son
        linealmente dependientes. Existe una relación
        \[
            \lambda_1v_1+\cdots+\lambda_rv_r=0
        \]
        con coeficientes no todos nulos.

        Sea \(i\) el mayor índice tal que \(\lambda_i\neq0\).
        Entonces
        \[
            \lambda_1v_1+\cdots+\lambda_iv_i=0,
        \]
        y, como \(\lambda_i\neq0\),
        \[
            v_i
            =
            -\lambda_i^{-1}
            \bigl(
                \lambda_1v_1+\cdots+\lambda_{i-1}v_{i-1}
            \bigr).
        \]
        Por tanto, \(v_i\) es combinación lineal de
        \(v_1,\ldots,v_{i-1}\).
    \end{proof}
\end{frame}


\begin{frame}
    \frametitle{Subfamilias e independencia lineal}

    \begin{proposition}
        Toda subfamilia de una familia linealmente independiente
        es linealmente independiente.
    \end{proposition}

    \begin{proof}
        Si una subfamilia fuese linealmente dependiente, existiría una
        combinación lineal no trivial de sus vectores igual a cero.

        Añadiendo coeficientes nulos para los restantes vectores,
        obtendríamos una relación de dependencia lineal de la familia
        original.
    \end{proof}

    \medskip

    Equivalentemente, toda familia que contiene una subfamilia
    linealmente dependiente es linealmente dependiente.
\end{frame}


\subsection{Sistemas generadores y bases}

\begin{frame}
    \frametitle{Sistemas generadores}

    \begin{definition}
        Sean \(v_1,\ldots,v_r\) vectores de un espacio vectorial \(V\).
        Decimos que \(v_1,\ldots,v_r\) \alert{generan} \(V\), o que
        \[
            \{v_1,\ldots,v_r\}
        \]
        es un \alert{sistema de generadores} de \(V\), si todo vector
        \(v\in V\) es combinación lineal de \(v_1,\ldots,v_r\).

        Es decir, para cada \(v\in V\) existen
        \(\lambda_1,\ldots,\lambda_r\in k\) tales que
        \[
            v=\lambda_1v_1+\cdots+\lambda_rv_r.
        \]
    \end{definition}
\end{frame}


\begin{frame}
    \frametitle{Un sistema de generadores de \(k^n\)}

    Consideremos los vectores
    \[
        e_1=(1,0,\ldots,0),\quad
        e_2=(0,1,\ldots,0),\quad\ldots,\quad
        e_n=(0,0,\ldots,1).
    \]

    Si
    \[
        v=(v_1,\ldots,v_n)\in k^n,
    \]
    entonces
    \[
        v=v_1e_1+\cdots+v_ne_n.
    \]

    Por tanto,
    \[
        e_1,\ldots,e_n
    \]
    generan \(k^n\).

    \medskip

    Estos vectores se llaman \alert{vectores canónicos} de \(k^n\).
\end{frame}


\begin{frame}
    \frametitle{Un sistema generador puede ser redundante}

    \begin{example}
        Los vectores
        \[
            v_1=(1,0),\qquad
            v_2=(0,1),\qquad
            v_3=(1,1)
        \]
        generan \(\mathbb{R}^2\).

        En efecto, si
        \[
            v=(v_1,v_2)\in\mathbb{R}^2,
        \]
        entonces
        \[
            v=v_1(1,0)+v_2(0,1).
        \]

        Sin embargo,
        \[
            v_3=v_1+v_2,
        \]
        de modo que \(v_3\) no es necesario para generar
        \(\mathbb{R}^2\).
    \end{example}

    \medskip

    Un sistema de generadores no tiene por qué ser linealmente independiente.
\end{frame}


\begin{frame}
    \frametitle{Generadores de espacios de polinomios}

    \begin{example}
        Los polinomios
        \[
            1,x,x^2,\ldots,x^n
        \]
        generan \(k_n[x]\).

        En efecto, cualquier
        \[
            p(x)=a_0+a_1x+\cdots+a_nx^n
        \]
        puede escribirse como
        \[
            p(x)
            =
            a_0\cdot1+a_1x+\cdots+a_nx^n.
        \]
    \end{example}

    \begin{example}
        También
        \[
            1,\quad 1+x,\quad x,\quad x^2
        \]
        generan \(k_2[x]\), pero uno de estos polinomios es redundante.
    \end{example}
\end{frame}


\begin{frame}
    \frametitle{Generadores de espacios de matrices}

    \begin{example}
        Las matrices
        \[
            E_{11}=
            \begin{pmatrix}
                1&0\\
                0&0
            \end{pmatrix},
            \quad
            E_{12}=
            \begin{pmatrix}
                0&1\\
                0&0
            \end{pmatrix},
        \]
        \[
            E_{21}=
            \begin{pmatrix}
                0&0\\
                1&0
            \end{pmatrix},
            \quad
            E_{22}=
            \begin{pmatrix}
                0&0\\
                0&1
            \end{pmatrix}
        \]
        generan \(M_{2\times2}(k)\).

        En efecto,
        \[
            \begin{pmatrix}
                a&b\\
                c&d
            \end{pmatrix}
            =
            aE_{11}+bE_{12}+cE_{21}+dE_{22}.
        \]
    \end{example}
\end{frame}


\begin{frame}
    \frametitle{Generadores y sistemas de ecuaciones}

    Sean
    \[
        v_1,\ldots,v_r\in k^n
    \]
    y sea
    \[
        A=
        \begin{pmatrix}
            \vert & & \vert\\
            v_1 & \cdots & v_r\\
            \vert & & \vert
        \end{pmatrix}.
    \]

    Los vectores \(v_1,\ldots,v_r\) generan \(k^n\) si y solo si,
    para todo \(u\in k^n\), el sistema
    \[
        A
        \begin{pmatrix}
            \lambda_1\\
            \vdots\\
            \lambda_r
        \end{pmatrix}
        =u
    \]
    es compatible.

    \medskip

    Es decir, el problema de determinar si unos vectores generan \(k^n\)
    vuelve a ser un problema de sistemas de ecuaciones lineales.
\end{frame}


\begin{frame}
    \frametitle{Criterio de rango para sistemas generadores}

    \begin{theorem}
        Sean \(v_1,\ldots,v_r\in k^n\) y consideramos la matriz de tamaño \(n\times r\)
        \[
            A=
            \begin{pmatrix}
                \vert & & \vert\\
                v_1 & \cdots & v_r\\
                \vert & & \vert
            \end{pmatrix}.
        \]
        Entonces \(v_1,\ldots,v_r\) generan \(k^n\) si y solo si
        \[
            \operatorname{rg}(A)=n.
        \]
    \end{theorem}

\end{frame}


\begin{frame}[allowframebreaks,t]
    \frametitle{Criterio de rango para sistemas generadores: demostración}

    \begin{proof}
        Sabemos que \(v_1,\ldots,v_r\) generan \(k^n\) si y solo si el sistema
        \(A\lambda=u\) es compatible para todo \(u\in k^n\).

        Sea \(A'\) una forma escalonada de \(A\). Por el Tema~1, existe una
        matriz invertible \(P\in M_n(k)\), producto de matrices elementales,
        tal que
        \[
            PA=A'.
        \]
        Para cada \(u\in k^n\), las mismas operaciones elementales
        transforman la matriz ampliada \((A\mid u)\) en
        \[
            P(A\mid u)=(PA\mid Pu)=(A'\mid Pu).
        \]
        Como las matrices ampliadas son equivalentes por filas, los sistemas
        \[
            A\lambda=u
            \qquad\text{y}\qquad
            A'\lambda=Pu
        \]
        tienen las mismas soluciones.

        \framebreak

        \alert{Si \(\operatorname{rg}(A)=n\), los vectores generan \(k^n\).}
        El rango es el número de filas no nulas de \(A'\), y \(A'\) tiene
        \(n\) filas, así que todas son no nulas: cada fila de \(A'\) tiene
        su pivote.

        Sea \(u\in k^n\) arbitrario. La matriz \((A'\mid Pu)\) está escalonada,
        y el primer elemento no nulo de cada fila está en la parte \(A'\).
        Por tanto, el sistema \(A\lambda=u\) es compatible.

        \framebreak

        \alert{Si \(\operatorname{rg}(A)<n\), los vectores no generan \(k^n\).}
        Ahora \(A'\) tiene menos de \(n\) filas no nulas. Como en una matriz
        escalonada las filas nulas están al final, así que si el rango de $A$ es $r<n$, la fila $r+1$ de $A'$ es nula.

        Buscamos \(u\) tal que el sistema escalonado tenga término
        independiente \(e_{r+1}=(0,\ldots,0,1, 0,\ldots,0)\). Como \(P\) es invertible,
        basta tomar
        \[
            u=P^{-1}e_{r+1},
            \qquad\text{pues entonces}\qquad
            Pu=e_{r+1}.
        \]
        La matriz\((A'\mid e_{r+1})\) está escalonada y tiene un pivote en la última columna, concretamente en la fila \(r+1\), que es:
        \[
            (0\ \cdots\ 0\mid 1).
        \]
        Por tanto \(A\lambda=u\)
        es incompatible. Así, \(u\) no es combinación lineal de \(v_1,\ldots,v_r\), y estos
        vectores no generan \(k^n\).
    \end{proof}
\end{frame}


\begin{frame}
    \frametitle{Una consecuencia inmediata}

    \begin{corollary}
        Si \(v_1,\ldots,v_r\) generan \(k^n\), entonces
        \[
            r\geq n.
        \]
    \end{corollary}

    \begin{proof}
        Si generan \(k^n\), la matriz
        \[
            A=
            \begin{pmatrix}
                \vert & & \vert\\
                v_1 & \cdots & v_r\\
                \vert & & \vert
            \end{pmatrix}
        \]
        tiene rango \(n\).

        Como una matriz con \(r\) columnas tiene rango a lo sumo \(r\),
        \[
            n=\operatorname{rg}(A)\leq r.
        \]
    \end{proof}

    \medskip

    Por tanto, no podemos generar \(k^n\) con menos de \(n\) vectores.
\end{frame}


\begin{frame}[allowframebreaks,t]
    \frametitle{Eliminar un generador redundante}

    \begin{proposition}
        Supongamos que \(v_1,\ldots,v_r\) generan \(V\) y que uno de ellos,
        por ejemplo \(v_r\), es combinación lineal de los demás:
        \[
            v_r=\lambda_1v_1+\cdots+\lambda_{r-1}v_{r-1}.
        \]
        Entonces
        \[
            v_1,\ldots,v_{r-1}
        \]
        también generan \(V\).
    \end{proposition}

    \begin{proof}
        Hay que ver que todo vector \(u\in V\) es combinación lineal de
        \(v_1,\ldots,v_{r-1}\).

        Sea \(u\in V\). Como \(v_1,\ldots,v_r\) generan \(V\), existen
        escalares \(\mu_1,\ldots,\mu_r\in k\) tales que
        \[
            u=\mu_1v_1+\cdots+\mu_{r-1}v_{r-1}+\mu_rv_r.
        \]
        En esta expresión aparece \(v_r\), que queremos eliminar.
        Sustituimos \(v_r\) por su expresión como combinación lineal
        de los demás:
        \[
            u=\mu_1v_1+\cdots+\mu_{r-1}v_{r-1}
            +\mu_r\bigl(\lambda_1v_1+\cdots+\lambda_{r-1}v_{r-1}\bigr).
        \]

        \framebreak

        Por la propiedad distributiva,
        \[
            \mu_r\bigl(\lambda_1v_1+\cdots+\lambda_{r-1}v_{r-1}\bigr)
            =(\mu_r\lambda_1)v_1+\cdots+(\mu_r\lambda_{r-1})v_{r-1}.
        \]
        Reordenando los sumandos (la suma es asociativa y conmutativa)
        y agrupando los que contienen el mismo \(v_i\), de nuevo por la
        propiedad distributiva,
        \[
            u=(\mu_1+\mu_r\lambda_1)v_1+\cdots
            +(\mu_{r-1}+\mu_r\lambda_{r-1})v_{r-1}.
        \]
        Los coeficientes \(\mu_i+\mu_r\lambda_i\) son escalares de \(k\),
        así que \(u\) es combinación lineal de \(v_1,\ldots,v_{r-1}\).

        Como \(u\in V\) era arbitrario, \(v_1,\ldots,v_{r-1}\) generan \(V\).
    \end{proof}
\end{frame}


\begin{frame}
    \frametitle{Reducir un sistema de generadores}

    \begin{example}
        Consideremos en \(\mathbb{R}^3\) los vectores
        \[
            v_1=(1,0,1),\quad
            v_2=(0,1,1),\quad
            v_3=(1,1,2),\quad
            v_4=(1,-1,0),\quad
            v_5=(0,0,1).
        \]
        \alert{Generan \(\mathbb{R}^3\).} Por el criterio de rango, basta
        escalonar la matriz que los tiene por columnas:
        \[
            \begin{pmatrix}
                1&0&1&1&0\\
                0&1&1&-1&0\\
                1&1&2&0&1
            \end{pmatrix}
            \xrightarrow{\;F_3-F_1\;}
            \begin{pmatrix}
                1&0&1&1&0\\
                0&1&1&-1&0\\
                0&1&1&-1&1
            \end{pmatrix}
            \xrightarrow{\;F_3-F_2\;}
            \begin{pmatrix}
                1&0&1&1&0\\
                0&1&1&-1&0\\
                0&0&0&0&1
            \end{pmatrix}.
        \]
        El rango es \(3\), así que \(v_1,\ldots,v_5\) generan \(\mathbb{R}^3\).
    \end{example}
\end{frame}


\begin{frame}
    \frametitle{Reducir un sistema de generadores}

    \begin{example}
        \alert{Eliminamos generadores redundantes} de
        \(v_1,\ldots,v_5\). Se observa que
        \[
            v_3=v_1+v_2,
            \qquad
            v_4=v_1-v_2.
        \]
        \begin{itemize}
            \item Como \(v_3\) es combinación lineal de los demás, por la
                  proposición anterior \(v_1,v_2,v_4,v_5\) generan \(\mathbb{R}^3\).
            \item Como \(v_4\) es combinación lineal de \(v_1,v_2,v_5\),
                  aplicando de nuevo la proposición, \(v_1,v_2,v_5\)
                  generan \(\mathbb{R}^3\).
        \end{itemize}

        No podemos eliminar ninguno más: hemos visto que para generar
        \(\mathbb{R}^3\) hacen falta al menos \(3\) vectores.
    \end{example}

    \medskip

    Este ejemplo muestra que un sistema de generadores puede reducirse
    eliminando los vectores que son combinación lineal de los demás.
\end{frame}


\begin{frame}
    \frametitle{Bases}

    Un sistema generador puede contener vectores redundantes.
    La independencia lineal expresa precisamente la ausencia de esa
    redundancia.

    \begin{definition}
        Una \alert{base} de un espacio vectorial \(V\) es un conjunto
        de vectores de \(V\) que

        \begin{enumerate}
            \item genera \(V\);
            \item es linealmente independiente.
        \end{enumerate}
    \end{definition}

    \medskip

    Por tanto, una base es un sistema de generadores en el que
    ningún vector sobra.
\end{frame}


\begin{frame}
    \frametitle{La base canónica de \(k^n\)}

    \begin{example}
        Los vectores
        \[
            e_1=(1,0,\ldots,0),\quad\ldots,\quad
            e_n=(0,\ldots,0,1)
        \]
        forman una base de \(k^n\).

        Ya sabemos que generan \(k^n\).

        Además, si
        \[
            \lambda_1e_1+\cdots+\lambda_ne_n=0,
        \]
        entonces
        \[
            (\lambda_1,\ldots,\lambda_n)=(0,\ldots,0),
        \]
        y, por tanto,
        \[
            \lambda_1=\cdots=\lambda_n=0.
        \]
    \end{example}

    Esta base se llama \alert{base canónica} de \(k^n\).
\end{frame}



\begin{frame}[allowframebreaks, t]
    \frametitle{Base de \(M_{m\times n}(k)\)}

    \begin{proposition}
        Para \(1\leq i\leq m\) y \(1\leq j\leq n\), sea \(E_{ij}\) la matriz
        \(m\times n\) cuya entrada \((i,j)\) es \(1\) y las demás son \(0\).
        Las \(mn\) matrices
        \[
            E_{ij},\quad 1\leq i\leq m,\quad 1\leq j\leq n,
        \]
        forman una base de \(M_{m\times n}(k)\).
    \end{proposition}

    \begin{proof}
        Toda \(A=(a_{ij})\in M_{m\times n}(k)\) satisface
        \[
            A=\sum_{i,j}a_{ij}E_{ij},
        \]
        pues la entrada \((i,j)\) de ambos lados es \(a_{ij}\).
        Por tanto las \(E_{ij}\) \alert{generan} \(M_{m\times n}(k)\).

        Para la \alert{independencia lineal}, si
        \[
            \sum_{i,j}\lambda_{ij}E_{ij}=0,
        \]
        la entrada \((i,j)\) del lado izquierdo es \(\lambda_{ij}\),
        luego \(\lambda_{ij}=0\) para todo \(i,j\).
    \end{proof}
\end{frame}


\begin{frame}
    \frametitle{Base de \(k_n[x]\)}

    \begin{proposition}
        Los monomios \(1,x,x^2,\ldots,x^n\) forman una base de \(k_n[x]\).
    \end{proposition}

    \begin{proof}
        Todo \(p(x)=a_0+a_1x+\cdots+a_nx^n\in k_n[x]\) se escribe
        \[
            p(x)=a_0\cdot1+a_1\cdot x+\cdots+a_n\cdot x^n,
        \]
        luego \(1,x,\ldots,x^n\) \alert{generan} \(k_n[x]\).

        Para la \alert{independencia lineal}, si
        \[
            \lambda_0\cdot1+\lambda_1 x+\cdots+\lambda_n x^n=0,
        \]
        el lado izquierdo es el polinomio nulo. El coeficiente de
        \(x^i\) en el polinomio nulo es \(0\) para todo \(i\),
        luego \(\lambda_i=0\) para todo \(0\leq i\leq n\).
    \end{proof}
\end{frame}


\begin{frame}[allowframebreaks, t]
    \frametitle{Base de \(k^X\) para \(X\) finito}

    \begin{proposition}
        Sea \(X=\{x_1,\ldots,x_n\}\) un conjunto finito.
        Para cada \(1\leq i\leq n\), sea \(\delta_i\colon X\to k\) la función
        \[
            \delta_i(x_j)=
            \begin{cases}
                1 & \text{si }j=i,\\
                0 & \text{si }j\neq i.
            \end{cases}
        \]
        Las funciones \(\delta_1,\ldots,\delta_n\) forman una base de \(k^X\).
    \end{proposition}

    \begin{proof}
        Dada \(f\colon X\to k\), la combinación lineal
        \(f(x_1)\delta_1+\cdots+f(x_n)\delta_n\)
        toma en \(x_j\) el valor
        \[
            f(x_j)\underbrace{\delta_j(x_j)}_{=\,1}
            \;+\!\sum_{i\neq j}f(x_i)\underbrace{\delta_i(x_j)}_{=\,0}
            \;=f(x_j).
        \]
        Por tanto coincide con \(f\), y las \(\delta_i\) \alert{generan} \(k^X\).

        Si \(\lambda_1\delta_1+\cdots+\lambda_n\delta_n=0\),
        evaluando en \(x_j\) obtenemos \(\lambda_j=0\) para todo \(j\),
        luego son \alert{linealmente independientes}.
    \end{proof}
\end{frame}



\begin{frame}
    \frametitle{Bases de \(k^n\) y rango}

    Sean
    \[
        v_1,\ldots,v_r\in k^n
    \]
    y sea \(A\) la matriz que tiene estos vectores como columnas.

    Para que formen una base de \(k^n\) necesitamos simultáneamente
    \[
        \operatorname{rg}(A)=r
    \]
    por independencia lineal, y
    \[
        \operatorname{rg}(A)=n
    \]
    porque han de generar \(k^n\).

    Por tanto:

    \begin{corollary}
        Una base de \(k^n\) formada por un número finito de vectores
        tiene exactamente \(n\) elementos.
    \end{corollary}
\end{frame}


\begin{frame}
    \frametitle{Criterios para ser una base de \(k^n\)}

\begin{proposition}
    Sean
    \(
        v_1,\ldots,v_n\in k^n
    \)
    y
    \[
        A=
        \begin{pmatrix}
            \vert & & \vert\\
            v_1 & \cdots & v_n\\
            \vert & & \vert
        \end{pmatrix}.
    \]

    Son equivalentes:

    \begin{enumerate}
        \item \(v_1,\ldots,v_n\) forman una base de \(k^n\);
        \item \(v_1,\ldots,v_n\) son linealmente independientes;
        \item \(v_1,\ldots,v_n\) generan \(k^n\);
        \item \(\operatorname{rg}(A)=n\);
        \item \(A\) es invertible;
        \item \(\det(A)\neq0\).
    \end{enumerate}
\end{proposition}
\end{frame}


\begin{frame}
    \frametitle{Ejemplo}

    Consideremos en \(\mathbb{R}^3\)
    \[
        v_1=(1,1,0),\qquad
        v_2=(1,0,1),\qquad
        v_3=(0,1,1).
    \]

    La matriz que tiene estos vectores como columnas es
    \[
        A=
        \begin{pmatrix}
            1&1&0\\
            1&0&1\\
            0&1&1
        \end{pmatrix}.
    \]

    Como
    \[
        \det(A)=-2\neq0,
    \]
    los vectores
    \[
        v_1,v_2,v_3
    \]
    forman una base de \(\mathbb{R}^3\).
\end{frame}


\begin{frame}
    \frametitle{Añadir un vector a una familia independiente}

    \begin{lemma}
        Sean \(v_1,\ldots,v_r\in V\) linealmente independientes y sea
        \(u\in V\) un vector que \alert{no} es combinación lineal de
        \(v_1,\ldots,v_r\). Entonces \(v_1,\ldots,v_r,u\) son linealmente
        independientes.
    \end{lemma}

    \begin{proof}
        Supongamos que
        \[
            \lambda_1v_1+\cdots+\lambda_rv_r+\mu u=0.
        \]
        Si fuera \(\mu\neq0\), podríamos despejar
        \[
            u=-\mu^{-1}\bigl(\lambda_1v_1+\cdots+\lambda_rv_r\bigr),
        \]
        y \(u\) sería combinación lineal de \(v_1,\ldots,v_r\), en contra
        de la hipótesis. Luego \(\mu=0\), y queda
        \(\lambda_1v_1+\cdots+\lambda_rv_r=0\). Por la independencia de
        \(v_1,\ldots,v_r\), \(\lambda_1=\cdots=\lambda_r=0\).
    \end{proof}
\end{frame}


\begin{frame}[allowframebreaks,t]
    \frametitle{Existencia de bases}

    \begin{theorem}
        Sea \(V\) un espacio vectorial y sea
        \[
            S=\{v_1,\ldots,v_r,v_{r+1},\ldots,v_{r+s}\}
        \]
        un sistema generador de \(V\) tal que \(v_1,\ldots,v_r\) son
        linealmente independientes (se admite \(r=0\) o \(s=0\)). Entonces \(S\)
        contiene una base \(B\) de \(V\) tal que
        \(\{v_1,\ldots,v_r\}\subseteq B\).
    \end{theorem}

    \medskip

    El caso \(r=0\) se traduce en lo siguiente.

    \begin{corollary}
        Todo espacio vectorial con un sistema generador finito tiene una base. Es más, se puede encontrar una base contenida en dicho sistema generador.
    \end{corollary}

    \framebreak

    \begin{proof}[Demostración del teorema]
        La idea es recorrer \(v_{r+1},\ldots,v_{r+s}\) en orden y añadir
        a \(v_1,\ldots,v_r\) cada uno que no sea combinación lineal de los
        ya elegidos.

        \medskip

        Empezamos con \(F_0=\{v_1,\ldots,v_r\}\), que es linealmente
        independiente. Para \(j=1,\ldots,s\), construimos \(F_j\) a partir
        de \(F_{j-1}\):
        \begin{itemize}
            \item si \(v_{r+j}\) es combinación lineal de los vectores de
                  \(F_{j-1}\), tomamos \(F_j=F_{j-1}\);
            \item si no lo es, tomamos \(F_j=F_{j-1}\cup\{v_{r+j}\}\), que es
                  linealmente independiente por el lema anterior.
        \end{itemize}
        (Si \(F_{j-1}\) es vacía, «ser combinación lineal» significa
        ser \(0\), y un único vector no nulo es independiente.)

        \framebreak

        Así, \(B=F_s\) es linealmente independiente y
        \(\{v_1,\ldots,v_r\}\subseteq B\subseteq S\). Falta ver que $B$ 
        \alert{genera} \(V\).

        \medskip

        Cada vector de \(S\) es combinación lineal de los vectores de
        \(B\): para \(i\leq r\), porque \(v_i\in B\); y para \(v_{r+j}\),
        bien \(v_{r+j}\in F_j\subseteq B\), o bien es combinación lineal
        de los de \(F_{j-1}\subseteq B\) (completando con coeficientes
        nulos). Por eliminación de generadores redundantes, \(B\) genera \(V\).
    \end{proof}
\end{frame}


\begin{frame}[allowframebreaks,t]
    \frametitle{Existencia de bases en \(k^n\): el método}

    Sean \(v_1,\ldots,v_r,v_{r+1},\ldots,v_{r+s}\in k^n\) vectores que
    generan \(k^n\), con \(v_1,\ldots,v_r\) linealmente independientes.

    \medskip

    \alert{Método.}
    \begin{enumerate}
        \item Formamos la matriz
              \(A=(v_1\mid\cdots\mid v_r\mid v_{r+1}\mid\cdots\mid v_{r+s})\).
        \item La escalonamos por filas hasta una forma escalonada \(E\).
        \item Nos quedamos con los \(v_j\) tales que la columna \(j\) de
              \(E\) contiene un pivote.
    \end{enumerate}
    Estos vectores forman una base de \(k^n\) que contiene a
    \(v_1,\ldots,v_r\).

    \framebreak

    \alert{Por qué funciona.} Sean \(j_1<\cdots<j_m\) las columnas de
    \(E\) con pivote.
    \begin{itemize}
        \item \alert{Son independientes.} Las mismas operaciones
              elementales transforman \((v_{j_1}\mid\cdots\mid v_{j_m})\) en
              las columnas \(j_1,\ldots,j_m\) de \(E\), una matriz
              escalonada con \(m\) pivotes. Su rango es \(m\).
        \item \alert{Son una base.} Como \(v_1,\ldots,v_{r+s}\) generan
              \(k^n\), \(m=\rg(A)=n\). Son \(n\) vectores independientes
              en \(k^n\).
        \item \alert{Contienen a \(v_1,\ldots,v_r\).} Las \(j-1\) primeras
              columnas de \(E\) son una forma escalonada de
              \((v_1\mid\cdots\mid v_{j-1})\). Por tanto, la columna \(j\)
              no tiene pivote si y solo si el sistema
              \((v_1\mid\cdots\mid v_{j-1})\,x=v_j\) es compatible, es decir,
              si y solo si \(v_j\) es combinación lineal de
              \(v_1,\ldots,v_{j-1}\). Como \(v_1,\ldots,v_r\) son
              independientes, ninguno es combinación lineal de los
              anteriores, y las columnas \(1,\ldots,r\) tienen pivote.
    \end{itemize}
\end{frame}

\begin{frame}
    \frametitle{Ejemplo: una base de \(\mathbb{R}^4\) entre unos generadores}

    \begin{example}
        En \(\mathbb{R}^4\) consideramos
        \[
            \begin{aligned}
                v_1&=(1,0,1,0), & v_2&=(0,1,0,1), & v_3&=(1,1,1,1),\\
                v_4&=(1,0,0,0), & v_5&=(2,0,1,0), & v_6&=(0,0,0,1).
            \end{aligned}
        \]
        donde \(v_1,v_2\) son independientes por no ser proporcionales.
        Escalonamos la matriz con columnas \(v_1,\ldots,v_6\):
        \[
            \begin{pmatrix}
                1&0&1&1&2&0\\
                0&1&1&0&0&0\\
                1&0&1&0&1&0\\
                0&1&1&0&0&1
            \end{pmatrix}
            \xrightarrow[\;F_4-F_2\;]{\;F_3-F_1\;}
            \begin{pmatrix}
                \smash{\overset{\raisebox{1.5ex}{$\scriptstyle\alert{1}$}}{\alert{1}}}&\smash{\overset{\raisebox{1.5ex}{$\scriptstyle\alert{2}$}}{0}}&1&\smash{\overset{\raisebox{1.5ex}{$\scriptstyle\alert{4}$}}{1}}&2&\smash{\overset{\raisebox{1.5ex}{$\scriptstyle\alert{6}$}}{0}}\\
                0&\alert{1}&1&0&0&0\\
                0&0&0&\alert{-1}&-1&0\\
                0&0&0&0&0&\alert{1}
            \end{pmatrix}.
        \]
        El rango es \(4\), luego los \(v_j\) generan \(\mathbb{R}^4\).
        Los pivotes están en las columnas \(1,2,4,6\):
        \(\{v_{\alert{1}},v_{\alert{2}},v_{\alert{4}},v_{\alert{6}}\}\) es una base que contiene a \(\{v_1,v_2\}\).
    \end{example}
\end{frame}



\subsection{Dimensión}

\begin{frame}[allowframebreaks, t]
    \frametitle{Bases y representación única}

    \begin{theorem}
        Sea $B=\{v_1,\ldots,v_n\}$ un conjunto de vectores de $V$.
        Es una base de $V$ si y solo si para todo $v\in V$ existen
        \alert{únicos} $\lambda_1,\ldots,\lambda_n\in k$ tales que
        \[
            v=\lambda_1v_1+\cdots+\lambda_nv_n.
        \]
    \end{theorem}

    \begin{proof}
        ($\Rightarrow$)
        Que $B$ genera $V$ garantiza la existencia.
        Para la unicidad, si
        \[
            \lambda_1v_1+\cdots+\lambda_nv_n
            =v
            =\mu_1v_1+\cdots+\mu_nv_n,
        \]
        restando se obtiene
        $(\lambda_1-\mu_1)v_1+\cdots+(\lambda_n-\mu_n)v_n=0$,
        y la independencia lineal de $B$ fuerza $\lambda_i=\mu_i$
        para todo~$i$.

        ($\Leftarrow$)
        La existencia de representación para todo $v\in V$ es la generación.
        Para la independencia lineal, si
        $\lambda_1v_1+\cdots+\lambda_nv_n=0$,
        la representación trivial $0\cdot v_1+\cdots+0\cdot v_n=0$
        y la unicidad fuerzan $\lambda_i=0$ para todo~$i$.
    \end{proof}
\end{frame}


\begin{frame}
    \frametitle{Coordenadas y aplicación de coordenadas}

    \begin{definition}
        Sea $B=\{v_1,\ldots,v_n\}$ una base de $V$.
        Los únicos $\lambda_1,\ldots,\lambda_n\in k$ tales que
        $v=\lambda_1v_1+\cdots+\lambda_nv_n$
        se llaman \alert{coordenadas} de $v$ en $B$.
        El vector columna
        \[
            v_B=
            \begin{pmatrix}
                \lambda_1\\\vdots\\\lambda_n
            \end{pmatrix}
            \in k^n
        \]
        es el \alert{vector de coordenadas} de $v$ en $B$,
        y la aplicación
        \[
            c_B\colon V\longrightarrow k^n,\quad v\longmapsto v_B,
        \]
        es la \alert{aplicación de coordenadas} respecto de $B$.
    \end{definition}

    \medskip

    En la base canónica $C=\{e_1,\ldots,e_n\}$ de $k^n$,
    el vector de coordenadas de $v=(v_1,\ldots,v_n)$
    son sus propias entradas: $v_C=(v_1,\ldots,v_n) $.
\end{frame}


\begin{frame}
    \frametitle{Propiedades de la aplicación de coordenadas}

    \begin{proposition}
        Sea $B$ una base de $V$.
        Para todo $u,v\in V$ y $\lambda\in k$,
        \[
            (u+v)_B=u_B+v_B,
            \qquad
            (\lambda v)_B=\lambda\,v_B.
        \]
    \end{proposition}

    \begin{proof}
        Sean $B=\{v_1,\ldots,v_n\}$, $u_B=(\alpha_i)$ y $v_B=(\beta_i)$.
        Entonces
        \[
            u+v=\sum_i(\alpha_i+\beta_i)v_i,
        \]
        y por unicidad de las coordenadas,
        $(u+v)_B=(\alpha_i+\beta_i)=u_B+v_B$.
        La segunda identidad es análoga.
    \end{proof}

    \medskip

    La aplicación $c_B$ convierte la suma y el producto por escalares
    de $V$ en las operaciones de $k^n$.
\end{frame}


\begin{frame}
    \frametitle{Trabajar en coordenadas: la idea}

    En \(\mathbb{R}_2[x]\), con la base \(B=\{1,x,x^2\}\), cada polinomio
    queda determinado por sus coordenadas:
    \[
        a+bx+cx^2
        \;\longmapsto\;
        \begin{pmatrix}a\\b\\c\end{pmatrix}\in\mathbb{R}^3.
    \]

    \medskip

    Para estudiar si unos polinomios son linealmente independientes,
    si generan \(\mathbb{R}_2[x]\) o si forman una base, parece razonable
    \alert{sustituirlos por sus vectores de coordenadas} y usar en
    \(\mathbb{R}^3\) la eliminación gaussiana y el rango.

    \medskip

    Para que esto sea legítimo hace falta que \(c_B\colon V\to k^n\)
    \begin{itemize}
        \item no pierda información: \alert{es biyectiva};
        \item traduzca cada combinación lineal en \(V\) en la misma
              combinación lineal en \(k^n\), y viceversa.
    \end{itemize}
    Es lo que probamos a continuación.
\end{frame}


\begin{frame}[allowframebreaks,t]
    \frametitle{La aplicación de coordenadas es biyectiva}

    \begin{proposition}
        Sea \(B=\{v_1,\ldots,v_n\}\) una base de \(V\). La aplicación
        de coordenadas \(c_B\colon V\to k^n\) es biyectiva, y su inversa es
        \[
            (\lambda_1,\ldots,\lambda_n)
            \longmapsto
            \lambda_1v_1+\cdots+\lambda_nv_n.
        \]
    \end{proposition}

    \begin{proof}
        \alert{Inyectiva.} Sean \(u,w\in V\) con \(u_B=w_B\), y llamemos
        \(\lambda_1,\ldots,\lambda_n\) a sus entradas. Por definición de
        coordenadas,
        \[
            u=\lambda_1v_1+\cdots+\lambda_nv_n=w.
        \]

        \alert{Sobreyectiva.} Sea
        \((\lambda_1,\ldots,\lambda_n) \in k^n\) y sea
        \[
            u=\lambda_1v_1+\cdots+\lambda_nv_n\in V.
        \]
        Las coordenadas de \(u\) en \(B\) son los \alert{únicos} escalares
        con los que \(u\) se escribe como combinación lineal de
        \(v_1,\ldots,v_n\). Como \(\lambda_1,\ldots,\lambda_n\) son unos
        de esos escalares, \(u_B=(\lambda_1,\ldots,\lambda_n) \).

        \framebreak

        \alert{Inversa.} El argumento anterior muestra que la aplicación
        \((\lambda_1,\ldots,\lambda_n) \mapsto\lambda_1v_1+\cdots+\lambda_nv_n\)
        seguida de \(c_B\) es la identidad de \(k^n\). Y si \(u\in V\)
        tiene \(u_B=(\lambda_1,\ldots,\lambda_n) \), entonces
        \(\lambda_1v_1+\cdots+\lambda_nv_n=u\), así que
        \(c_B\) seguida de esa aplicación es la identidad de \(V\).
    \end{proof}

    \medskip

    En particular, \(0_B=0\): el vector nulo de \(V\) es el único vector
    cuyas coordenadas son todas nulas, pues \(0=0v_1+\cdots+0v_n\).
\end{frame}


\begin{frame}[allowframebreaks,t]
    \frametitle{Coordenadas de una combinación lineal}

    \begin{proposition}
        Sea \(B\) una base de \(V\) y sean \(u_1,\ldots,u_r\in V\) y
        \(\lambda_1,\ldots,\lambda_r\in k\). Entonces
        \[
            (\lambda_1u_1+\cdots+\lambda_ru_r)_B
            =\lambda_1(u_1)_B+\cdots+\lambda_r(u_r)_B.
        \]
    \end{proposition}

    \begin{proof}
        Por inducción sobre \(r\). Para \(r=1\) es la propiedad
        \((\lambda u)_B=\lambda u_B\) ya demostrada.

        Si \(r\geq2\) y el resultado es cierto para \(r-1\) vectores,
        escribimos
        \(\lambda_1u_1+\cdots+\lambda_ru_r
        =(\lambda_1u_1+\cdots+\lambda_{r-1}u_{r-1})+\lambda_ru_r\).
        Usando primero \((u+w)_B=u_B+w_B\), después la hipótesis de
        inducción y \((\lambda u)_B=\lambda u_B\),
        \begin{align*}
            (\lambda_1u_1+\cdots+\lambda_ru_r)_B
            &=(\lambda_1u_1+\cdots+\lambda_{r-1}u_{r-1})_B+(\lambda_ru_r)_B\\
            &=\lambda_1(u_1)_B+\cdots+\lambda_{r-1}(u_{r-1})_B
              +\lambda_r(u_r)_B.\qedhere
        \end{align*}
    \end{proof}

    \medskip

    Decimos que \(c_B\) \alert{preserva} las combinaciones lineales:
    las coordenadas de una combinación lineal son la misma combinación
    lineal, \alert{con los mismos coeficientes}, de las coordenadas.
\end{frame}


\begin{frame}
    \frametitle{Las coordenadas también reflejan combinaciones lineales}

    \begin{corollary}
        Sea \(B\) una base de \(V\) y sean \(v,u_1,\ldots,u_r\in V\) y
        \(\lambda_1,\ldots,\lambda_r\in k\). Entonces
        \[
            v=\lambda_1u_1+\cdots+\lambda_ru_r
            \iff
            v_B=\lambda_1(u_1)_B+\cdots+\lambda_r(u_r)_B.
        \]
    \end{corollary}

    \begin{proof}
        (\(\Rightarrow\)) Se toman coordenadas y se aplica la proposición
        anterior.

        (\(\Leftarrow\)) Por la proposición anterior, la hipótesis dice que
        \[
            v_B=(\lambda_1u_1+\cdots+\lambda_ru_r)_B.
        \]
        Como \(c_B\) es inyectiva, \(v=\lambda_1u_1+\cdots+\lambda_ru_r\).
    \end{proof}

    \medskip

    Así, una igualdad entre vectores de \(V\) es cierta si y solo si lo es
    la igualdad correspondiente entre sus coordenadas. Decimos que \(c_B\)
    \alert{refleja} las combinaciones lineales.
\end{frame}


\begin{frame}
    \frametitle{Traducción a coordenadas}

    \begin{theorem}
        Sea \(B\) una base de \(V\) con \(n\) vectores y sean
        \(u_1,\ldots,u_r\in V\). Entonces:
        \begin{enumerate}
            \item \(u_1,\ldots,u_r\) son linealmente independientes en \(V\)
                  si y solo si \((u_1)_B,\ldots,(u_r)_B\) lo son en \(k^n\);
            \item \(u_1,\ldots,u_r\) generan \(V\) si y solo si
                  \((u_1)_B,\ldots,(u_r)_B\) generan \(k^n\);
            \item \(u_1,\ldots,u_r\) son una base de \(V\) si y solo si
                  \((u_1)_B,\ldots,(u_r)_B\) son una base de \(k^n\).
        \end{enumerate}
        Además, las relaciones de dependencia lineal entre
        \(u_1,\ldots,u_r\) y entre \((u_1)_B,\ldots,(u_r)_B\) son
        \alert{las mismas}, con los mismos coeficientes.
    \end{theorem}

    \medskip

    En resumen: todo lo que se exprese mediante combinaciones lineales
    puede estudiarse indistintamente en \(V\) o en coordenadas.
\end{frame}


\begin{frame}[allowframebreaks,t]
    \frametitle{Traducción a coordenadas: demostración}

    \begin{proof}
        \alert{Relaciones de dependencia y (1).}
        Sean \(\lambda_1,\ldots,\lambda_r\in k\). Aplicando el corolario
        anterior con \(u=0\), y teniendo en cuenta que \(0_B=0\),
        \[
            \lambda_1u_1+\cdots+\lambda_ru_r=0
            \iff
            \lambda_1(u_1)_B+\cdots+\lambda_r(u_r)_B=0.
        \]
        Por tanto, los escalares que anulan una combinación lineal de
        \(u_1,\ldots,u_r\) son exactamente los que anulan la misma
        combinación de \((u_1)_B,\ldots,(u_r)_B\).

        En particular, la única solución es
        \(\lambda_1=\cdots=\lambda_r=0\) en un caso si y solo si lo es en
        el otro, lo que prueba (1).

        \framebreak

        \alert{(2), ``\(\Rightarrow\)''.}
        Supongamos que \(u_1,\ldots,u_r\) generan \(V\) y sea
        \(x\in k^n\). Como \(c_B\) es sobreyectiva, \(x=v_B\) para algún
        \(v\in V\). Escribimos \(v=\lambda_1u_1+\cdots+\lambda_ru_r\), y
        tomando coordenadas,
        \[
            x=v_B=\lambda_1(u_1)_B+\cdots+\lambda_r(u_r)_B.
        \]

        \alert{(2), ``\(\Leftarrow\)''.}
        Supongamos que \((u_1)_B,\ldots,(u_r)_B\) generan \(k^n\) y sea
        \(v\in V\). Entonces existen escalares tales que
        \[
            v_B=\lambda_1(u_1)_B+\cdots+\lambda_r(u_r)_B,
        \]
        y como \(c_B\) refleja combinaciones lineales,
        \(v=\lambda_1u_1+\cdots+\lambda_ru_r\).

        \alert{(3)} Una base es un sistema generador linealmente
        independiente, así que (3) se sigue de (1) y (2).
    \end{proof}
\end{frame}


\begin{frame}[allowframebreaks,t]
    \frametitle{Consecuencia práctica: el rango en coordenadas}

    Sea \(B\) una base de \(V\) con \(n=\dim V\) vectores, sean
    \(u_1,\ldots,u_r\in V\) y sea
    \[
        A=\bigl((u_1)_B\mid\cdots\mid(u_r)_B\bigr)\in M_{n\times r}(k)
    \]
    la matriz $n\times r$ que tiene por columnas sus coordenadas.

    \medskip

    Combinando la traducción a coordenadas con lo que sabemos en \(k^n\):
    \begin{itemize}
        \item \(u_1,\ldots,u_r\) son linealmente independientes
              \(\iff\) \(\rg(A)=r\);
        \item \(u_1,\ldots,u_r\) generan \(V\)
              \(\iff\) \(\rg(A)=\dim V\);
        \item \(u_1,\ldots,u_r\) son una base de \(V\)
              \(\iff\) \(\rg(A)=r=\dim V\).
    \end{itemize}

    \medskip

    Además, cada solución no trivial de \(Ax=0\) da una relación de
    dependencia lineal entre \(u_1,\ldots,u_r\) \alert{en \(V\)}.

    \framebreak

    Por tanto:
    \begin{itemize}
        \item si \(u_1,\ldots,u_r\) son linealmente independientes,
              entonces \(r\leq \dim V\);
        \item si \(u_1,\ldots,u_r\) generan \(V\),
              entonces \(\dim V\leq r\);
        \item si \(r=\dim V\), entonces las siguientes condiciones son equivalentes:
        \begin{itemize}
            \item \(u_1,\ldots,u_r\) son linealmente independientes;
            \item \(u_1,\ldots,u_r\) generan \(V\);
            \item \(u_1,\ldots,u_r\) son una base de \(V\).
        \end{itemize}
    \end{itemize}

    \medskip

    Es lo mismo que vimos en \(k^n\) con los criterios para ser una base,
    ahora en cualquier espacio con una base de \(n\) vectores.
\end{frame}


\begin{frame}
    \frametitle{Ejemplo: polinomios en coordenadas}

    \begin{example}
        En \(\mathbb{R}_2[x]\), con \(B=\{1,x,x^2\}\), consideremos
        \[
            p_1=1+2x,\qquad
            p_2=x-x^2,\qquad
            p_3=1+x+x^2.
        \]
        Sus coordenadas son \((1,2,0) \), \((0,1,-1) \) y \((1,1,1) \).
        Escalonamos la matriz que las tiene por columnas:
        \[
            \begin{pmatrix}
                1&0&1\\
                2&1&1\\
                0&-1&1
            \end{pmatrix}
            \xrightarrow{\;F_2-2F_1\;}
            \begin{pmatrix}
                1&0&1\\
                0&1&-1\\
                0&-1&1
            \end{pmatrix}
            \xrightarrow{\;F_3+F_2\;}
            \begin{pmatrix}
                1&0&1\\
                0&1&-1\\
                0&0&0
            \end{pmatrix}.
        \]
        El rango es \(2<3\): los polinomios son linealmente
        \alert{dependientes} y \alert{no generan} \(\mathbb{R}_2[x]\).
    \end{example}
\end{frame}


\begin{frame}
    \frametitle{Ejemplo: polinomios en coordenadas}

    \begin{example}
        Para encontrar una relación de dependencia entre
        \(p_1=1+2x\), \(p_2=x-x^2\) y \(p_3=1+x+x^2\), resolvemos el
        sistema homogéneo escalonado del frame anterior. La incógnita libre es \(\lambda_3\);
        tomando \(\lambda_3=1\), la segunda fila da \(\lambda_2=1\) y la
        primera, \(\lambda_1=-1\). Así,
        \[
            -(p_1)_B+(p_2)_B+(p_3)_B=0
            \quad\text{en }\mathbb{R}^3.
        \]
        Como las coordenadas reflejan combinaciones lineales, la misma
        relación se cumple en \(\mathbb{R}_2[x]\):
        \[
            -p_1+p_2+p_3=-(1+2x)+(x-x^2)+(1+x+x^2)=0.
        \]
    \end{example}
\end{frame}


\begin{frame}
    \frametitle{Ejemplo: matrices en coordenadas}

    \begin{example}
        En \(M_{2\times2}(\mathbb{R})\), con la base
        \(B=\{E_{11},E_{12},E_{21},E_{22}\}\), consideremos
        \[
            A_1=\begin{pmatrix}1&1\\0&0\end{pmatrix},\quad
            A_2=\begin{pmatrix}0&1\\1&0\end{pmatrix},\quad
            A_3=\begin{pmatrix}0&0\\1&1\end{pmatrix},\quad
            A_4=\begin{pmatrix}1&0\\0&0\end{pmatrix}.
        \]
        Las coordenadas de \(\begin{pmatrix}a&b\\c&d\end{pmatrix}\) respecto de
        \(B\) son \((a,b,c,d) \). Por tanto, la matriz que tiene por
        columnas las coordenadas de \(A_1,\ldots,A_4\) es
        \[
            \begin{pmatrix}
                1&0&0&1\\
                1&1&0&0\\
                0&1&1&0\\
                0&0&1&0
            \end{pmatrix}.
        \]
    \end{example}
\end{frame}


\begin{frame}
    \frametitle{Ejemplo: matrices en coordenadas}

    \begin{example}
        Escalonamos:
        \vspace{-1ex}
        \[
            \begin{pmatrix}
                1&0&0&1\\
                1&1&0&0\\
                0&1&1&0\\
                0&0&1&0
            \end{pmatrix}
            \xrightarrow{\;F_2-F_1\;}
            \begin{pmatrix}
                1&0&0&1\\
                0&1&0&-1\\
                0&1&1&0\\
                0&0&1&0
            \end{pmatrix}
            \xrightarrow{\;F_3-F_2\;}
            \begin{pmatrix}
                1&0&0&1\\
                0&1&0&-1\\
                0&0&1&1\\
                0&0&1&0
            \end{pmatrix}
        \]
        \[
            \xrightarrow{\;F_4-F_3\;}
            \begin{pmatrix}
                1&0&0&1\\
                0&1&0&-1\\
                0&0&1&1\\
                0&0&0&-1
            \end{pmatrix}.
        \]
        El rango es \(4\), que coincide con el número de matrices y con el
        número de vectores de \(B\). Por tanto, \(A_1,A_2,A_3,A_4\) son
        una \alert{base} de \(M_{2\times2}(\mathbb{R})\).
        No ha hecho falta plantear ninguna ecuación entre matrices:
        todo el trabajo se ha hecho en \(\mathbb{R}^4\).
    \end{example}
\end{frame}

\begin{frame}
    \frametitle{Existencia de bases en coordenadas: el método}

    Sea \(B\) una base de \(V\) con \(n=\dim V\) vectores y sean
    \(v_1,\ldots,v_r,v_{r+1},\ldots,v_{r+s}\in V\) vectores que generan
    \(V\), con \(v_1,\ldots,v_r\) linealmente independientes.

    \medskip

    \alert{Método.}
    \begin{enumerate}
        \item Calculamos las coordenadas \((v_1)_B,\ldots,(v_{r+s})_B\in k^n\).
        \item Formamos la matriz
              \(A=\bigl((v_1)_B\mid\cdots\mid(v_{r+s})_B\bigr)\) y la
              escalonamos por filas.
        \item Nos quedamos con los \(v_j\) tales que la columna \(j\) de
              la forma escalonada contiene un pivote.
    \end{enumerate}
    Estos vectores forman una base de \(V\) que contiene a
    \(v_1,\ldots,v_r\).

    \medskip

    \alert{Por qué funciona.} La traducción a coordenadas conserva la
    generación, la independencia lineal y las bases. Basta, pues, aplicar
    el método en \(k^n\) a \((v_1)_B,\ldots,(v_{r+s})_B\).
\end{frame}


\begin{frame}
    \frametitle{Ejemplo: una base de \(\mathbb{R}_2[x]\) entre unos generadores}

    \begin{example}
        En \(\mathbb{R}_2[x]\), buscamos una base que contenga a \(1+x\)
        formada por algunos de los polinomios

        \vspace{-10pt}

        \[
            v_1=1+x,\quad v_2=2+2x,\quad v_3=x^2,\quad v_4=1+x+x^2,\quad v_5=x,
        \]

        \vspace{-0pt}

        y que contenga al primero, que al ser no nulo forma en sí mismo un conjunto linealmente independiente.
        Sus coordenadas en \(\{1,x,x^2\}\) son las columnas de

        \medskip

        \[
            \begin{pmatrix}
                1&2&0&1&0\\
                1&2&0&1&1\\
                0&0&1&1&0
            \end{pmatrix}
            \xrightarrow{\;F_2-F_1\;}
            \begin{pmatrix}
                1&2&0&1&0\\
                0&0&0&0&1\\
                0&0&1&1&0
            \end{pmatrix}
            \xrightarrow{\;F_2\leftrightarrow F_3\;}
            \begin{pmatrix}
                \smash{\overset{\raisebox{1.5ex}{$\scriptstyle\alert{1}$}}{\alert{1}}}&2&\smash{\overset{\raisebox{1.5ex}{$\scriptstyle\alert{3}$}}{0}}&1&\smash{\overset{\raisebox{1.5ex}{$\scriptstyle\alert{5}$}}{0}}\\
                0&0&\alert{1}&1&0\\
                0&0&0&0&\alert{1}
            \end{pmatrix}.
        \]
        El rango es \(3\), así que los cinco polinomios generan
        \(\mathbb{R}_2[x]\). Los pivotes están en las columnas \(1\),
        \(3\) y \(5\), luego el siguiente conjunto es una base que contiene a \(1+x\):
        \[
            \{v_{\alert{1}}, v_{\alert{3}}, v_{\alert{5}}\} = \{1+x, x^2, x\}.
        \]
    \end{example}
\end{frame}


\begin{frame}
    \frametitle{Ejemplo: una base de \(M_{2\times2}(\mathbb{R})\) entre unos generadores}

    \begin{example}
        En \(M_{2\times2}(\mathbb{R})\) consideramos

        \vspace{-3pt}

        \[
            v_1=\begin{pmatrix}1&0\\0&1\end{pmatrix},\quad
            v_2=\begin{pmatrix}0&1\\1&0\end{pmatrix},\quad
            v_3=\begin{pmatrix}1&1\\1&1\end{pmatrix},
        \]
        \[
            v_4=\begin{pmatrix}1&0\\0&-1\end{pmatrix},\quad
            v_5=\begin{pmatrix}0&1\\-1&0\end{pmatrix},\quad
            v_6=\begin{pmatrix}2&1\\3&0\end{pmatrix}.
        \]
        Los vectores \(v_1,v_2\) son linealmente independientes, pues
        ninguno es múltiplo del otro. Buscamos una base de
        \(M_{2\times2}(\mathbb{R})\) que los contenga, formada por algunos
        de los \(v_j\).

        Tomamos coordenadas respecto de la base
        \(B=\{E_{11},E_{12},E_{21},E_{22}\}\):
        \[
            \begin{pmatrix}a&b\\c&d\end{pmatrix}_B=\begin{pmatrix}a\\b\\c\\d\end{pmatrix}.
        \]
    \end{example}
\end{frame}


\begin{frame}
    \frametitle{Ejemplo: una base de \(M_{2\times2}(\mathbb{R})\) entre unos generadores}

    \begin{example}
        Escalonamos la matriz que tiene por columnas las coordenadas de
        \(v_1,\ldots,v_6\):
        \begin{align*}
            \begin{pmatrix}
                1&0&1&1&0&2\\
                0&1&1&0&1&1\\
                0&1&1&0&-1&3\\
                1&0&1&-1&0&0
            \end{pmatrix}
            &\xrightarrow[\;F_4-F_1\;]{\;F_3-F_2\;}
            \begin{pmatrix}
                1&0&1&1&0&2\\
                0&1&1&0&1&1\\
                0&0&0&0&-2&2\\
                0&0&0&-2&0&-2
            \end{pmatrix}\\[4ex]
            &\xrightarrow{\;F_3\leftrightarrow F_4\;}
            \begin{pmatrix}
                \smash{\overset{\raisebox{1.5ex}{$\scriptstyle\alert{1}$}}{\alert{1}}}&\smash{\overset{\raisebox{1.5ex}{$\scriptstyle\alert{2}$}}{0}}&1&\smash{\overset{\raisebox{1.5ex}{$\scriptstyle\alert{4}$}}{1}}&\smash{\overset{\raisebox{1.5ex}{$\scriptstyle\alert{5}$}}{0}}&2\\
                0&\alert{1}&1&0&1&1\\
                0&0&0&\alert{-2}&0&-2\\
                0&0&0&0&\alert{-2}&2
            \end{pmatrix}.
        \end{align*}
        El rango es \(4\), luego los \(v_j\) generan. Los pivotes están
        en las columnas \(1,2,4,5\): \(\{v_{\alert{1}}, v_{\alert{2}}, v_{\alert{4}}, v_{\alert{5}}\}\) es una base
        que contiene a \(v_1,v_2\).
    \end{example}
\end{frame}








\begin{frame}
    \frametitle{Matriz de cambio de base}

    Sean $B=\{u_1,\ldots,u_n\}$ y $B'=\{u'_1,\ldots,u'_m\}$
    dos bases de $V$.

    \begin{definition}
        La \alert{matriz de cambio de base de $B$ a $B'$} es la
        matriz $m\times n$
        \[
            M(B,B')=\bigl((u_1)_{B'}\mid\cdots\mid(u_n)_{B'}\bigr),
        \]
        cuya $j$-ésima columna es el vector de coordenadas de $u_j$ en $B'$.
    \end{definition}
\end{frame}

\begin{frame}
    \frametitle{Propiedades de la matriz de cambio de base}

    \begin{proposition}
        Para todo $v\in V$,
        \[
            v_{B'}=M(B,B')\cdot v_B.
        \]
    \end{proposition}

    \begin{proof}
        Si $v_B=(\lambda_1,\ldots,\lambda_n) $, entonces $v=\sum_j\lambda_j u_j$.
        Por las propiedades de $c_{B'}$,
        \[
            v_{B'}
            =\sum_j\lambda_j(u_j)_{B'}
            =M(B,B')\,v_B.\qedhere
        \]
    \end{proof}
\end{frame}


\begin{frame}
    \frametitle{Ejemplo de matriz de cambio de base}

    En $\mathbb{R}^2$, sean $C=\{e_1,e_2\}$ y $B=\{(1,1),(1,-1)\}$.

    Las columnas de $M(C,B)$ son las coordenadas de $e_1$ y $e_2$ en $B$.
    Resolviendo:
    \[
        e_1=\tfrac{1}{2}(1,1)+\tfrac{1}{2}(1,-1),
        \qquad
        e_2=\tfrac{1}{2}(1,1)-\tfrac{1}{2}(1,-1),
    \]
    de modo que
    \[
        M(C,B)=
        \begin{pmatrix}
            \tfrac{1}{2}&\tfrac{1}{2}\\[4pt]
            \tfrac{1}{2}&-\tfrac{1}{2}
        \end{pmatrix}.
    \]
    Para $v=(3,1)$, $v_C=(3,1) $ y
    \[
        v_B=M(C,B)\,v_C=
        \begin{pmatrix}
            \tfrac{1}{2}&\tfrac{1}{2}\\[4pt]
            \tfrac{1}{2}&-\tfrac{1}{2}
        \end{pmatrix}
        \begin{pmatrix}3\\1\end{pmatrix}
        =\begin{pmatrix}2\\1\end{pmatrix}.
    \]
    Comprobación: $2(1,1)+1\cdot(1,-1)=(3,1)$. \checkmark
\end{frame}


\begin{frame}[allowframebreaks,t]
    \frametitle{Propiedades de la matriz de cambio de base}

    \begin{proposition}
        Sean $B$, $B'$, $B''$ bases de $V$ con $n$, $m$ y $p$
        vectores respectivamente.
        \begin{enumerate}
            \item $M(B,B)=I_n$.
            \item $M(B,B'')=M(B',B'')\cdot M(B,B')$.
        \end{enumerate}
    \end{proposition}

    \begin{proof}
        Para (1): la $j$-ésima columna de $M(B,B)$ es $(u_j)_B=e_j$,
        pues $u_j=0\cdot u_1+\cdots+1\cdot u_j+\cdots+0\cdot u_n$.
        Por tanto $M(B,B)=I_n$.

        Para (2): sea $B=\{u_1,\ldots,u_n\}$. Por definición,
        $M(B,B')$ tiene por columnas las coordenadas de los $u_j$ en $B'$:
        \[
            M(B,B')=\bigl((u_1)_{B'}\mid\cdots\mid(u_n)_{B'}\bigr).
        \]
        Multiplicando por bloques, con $M(B',B'')$ como un único bloque y
        $M(B,B')$ dividida en columnas,
        \[
            M(B',B'')\cdot M(B,B')
            =\bigl(M(B',B'')\,(u_1)_{B'}\mid\cdots\mid M(B',B'')\,(u_n)_{B'}\bigr).
        \]

        Por la proposición anterior, aplicada a la matriz de cambio de
        base de $B'$ a $B''$ y al vector $u_j$,
        \[
            M(B',B'')\,(u_j)_{B'}=(u_j)_{B''},
            \qquad j=1,\ldots,n.
        \]
        Sustituyendo en cada columna,
        \[
            M(B',B'')\cdot M(B,B')
            =\bigl((u_1)_{B''}\mid\cdots\mid(u_n)_{B''}\bigr),
        \]
        que, por definición, es la matriz $M(B,B'')$.
    \end{proof}
\end{frame}


\begin{frame}[allowframebreaks]
    \frametitle{Invertibilidad y dimensión}

    \begin{corollary}
        Sean $B$ y $B'$ bases de $V$ con $n$ y $m$ vectores respectivamente.
        La matriz $M(B,B')$ es invertible, con $M(B,B')^{-1}=M(B',B)$,
        y en particular $n=m$.
    \end{corollary}

    \begin{proof}
        Por las propiedades (1) y (2),
        \[
            M(B',B)\cdot M(B,B')=M(B,B)=I_n,
        \]
        \[
            M(B,B')\cdot M(B',B)=M(B',B')=I_m.
        \]
        Para que estos dos productos expresen que $M(B',B)$ es la inversa
        de $M(B,B')$ —es decir, para que ambos den la misma identidad—
        es necesario que $I_n=I_m$, y por tanto $n=m$.
        Siendo $M(B,B')$ una matriz $n\times n$, su inversa es $M(B',B)$.
    \end{proof}

    \begin{definition}
        Si $V$ admite una base finita, el número de vectores de cualquier
        base se llama \alert{dimensión} de $V$ sobre $k$, y se denota
        $\dim_k V$ o $\dim V$.
        Si $V$ no admite ninguna base finita, decimos que $V$ tiene
        \alert{dimensión infinita}.
    \end{definition}

    A partir de las bases ya conocidas: $\dim k^n=n$,
    $\dim k_n[x]=n+1$, $\dim M_{m\times n}(k)=mn$, $\dim k^X=|X|$.
    El espacio $k[x]$ tiene dimensión infinita: toda combinación lineal
    finita de polinomios tiene grado acotado y no puede generar $k[x]$.

    Según se ha visto anteriormente, cualquier espacio vectorial con un conjunto generador finito tiene dimensión finita.
\end{frame}


\begin{frame}[allowframebreaks,t]
    \frametitle{Ampliar una familia independiente hasta una base}

    \begin{corollary}
        Sea \(V\) un espacio vectorial de dimensión \(n\) y sean
        \(v_1,\ldots,v_r\in V\) linealmente independientes. Entonces
        \(r\leq n\), y existen \(v_{r+1},\ldots,v_n\in V\) tales que
        \(v_1,\ldots,v_n\) forman una base de \(V\). Además,
        \(v_{r+1},\ldots,v_n\) pueden elegirse entre los vectores de
        cualquier base dada de \(V\).
    \end{corollary}

    \begin{proof}
        Sea \(\{w_1,\ldots,w_n\}\) una base de \(V\). El conjunto
        \[
            S=\{v_1,\ldots,v_r,w_1,\ldots,w_n\}
        \]
        genera \(V\), porque ya lo generan los \(w_j\). Por el teorema de
        existencia de bases, \(S\) contiene una base \(B\) de \(V\) que
        contiene a \(v_1,\ldots,v_r\); los demás vectores de \(B\) son
        algunos de los \(w_j\). Como toda base de \(V\) tiene \(n\)
        vectores, se han añadido exactamente \(n-r\) de ellos; en
        particular, \(r\leq n\).
    \end{proof}

    \medskip

    En la práctica, basta encontrar \(n-r\) vectores que, junto con
    \(v_1,\ldots,v_r\), sean \alert{linealmente independientes}: al ser
    \(n\) vectores en un espacio de dimensión \(n\), ya forman una base.
\end{frame}


\begin{frame}
    \frametitle{Ejemplo: ampliar a una base de \(\mathbb{R}^4\)}

    \begin{example}
        Los vectores \(v_1=(1,2,0,1)\) y \(v_2=(0,1,1,1)\) de
        \(\mathbb{R}^4\) son linealmente independientes, pues ninguno es
        múltiplo del otro. Para ampliarlos a una base buscamos dos vectores
        más; probamos con \(e_3\) y \(e_4\). La matriz que tiene por
        columnas \(v_1,v_2,e_3,e_4\) es
        \[
            \begin{pmatrix}
                1&0&0&0\\
                2&1&0&0\\
                0&1&1&0\\
                1&1&0&1
            \end{pmatrix},
        \]
        triangular con unos en la diagonal. Su determinante es \(1\neq0\),
        luego \(v_1,v_2,e_3,e_4\) forman una base de \(\mathbb{R}^4\).
    \end{example}

    \medskip

    Hemos elegido vectores canónicos que hacen \alert{triangular} la
    matriz, para que el determinante se calcule sin esfuerzo.
\end{frame}

\begin{frame}
    \frametitle{Ejemplo: ampliar a una base de \(\mathbb{R}^4\) sistemáticamente}

    \begin{example}

        \vspace{-5pt}

        Ampliamos de nuevo \(v_1=(1,2,0,1)\), \(v_2=(0,1,1,1)\), ahora sin
        tantear. Como en la demostración del corolario, añadimos la base
        canónica: \(v_1,v_2,e_1,e_2,e_3,e_4\) generan \(\mathbb{R}^4\), y
        aplicamos el método de existencia de bases en \(k^n\):

        \vspace{-15pt}

        \begin{align*}
            \begin{pmatrix}
                1&0&1&0&0&0\\
                2&1&0&1&0&0\\
                0&1&0&0&1&0\\
                1&1&0&0&0&1
            \end{pmatrix}
            &\xrightarrow[\;F_4-F_1\;]{\;F_2-2F_1\;}
            \begin{pmatrix}
                1&0&1&0&0&0\\
                0&1&-2&1&0&0\\
                0&1&0&0&1&0\\
                0&1&-1&0&0&1
            \end{pmatrix}
            \xrightarrow[\;F_4-F_2\;]{\;F_3-F_2\;}
            \begin{pmatrix}
                1&0&1&0&0&0\\
                0&1&-2&1&0&0\\
                0&0&2&-1&1&0\\
                0&0&1&-1&0&1
            \end{pmatrix}\\[2.5ex]
            &\xrightarrow[\;F_4-2F_3\;]{\;F_3\leftrightarrow F_4\;}
            \begin{pmatrix}
                \smash{\overset{\raisebox{1ex}{$\scriptstyle\alert{1}$}}{\alert{1}}}&\smash{\overset{\raisebox{1ex}{$\scriptstyle\alert{2}$}}{0}}&\smash{\overset{\raisebox{1ex}{$\scriptstyle\alert{3}$}}{1}}&\smash{\overset{\raisebox{1ex}{$\scriptstyle\alert{4}$}}{0}}&0&0\\
                0&\alert{1}&-2&1&0&0\\
                0&0&\alert{1}&-1&0&1\\
                0&0&0&\alert{1}&1&-2
            \end{pmatrix}.
        \end{align*}

        \vspace{-8pt}

        Los pivotes están en las columnas \(1,2,3,4\): \(\{v_1,v_2,e_1,e_2\}\)
        es una base de \(\mathbb{R}^4\) que contiene a \(v_1,v_2\).
    \end{example}
\end{frame}



\begin{frame}
    \frametitle{Ejemplo: ampliar a una base de \(\mathbb{R}_2[x]\)}

    \begin{example}
        Los polinomios \(1+x\) y \(x-x^2\) son linealmente independientes,
        pues ninguno es múltiplo del otro. Para ampliarlos a una base de
        \(\mathbb{R}_2[x]\), que tiene dimensión \(3\), basta añadir un
        polinomio adecuado.

        Trabajamos en coordenadas respecto de \(\{1,x,x^2\}\). Probamos
        con \(x^2\): la matriz de coordenadas de \(1+x,\ x-x^2,\ x^2\) es
        \[
            \begin{pmatrix}
                1&0&0\\
                1&1&0\\
                0&-1&1
            \end{pmatrix},
        \]
        triangular con determinante \(1\neq0\). Por tanto
        \(\{1+x,\ x-x^2,\ x^2\}\) es una base de \(\mathbb{R}_2[x]\).
    \end{example}
\end{frame}

\begin{frame}
    \frametitle{Ejemplo: ampliar a una base de \(\mathbb{R}_2[x]\) sistemáticamente}

    \begin{example}
        Ampliamos de nuevo \(1+x\), \(x-x^2\), ahora sin tantear. Añadimos
        la base \(\{1,x,x^2\}\): los polinomios
        \[
            v_1=1+x,\quad v_2=x-x^2,\quad v_3=1,\quad v_4=x,\quad v_5=x^2
        \]
        generan \(\mathbb{R}_2[x]\), y aplicamos el método de existencia de
        bases en coordenadas respecto de \(\{1,x,x^2\}\):
        \[
            \begin{pmatrix}
                1&0&1&0&0\\
                1&1&0&1&0\\
                0&-1&0&0&1
            \end{pmatrix}
            \xrightarrow{\;F_2-F_1\;}
            \begin{pmatrix}
                1&0&1&0&0\\
                0&1&-1&1&0\\
                0&-1&0&0&1
            \end{pmatrix}
            \xrightarrow{\;F_3+F_2\;}
            \begin{pmatrix}
                \smash{\overset{\raisebox{1.5ex}{$\scriptstyle\alert{1}$}}{\alert{1}}}&\smash{\overset{\raisebox{1.5ex}{$\scriptstyle\alert{2}$}}{0}}&\smash{\overset{\raisebox{1.5ex}{$\scriptstyle\alert{3}$}}{1}}&0&0\\
                0&\alert{1}&-1&1&0\\
                0&0&\alert{-1}&1&1
            \end{pmatrix}.
        \]
        Los pivotes están en las columnas \(1,2,3\):
        \(\{v_{\alert{1}}, v_{\alert{2}}, v_{\alert{3}}\}=\{1+x,\ x-x^2,\ 1\}\) es una base de \(\mathbb{R}_2[x]\) que
        contiene a \(1+x\) y \(x-x^2\).
    \end{example}
\end{frame}

\subsection{Producto de espacios vectoriales}

\begin{frame}
    \frametitle{Un ejemplo antes de la definición}

    Supongamos que queremos manejar a la vez un vector de $\mathbb{R}^2$
    y un polinomio de $\mathbb{R}_1[x]$, es decir, pares
    \[
        \bigl((a,b),\,c+dx\bigr),
        \qquad
        (a,b)\in\mathbb{R}^2,
        \quad
        c+dx\in\mathbb{R}_1[x].
    \]

    Es natural sumar y multiplicar por escalares
    \alert{componente a componente}:
    \[
        \bigl((1,2),\,3-x\bigr)+\bigl((0,-1),\,1+4x\bigr)
        =
        \bigl((1,1),\,4+3x\bigr),
    \]
    \[
        2\,\bigl((1,2),\,3-x\bigr)
        =
        \bigl((2,4),\,6-2x\bigr).
    \]

    \medskip

    Cada componente se opera en su propio espacio vectorial,
    con independencia de la otra.
\end{frame}


\begin{frame}
    \frametitle{Producto de dos espacios vectoriales}

    \begin{definition}
        Sean $V$ y $W$ espacios vectoriales sobre el mismo cuerpo $k$.
        El \alert{producto} de $V$ y $W$ es el conjunto
        \[
            V\times W=\{(v,w):v\in V,\ w\in W\}
        \]
        con las operaciones
        \[
            (v,w)+(v',w')=(v+v',\,w+w'),
            \qquad
            \lambda(v,w)=(\lambda v,\,\lambda w),
        \]
        para $v,v'\in V$, $w,w'\in W$ y $\lambda\in k$.
    \end{definition}

    \medskip

    En la primera componente se usan las operaciones de $V$,
    y en la segunda, las de $W$.
\end{frame}


\begin{frame}[allowframebreaks]
    \frametitle{El producto es un espacio vectorial}

    \begin{proposition}
        Con las operaciones anteriores, $V\times W$ es un espacio vectorial
        sobre $k$. Su vector cero es $(0,0)$ y el opuesto de $(v,w)$ es
        $(-v,-w)$.
    \end{proposition}

    \begin{proof}
        Hay que comprobar los axiomas. La idea es la misma en todos:
        cada axioma en $V\times W$ se reduce al mismo axioma en $V$
        (primera componente) y en $W$ (segunda componente).

        Por ejemplo, la distributividad respecto de la suma de vectores:
        \begin{align*}
            \lambda\bigl((v,w)+(v',w')\bigr)
            &=\lambda(v+v',\,w+w')
            =\bigl(\lambda(v+v'),\,\lambda(w+w')\bigr)\\
            &=(\lambda v+\lambda v',\,\lambda w+\lambda w')
            =(\lambda v,\lambda w)+(\lambda v',\lambda w')\\
            &=\lambda(v,w)+\lambda(v',w'),
        \end{align*}
        donde en la segunda línea se han usado los axiomas en $V$ y en $W$.

        Para el vector cero,
        \[
            (v,w)+(0,0)=(v+0,\,w+0)=(v,w),
        \]
        y para el opuesto,
        \[
            (v,w)+(-v,-w)=\bigl(v+(-v),\,w+(-w)\bigr)=(0,0).
        \]

        Los axiomas restantes se comprueban del mismo modo, componente
        a componente.
    \end{proof}
\end{frame}


\begin{frame}
    \frametitle{Combinaciones lineales en el producto}

    Como las operaciones se hacen componente a componente, lo mismo
    ocurre con las combinaciones lineales.

    \begin{proposition}
        Para $(v_1,w_1),\ldots,(v_r,w_r)\in V\times W$ y
        $\lambda_1,\ldots,\lambda_r\in k$,
        \[
            \lambda_1(v_1,w_1)+\cdots+\lambda_r(v_r,w_r)
            =
            \bigl(\lambda_1v_1+\cdots+\lambda_rv_r,\;
                  \lambda_1w_1+\cdots+\lambda_rw_r\bigr).
        \]
    \end{proposition}

    \begin{proof}
        Por la definición del producto por escalares,
        $\lambda_i(v_i,w_i)=(\lambda_iv_i,\lambda_iw_i)$ para cada $i$.
        Sumando estos $r$ pares y aplicando repetidamente la definición
        de la suma en $V\times W$, se obtiene el par cuyas componentes son
        las sumas de las componentes.
    \end{proof}

    \medskip

    En particular, una combinación lineal en $V\times W$ es nula
    si y solo si \alert{ambas} componentes lo son.
\end{frame}


\begin{frame}
    \frametitle{Ejemplo: $k^m\times k^n$ y $k^{m+n}$}

    \begin{example}
        Eliminando los paréntesis interiores, a cada par de
        $k^m\times k^n$ le corresponde un vector de $k^{m+n}$:
        \[
            \bigl((a_1,\ldots,a_m),\,(b_1,\ldots,b_n)\bigr)
            \;\longleftrightarrow\;
            (a_1,\ldots,a_m,b_1,\ldots,b_n).
        \]
        Esta correspondencia es biyectiva y respeta las operaciones:
        sumar o multiplicar por escalares antes o después de eliminar
        los paréntesis da el mismo resultado, porque en ambos casos
        se opera coordenada a coordenada.
    \end{example}

    \medskip

    Por eso, en la práctica, identificamos
    \[
        k^m\times k^n=k^{m+n}.
    \]
    En el Tema~4 daremos un sentido preciso a esta identificación.
\end{frame}


\begin{frame}
    \frametitle{Producto de varios espacios vectoriales}

    \begin{definition}
        Sean $V_1,\ldots,V_r$ espacios vectoriales sobre $k$.
        Su \alert{producto} es
        \[
            V_1\times\cdots\times V_r
            =
            \{(u_1,\ldots,u_r):u_i\in V_i\},
        \]
        con la suma y el producto por escalares componente a componente:
        \[
            (u_1,\ldots,u_r)+(u'_1,\ldots,u'_r)=(u_1+u'_1,\ldots,u_r+u'_r),
        \]
        \[
            \lambda(u_1,\ldots,u_r)=(\lambda u_1,\ldots,\lambda u_r).
        \]
    \end{definition}

    Como en el caso de dos factores, es un espacio vectorial.
    El caso más conocido es $k^n=k\times\cdots\times k$, con $n$ factores.
\end{frame}


\begin{frame}
    \frametitle{Una base del producto: un ejemplo}

    En $\mathbb{R}^2\times\mathbb{R}_1[x]$, todo par se descompone como
    \[
        \bigl((a,b),\,c+dx\bigr)
        =
        \bigl((a,b),\,0\bigr)+\bigl((0,0),\,c+dx\bigr),
    \]
    y por tanto
    \[
        \bigl((a,b),\,c+dx\bigr)
        =
        a\bigl((1,0),0\bigr)+b\bigl((0,1),0\bigr)
        +c\bigl((0,0),1\bigr)+d\bigl((0,0),x\bigr).
    \]

    \medskip

    Los cuatro vectores que aparecen se obtienen
    \begin{itemize}
        \item de la base canónica de $\mathbb{R}^2$, poniendo $0$
              en la segunda componente;
        \item de la base $\{1,x\}$ de $\mathbb{R}_1[x]$, poniendo $(0,0)$
              en la primera componente.
    \end{itemize}

    \medskip

    \begin{center}
        \alert{¿Forman siempre una base del producto los vectores así obtenidos?}
    \end{center}
\end{frame}


\begin{frame}
    \frametitle{Bases del producto}

    \begin{theorem}
        Sean $\{v_1,\ldots,v_m\}$ una base de $V$ y
        $\{w_1,\ldots,w_n\}$ una base de $W$. Entonces
        \[
            B=\{(v_1,0),\ldots,(v_m,0),\,(0,w_1),\ldots,(0,w_n)\}
        \]
        es una base de $V\times W$. En particular,
        \[
            \dim(V\times W)=\dim V+\dim W.
        \]
    \end{theorem}

    \medskip

    Hay que probar dos cosas:
    \begin{enumerate}
        \item $B$ genera $V\times W$;
        \item $B$ es linealmente independiente.
    \end{enumerate}
    En ambas usaremos que las combinaciones lineales en $V\times W$
    se calculan componente a componente.
\end{frame}


\begin{frame}[allowframebreaks]
    \frametitle{Bases del producto: demostración}

    \begin{proof}
        \alert{$B$ genera $V\times W$.}
        Sea $(v,w)\in V\times W$. Como $\{v_1,\ldots,v_m\}$ genera $V$ y
        $\{w_1,\ldots,w_n\}$ genera $W$, existen escalares tales que
        \[
            v=\lambda_1v_1+\cdots+\lambda_mv_m,
            \qquad
            w=\mu_1w_1+\cdots+\mu_nw_n.
        \]
        Calculando componente a componente,
        \begin{align*}
            &\lambda_1(v_1,0)+\cdots+\lambda_m(v_m,0)
            +\mu_1(0,w_1)+\cdots+\mu_n(0,w_n)\\
            &\qquad=
            \bigl(\lambda_1v_1+\cdots+\lambda_mv_m,\;
                  \mu_1w_1+\cdots+\mu_nw_n\bigr)
            =(v,w),
        \end{align*}
        ya que en la primera componente los sumandos $\mu_j\cdot0$ son
        nulos, y en la segunda lo son los $\lambda_i\cdot0$.

        \framebreak

        \alert{$B$ es linealmente independiente.}
        Supongamos que
        \[
            \lambda_1(v_1,0)+\cdots+\lambda_m(v_m,0)
            +\mu_1(0,w_1)+\cdots+\mu_n(0,w_n)=(0,0).
        \]
        El mismo cálculo de antes muestra que el miembro izquierdo es
        \[
            \bigl(\lambda_1v_1+\cdots+\lambda_mv_m,\;
                  \mu_1w_1+\cdots+\mu_nw_n\bigr).
        \]
        Igualando componentes,
        \[
            \lambda_1v_1+\cdots+\lambda_mv_m=0
            \quad\text{en } V,
            \qquad
            \mu_1w_1+\cdots+\mu_nw_n=0
            \quad\text{en } W.
        \]
        Como $v_1,\ldots,v_m$ son linealmente independientes,
        $\lambda_1=\cdots=\lambda_m=0$; y como $w_1,\ldots,w_n$ lo son,
        $\mu_1=\cdots=\mu_n=0$.

        Por tanto $B$ es una base de $V\times W$. Tiene $m+n$ vectores
        (son distintos dos a dos, por serlo los $v_i$ y los $w_j$, y por
        no ser nulos), luego
        \[
            \dim(V\times W)=m+n=\dim V+\dim W.\qedhere
        \]
    \end{proof}
\end{frame}


\begin{frame}
    \frametitle{Dimensión de un producto de varios espacios}

    \begin{corollary}
        Si $V_1,\ldots,V_r$ tienen dimensión finita, entonces
        \[
            \dim(V_1\times\cdots\times V_r)
            =
            \dim V_1+\cdots+\dim V_r.
        \]
    \end{corollary}

    \begin{proof}
        Se elige una base de cada $V_i$ y se coloca cada uno de sus vectores
        en la componente $i$-ésima, con $0$ en las demás. La misma
        demostración que para dos factores, ahora con $r$ componentes,
        prueba que todos estos vectores forman una base del producto.
        Su número es $\dim V_1+\cdots+\dim V_r$.
    \end{proof}

    \medskip

    Por ejemplo, $\dim k^n=\dim k+\cdots+\dim k=n$, y la base obtenida
    es precisamente la base canónica $\{e_1,\ldots,e_n\}$.
\end{frame}


\begin{frame}
    \frametitle{Ejemplos de dimensión de productos}

    \begin{example}
        \begin{itemize}
            \item $\dim\bigl(\mathbb{R}^2\times\mathbb{R}_1[x]\bigr)=2+2=4$.
            \item $\dim\bigl(M_{2\times2}(k)\times k_3[x]\bigr)=4+4=8$.
            \item $\dim\bigl(k^m\times k^n\bigr)=m+n$, de acuerdo con la
                  identificación $k^m\times k^n=k^{m+n}$.
            \item $\dim\bigl(M_{m\times n}(k)\times M_{m\times n}(k)\bigr)=2mn$.
        \end{itemize}
    \end{example}

    En $\mathbb{R}^2\times\mathbb{R}_1[x]$, con la base
    \[
        B=\bigl\{((1,0),0),\ ((0,1),0),\ ((0,0),1),\ ((0,0),x)\bigr\},
    \]
    el vector $\bigl((3,-1),\,2-5x\bigr)$ tiene coordenadas
    \[
        \bigl((3,-1),\,2-5x\bigr)_B
        =
        \begin{pmatrix}3\\-1\\2\\-5\end{pmatrix}.
    \]
\end{frame}


\begin{frame}
    \frametitle{No toda base del producto es de esta forma}

    \begin{example}
        En $\mathbb{R}\times\mathbb{R}=\mathbb{R}^2$, los vectores
        \[
            (1,1),
            \qquad
            (1,-1)
        \]
        forman una base, pues la matriz
        \[
            \begin{pmatrix}
                1&1\\
                1&-1
            \end{pmatrix}
        \]
        tiene rango $2$.

        Sin embargo, ninguno de ellos es de la forma $(v,0)$ ni $(0,w)$.
    \end{example}

    \medskip

    El teorema proporciona \alert{una} base de $V\times W$ a partir
    de bases de $V$ y de $W$, pero $V\times W$ tiene muchas otras bases
    que no se obtienen así.
\end{frame}

% \begin{frame}
    \frametitle{Ejercicio}

    Estudia la dependencia o independencia lineal de los siguientes vectores
    de \(\mathbb{R}^3\):
    \[
        v_1=(1,2,1),\qquad
        v_2=(2,1,0),\qquad
        v_3=(3,3,1).
    \]

    Si son linealmente dependientes, encuentra una relación de dependencia
    lineal.
\end{frame}


\begin{frame}
    \frametitle{Ejercicio}

    Determina para qué valores de \(a\in\mathbb{R}\) son linealmente
    independientes los vectores
    \[
        v_1=(1,0,1),\qquad
        v_2=(0,1,1),\qquad
        v_3=(1,1,a).
    \]

    Para los restantes valores de \(a\), encuentra una relación de
    dependencia lineal.
\end{frame}


\begin{frame}
    \frametitle{Ejercicio}

    Estudia la dependencia o independencia lineal en \(\mathbb{R}_2[x]\)
    de los polinomios
    \[
        p_1(x)=1+x,
        \qquad
        p_2(x)=1-x,
        \qquad
        p_3(x)=1+x^2.
    \]

    \medskip

    Para ello, estudia cuándo
    \[
        \lambda_1p_1(x)+
        \lambda_2p_2(x)+
        \lambda_3p_3(x)=0.
    \]
\end{frame}


\begin{frame}
    \frametitle{Ejercicio}

    En \(M_{2\times2}(\mathbb{R})\), consideremos
    \[
        A_1=
        \begin{pmatrix}
            1&0\\
            0&1
        \end{pmatrix},
        \qquad
        A_2=
        \begin{pmatrix}
            0&1\\
            1&0
        \end{pmatrix},
        \qquad
        A_3=
        \begin{pmatrix}
            1&1\\
            1&1
        \end{pmatrix}.
    \]

    Estudia si \(A_1,A_2,A_3\) son linealmente dependientes.

    Si lo son, encuentra una relación de dependencia lineal.
\end{frame}

% \begin{frame}
    \frametitle{Ejercicio}

    Determina si los vectores
    \[
        v_1=(1,0,1),\qquad
        v_2=(0,1,1),\qquad
        v_3=(1,1,1)
    \]
    forman una base de \(\mathbb{R}^3\).

    \medskip

    Resuelve el problema de dos formas:

    \begin{enumerate}
        \item mediante eliminación gaussiana;
        \item calculando un determinante.
    \end{enumerate}
\end{frame}


\begin{frame}
    \frametitle{Ejercicio}

    Determina para qué valores de \(a\in\mathbb{R}\) los vectores
    \[
        v_1=(1,0,1),\qquad
        v_2=(0,1,1),\qquad
        v_3=(1,1,a)
    \]
    forman una base de \(\mathbb{R}^3\).

    \medskip

    Para los restantes valores de \(a\), determina:

    \begin{enumerate}
        \item una relación de dependencia lineal;
        \item cuál de los vectores puede eliminarse sin cambiar
        las combinaciones lineales que se pueden obtener.
    \end{enumerate}
\end{frame}

% \begin{frame}
    \frametitle{Ejercicio}

    \begin{exercise}
        Calcula la dimensión de los siguientes espacios vectoriales
        reales y da una base de cada uno:
        \begin{enumerate}
            \item $\mathbb{R}^3\times\mathbb{R}^2$;
            \item $M_{2\times3}(\mathbb{R})\times\mathbb{R}_3[x]$;
            \item $\mathbb{R}_1[x]\times\mathbb{R}_1[x]\times\mathbb{R}_1[x]$.
        \end{enumerate}
    \end{exercise}
\end{frame}


\begin{frame}[allowframebreaks]
    \frametitle{Ejercicio}

    \begin{solution}
        Aplicamos el teorema sobre bases del producto y su corolario.

        \begin{enumerate}
            \item $\dim=3+2=5$. Una base es
                  \[
                      \{(e_1,0),(e_2,0),(e_3,0),(0,e'_1),(0,e'_2)\},
                  \]
                  donde $e_i$ y $e'_j$ son los vectores canónicos de
                  $\mathbb{R}^3$ y $\mathbb{R}^2$. Con la identificación
                  $\mathbb{R}^3\times\mathbb{R}^2=\mathbb{R}^5$, es la base
                  canónica de $\mathbb{R}^5$.

            \item $\dim=6+4=10$. Una base está formada por los seis pares
                  $(E_{ij},0)$, con $i=1,2$ y $j=1,2,3$, y los cuatro pares
                  $(0,1),(0,x),(0,x^2),(0,x^3)$.

            \item $\dim=2+2+2=6$. Una base es
                  \[
                      (1,0,0),\ (x,0,0),\ (0,1,0),\ (0,x,0),\ (0,0,1),\ (0,0,x).\qedhere
                  \]
        \end{enumerate}
    \end{solution}
\end{frame}


\begin{frame}
    \frametitle{Ejercicio}

    \begin{exercise}
        En $\mathbb{R}^2\times\mathbb{R}_1[x]$ se consideran los vectores
        \[
            u_1=\bigl((1,0),\,1+x\bigr),
            \qquad
            u_2=\bigl((0,1),\,x\bigr),
        \]
        \[
            u_3=\bigl((1,1),\,0\bigr),
            \qquad
            u_4=\bigl((0,0),\,1\bigr).
        \]
        Decide si $\{u_1,u_2,u_3,u_4\}$ es una base de
        $\mathbb{R}^2\times\mathbb{R}_1[x]$.
    \end{exercise}

    \medskip

    \alert{Indicación:} trabaja con coordenadas respecto de la base
    \[
        B=\bigl\{((1,0),0),\ ((0,1),0),\ ((0,0),1),\ ((0,0),x)\bigr\}.
    \]
\end{frame}


\begin{frame}[allowframebreaks]
    \frametitle{Ejercicio}

    \begin{solution}
        Los vectores de coordenadas en $B$ son
        \[
            (u_1)_B=\begin{pmatrix}1\\0\\1\\1\end{pmatrix},\quad
            (u_2)_B=\begin{pmatrix}0\\1\\0\\1\end{pmatrix},\quad
            (u_3)_B=\begin{pmatrix}1\\1\\0\\0\end{pmatrix},\quad
            (u_4)_B=\begin{pmatrix}0\\0\\1\\0\end{pmatrix}.
        \]
        Sea $A$ la matriz $4\times4$ que tiene estos vectores por columnas.

        Dado $v\in\mathbb{R}^2\times\mathbb{R}_1[x]$, por las propiedades
        de la aplicación de coordenadas y la unicidad de las coordenadas,
        \[
            v=\lambda_1u_1+\cdots+\lambda_4u_4
            \iff
            v_B=\lambda_1(u_1)_B+\cdots+\lambda_4(u_4)_B
            \iff
            A\begin{pmatrix}\lambda_1\\\vdots\\\lambda_4\end{pmatrix}=v_B.
        \]

        \framebreak

        Escalonamos $A$:
        \[
            \begin{pmatrix}
                1&0&1&0\\
                0&1&1&0\\
                1&0&0&1\\
                1&1&0&0
            \end{pmatrix}
            \xrightarrow[F_4-F_1]{F_3-F_1}
            \begin{pmatrix}
                1&0&1&0\\
                0&1&1&0\\
                0&0&-1&1\\
                0&1&-1&0
            \end{pmatrix}
            \xrightarrow[F_4-2F_3]{F_4-F_2}
            \begin{pmatrix}
                1&0&1&0\\
                0&1&1&0\\
                0&0&-1&1\\
                0&0&0&-2
            \end{pmatrix}.
        \]
        (La segunda flecha indica que primero se hace $F_4-F_2$ y después
        $F_4-2F_3$.)

        Así, $\rg(A)=4$: hay un pivote en cada fila y en cada columna.
        Por tanto, para cada $v$, el sistema $A\lambda=v_B$ tiene
        \alert{solución única}, es decir, $v$ se escribe de forma única
        como combinación lineal de $u_1,u_2,u_3,u_4$.

        Por el teorema de representación única, $\{u_1,u_2,u_3,u_4\}$
        es una base de $\mathbb{R}^2\times\mathbb{R}_1[x]$.
    \end{solution}
\end{frame}


\begin{frame}
    \frametitle{Ejercicio}

    \begin{exercise}
        Sean $V$ y $W$ espacios vectoriales sobre $k$.
        \begin{enumerate}
            \item Demuestra que si $v_1,\ldots,v_r\in V$ son linealmente
                  independientes, entonces, para cualesquiera
                  $w_1,\ldots,w_r\in W$, los vectores
                  \[
                      (v_1,w_1),\ldots,(v_r,w_r)
                  \]
                  son linealmente independientes en $V\times W$.
            \item Muestra con un ejemplo que el recíproco es falso.
            \item Si $v_1,\ldots,v_r$ generan $V$ y $w_1,\ldots,w_r$
                  generan $W$, ¿generan necesariamente
                  $(v_1,w_1),\ldots,(v_r,w_r)$ el espacio $V\times W$?
        \end{enumerate}
    \end{exercise}
\end{frame}


\begin{frame}[allowframebreaks]
    \frametitle{Ejercicio}

    \begin{solution}
        \alert{(1)} Supongamos que
        $\lambda_1(v_1,w_1)+\cdots+\lambda_r(v_r,w_r)=(0,0)$.
        Calculando componente a componente,
        \[
            \bigl(\lambda_1v_1+\cdots+\lambda_rv_r,\;
                  \lambda_1w_1+\cdots+\lambda_rw_r\bigr)=(0,0).
        \]
        En particular, la primera componente da
        $\lambda_1v_1+\cdots+\lambda_rv_r=0$ en $V$, y la independencia
        lineal de $v_1,\ldots,v_r$ implica $\lambda_1=\cdots=\lambda_r=0$.
        No ha hecho falta usar la segunda componente.

        \alert{(2)} En $\mathbb{R}\times\mathbb{R}=\mathbb{R}^2$,
        tomemos $v_1=v_2=1$, $w_1=0$ y $w_2=1$. Los vectores
        \[
            (v_1,w_1)=(1,0),
            \qquad
            (v_2,w_2)=(1,1)
        \]
        son linealmente independientes, pues
        $\rg\begin{pmatrix}1&1\\0&1\end{pmatrix}=2$.
        Sin embargo, $v_1=v_2$, así que $v_1-v_2=0$ es una relación
        de dependencia lineal entre $v_1,v_2$ en $\mathbb{R}$.

        \framebreak

        \alert{(3)} No. En $\mathbb{R}\times\mathbb{R}=\mathbb{R}^2$,
        tomemos $r=1$ y $v_1=w_1=1$. Entonces $v_1$ genera $\mathbb{R}$
        y $w_1$ genera $\mathbb{R}$, pero el único vector
        \[
            (v_1,w_1)=(1,1)
        \]
        no puede generar $\mathbb{R}^2$, ya que para generar $k^n$
        hacen falta al menos $n$ vectores. Explícitamente, las
        combinaciones lineales de $(1,1)$ son los vectores $(\lambda,\lambda)$,
        y $(1,0)$ no es de esa forma.

        \medskip

        Lo que sí genera $V\times W$ es la familia
        \[
            (v_1,0),\ldots,(v_r,0),\,(0,w_1),\ldots,(0,w_r),
        \]
        como muestra la primera parte de la demostración del teorema
        sobre bases del producto.
    \end{solution}
\end{frame}


\section{Subespacios}

\begin{frame}
    \frametitle{Subespacios vectoriales}

    \begin{definition}
        Sea $V$ un espacio vectorial sobre $k$. Un subconjunto \textbf{no vacío}
        $W\subseteq V$ es un \alert{subespacio vectorial} de $V$ si es
        cerrado para la suma y el producto por escalares:
        \begin{enumerate}
            \item para todo $u,v\in W$, se tiene $u+v\in W$,
            \item para todo $\lambda\in k$ y $v\in W$, se tiene $\lambda v\in W$.
        \end{enumerate}
    \end{definition}

\end{frame}

\begin{frame}
    \frametitle{Subespacios como espacios vectoriales}

    \begin{proposition}
        Todo subespacio $W\subseteq V$ es en sí mismo un espacio vectorial
        con las operaciones heredadas de $V$.
    \end{proposition}

    \begin{proof}
        Tomando $\lambda=0$ y un vector cualquier $v\in W$ en (2), lo cual es posible porque $W$ no es vacío, se obtiene $0\in W$.
        Tomando $\lambda=-1$, se obtiene $-v\in W$ para todo $v\in W$.
        Los axiomas restantes —asociatividad y conmutatividad de la suma,
        distributividades y compatibilidad del producto por escalares—
        valen en $W$ por ser válidos en $V$ para cualesquiera vectores.
    \end{proof}
\end{frame}

\begin{frame}
    \frametitle{Una condición necesaria}
    
    \begin{corollary}
        Si $W\subseteq V$ es un subespacio, entonces $0\in W$.
    \end{corollary}

    \medskip

    Esto tiene dos usos prácticos:
    \begin{itemize}
        \item para ver que $W$ \alert{no es vacío}, lo habitual es comprobar
        que $0\in W$;
        \item si $0\notin W$, podemos concluir directamente que $W$
        \alert{no} es un subespacio.
    \end{itemize}

    \medskip

    El recíproco es falso: veremos subconjuntos que contienen al $0$
    y no son subespacios.
\end{frame}

\begin{frame}
    \frametitle{Caracterización mediante combinaciones lineales}

    \begin{proposition}
        Un subconjunto no vacío $W\subseteq V$ es un subespacio si y solo si
        para todo $u,w\in W$ y todo $\lambda,\mu\in k$ se tiene
        \[
            \lambda u+\mu w\in W.
        \]
    \end{proposition}

    \begin{proof}
        Si $W$ es un subespacio, entonces $\lambda u\in W$ y $\mu w\in W$
        por la condición (2), y su suma está en $W$ por la condición (1).

        \medskip

        Recíprocamente, supongamos que $W$ cumple la condición del enunciado.
        Tomando $\lambda=\mu=1$ obtenemos $u+w\in W$, que es la condición (1).
        Tomando $\mu=0$ obtenemos $\lambda u=\lambda u+0w\in W$, que es la
        condición (2).
    \end{proof}
\end{frame}

\begin{frame}
    \frametitle{Cierre para combinaciones lineales}

    \begin{proposition}
        Si $W\subseteq V$ es un subespacio, $v_1,\ldots,v_r\in W$ y
        $\lambda_1,\ldots,\lambda_r\in k$, entonces
        \[
            \lambda_1v_1+\cdots+\lambda_rv_r\in W.
        \]
    \end{proposition}

    \vspace{-10pt}

    \begin{proof}
        Por inducción sobre $r$. Para $r=1$, $\lambda_1v_1\in W$ por la
        condición (2). Si $r>1$, por hipótesis de inducción
        $u=\lambda_1v_1+\cdots+\lambda_{r-1}v_{r-1}\in W$, y entonces
        \[
            \lambda_1v_1+\cdots+\lambda_rv_r=u+\lambda_rv_r\in W
        \]
        por las condiciones (2) y (1).
    \end{proof}

    \medskip

    En otras palabras: un subespacio es un subconjunto no vacío
    \alert{cerrado para combinaciones lineales}.
\end{frame}

\begin{frame}
    \frametitle{Los subespacios triviales}

    \begin{example}
        Todo espacio vectorial $V$ tiene, al menos, dos subespacios:
        \begin{itemize}
            \item el \alert{total}, que es el propio $V$;
            \item el \alert{subespacio trivial}, \alert{cero} o \alert{nulo} $\{0\}$.
        \end{itemize}
    \end{example}

    \medskip

    Para $\{0\}$: no es vacío y, como el único vector es $0$,
    \[
        0+0=0\in\{0\},
        \qquad
        \lambda 0=0\in\{0\}
        \quad\text{para todo }\lambda\in k.
    \]

    \medskip

    Si $V\neq\{0\}$, estos dos subespacios son distintos.
    Un subespacio $W\neq V$ se llama \alert{propio}, y uno $W\neq\{0\}$
    se llama \alert{no trivial}.
\end{frame}

\begin{frame}
    \frametitle{Rectas que pasan por el origen}

    \begin{example}
        Sea $u\in V$ un vector no nulo. El conjunto
        \[
            W=\{\lambda u:\lambda\in k\}
        \]
        de los múltiplos de $u$ es un subespacio de $V$.
    \end{example}

    \medskip

    En efecto, $0=0u\in W$. Si $\lambda u,\mu u\in W$ y $\alpha\in k$, entonces
    \[
        \lambda u+\mu u=(\lambda+\mu)u\in W,
        \qquad
        \alpha(\lambda u)=(\alpha\lambda)u\in W.
    \]

    \medskip

    El subespacio $W$ es la \alert{recta} que pasa por el origen
    con dirección $u$. Esto se ve con más claridad en el caso $V=\R^2$ o $V=\R^3$, pero es cierto en cualquier espacio vectorial.
\end{frame}

\begin{frame}
    \frametitle{Un subespacio de $k^n$}

    \begin{example}
        El conjunto
        \[
            W=\{(v_1,\ldots,v_n)\in k^n:\ v_1+\cdots+v_n=0\}
        \]
        es un subespacio de $k^n$.
    \end{example}

    \medskip

    Se tiene $0=(0,\ldots,0)\in W$, pues $0+\cdots+0=0$.
    Sean $u=(u_1,\ldots,u_n)$ y $w=(w_1,\ldots,w_n)$ en $W$, es decir,
    \[
        u_1+\cdots+u_n=0,
        \qquad
        w_1+\cdots+w_n=0.
    \]
    Entonces $u+w=(u_1+w_1,\ldots,u_n+w_n)$ y
    \[
        (u_1+w_1)+\cdots+(u_n+w_n)
        =(u_1+\cdots+u_n)+(w_1+\cdots+w_n)=0+0=0,
    \]
    luego $u+w\in W$. Análogamente, para $\lambda\in k$,
    \[
        \lambda u_1+\cdots+\lambda u_n=\lambda(u_1+\cdots+u_n)=\lambda0=0,
    \]
    luego $\lambda u\in W$.
\end{frame}

\begin{frame}
    \frametitle{Subespacios de polinomios}

    \begin{example}
        Los siguientes conjuntos son subespacios de $k[x]$:
        \begin{enumerate}
            \item $k_n[x]$, para cada $n\geq0$;
            \item $W_a=\{p(x)\in k[x]:\ p(a)=0\}$, el espacio de polinomios que tienen raíz $a\in k$ fijada.
        \end{enumerate}
    \end{example}

    \medskip

    (1) El polinomio nulo está en $k_n[x]$, y si $p,q$ tienen grado
    $\leq n$, entonces $p+q$ y $\lambda p$ también, pues sumar o multiplicar
    por un escalar no crea monomios de grado mayor que $n$.

    \medskip

    (2) El polinomio nulo se anula en $a$. Si $p(a)=q(a)=0$ y $\lambda\in k$,
    \[
        (p+q)(a)=p(a)+q(a)=0,
        \qquad
        (\lambda p)(a)=\lambda p(a)=0.
    \]

    \medskip

    Además, $k_m[x]$ es un subespacio de $k_n[x]$ si $m\leq n$.
\end{frame}

\begin{frame}
    \frametitle{Subespacios de matrices}

    \begin{example}
        Los siguientes conjuntos son subespacios de $M_n(k)$:
        \begin{enumerate}
            \item las \alert{matrices simétricas}, $\{A\in M_n(k):\ A^t=A\}$;
            \item las \alert{matrices triangulares superiores},
            $\{A=(a_{ij})\in M_n(k):\ a_{ij}=0 \text{ si } i>j\}$.
        \end{enumerate}
    \end{example}

    \medskip

    (1) La matriz nula es simétrica. Si $A^t=A$, $B^t=B$ y $\lambda\in k$,
    \[
        (A+B)^t=A^t+B^t=A+B,
        \qquad
        (\lambda A)^t=\lambda A^t=\lambda A.
    \]

    \medskip

    (2) La matriz nula es triangular superior. Si $i>j$, las entradas
    $(i,j)$ de $A+B$ y de $\lambda A$ son
    \[
        a_{ij}+b_{ij}=0+0=0,
        \qquad
        \lambda a_{ij}=\lambda0=0.
    \]
\end{frame}

\begin{frame}
    \frametitle{Subespacios de funciones}

    \begin{example}
        En el espacio $\R^\R$ de todas las funciones $f\colon\R\to\R$,
        los siguientes conjuntos son subespacios:
        \begin{enumerate}
            \item las funciones \alert{continuas};
            \item las funciones \alert{derivables};
            \item las funciones \alert{polinómicas}: $\R[x]$.
        \end{enumerate}
    \end{example}

    \medskip

    En los tres casos, la función constante $0$ pertenece al conjunto,
    y la suma de dos funciones del tipo indicado y el producto de una de
    ellas por un número real vuelven a ser del mismo tipo.

    \medskip

    Cada uno de ellos es un subespacio del anterior:
    \[
        \R[x]=\{\text{polinómicas}\}\subseteq\{\text{derivables}\}
        \subseteq\{\text{continuas}\}\subseteq\R^\R.
    \]
\end{frame}

\begin{frame}
    \frametitle{Contraejemplo: rectas que no pasan por el origen}

    \begin{example}
        El conjunto
        \[
            W=\{(v_1,v_2)\in\R^2:\ v_2=v_1+1\}
        \]
        \alert{no} es un subespacio de $\R^2$.
    \end{example}

    \begin{columns}[T]
        \begin{column}{0.58\textwidth}
            Basta observar que $(0,0)\notin W$, pues $0\neq0+1$.

            \medskip

            De hecho $W$ falla también las dos condiciones de cierre:
            $(0,1),(1,2)\in W$, pero
            \[
                (0,1)+(1,2)=(1,3)\notin W,
            \]
            \[
                2(0,1)=(0,2)\notin W.
            \]

            Las rectas de $\R^2$ que son subespacios son exactamente las que
            \alert{pasan por el origen}.
        \end{column}
        \begin{column}{0.38\textwidth}
            \centering
            \begin{tikzpicture}[scale=0.75]
                \draw[->] (-2.2,0) -- (2.4,0) node[below] {\footnotesize $v_1$};
                \draw[->] (0,-1.2) -- (0,3.6) node[left] {\footnotesize $v_2$};
                \draw[very thick,alertcolor] (-2.2,-1.2) -- (2.4,3.4)
                    node[right] {$W$};
                \fill (0,0) circle (2pt) node[below left] {\footnotesize $0$};
                \fill[alertcolor] (0,1) circle (2.5pt)
                    node[left] {\footnotesize $(0,1)$};
                \fill[alertcolor] (1,2) circle (2.5pt)
                    node[right] {\footnotesize $(1,2)$};
                \draw[thick] (1,3) circle (2.5pt)
                    node[above] {\footnotesize $(1,3)$};
                \draw[thick] (0,2) circle (2.5pt)
                    node[left] {\footnotesize $(0,2)$};
            \end{tikzpicture}
        \end{column}
    \end{columns}
\end{frame}

\begin{frame}
    \frametitle{Contraejemplo: cerrado para escalares, no para la suma}

    \begin{example}
        El conjunto $W=\{(v_1,v_2)\in\R^2:\ v_1v_2=0\}$, formado por los dos
        ejes de coordenadas, \alert{no} es un subespacio de $\R^2$.
    \end{example}

    \begin{columns}[T]
        \begin{column}{0.62\textwidth}
            Contiene a $(0,0)$ y es cerrado para el producto por escalares:
            si $v_1v_2=0$, entonces $(\lambda v_1)(\lambda v_2)=\lambda^2v_1v_2=0$.

            \medskip

            Sin embargo, no es cerrado para la suma: $(1,0),(0,1)\in W$, pero
            \[
                (1,0)+(0,1)=(1,1)\notin W.
            \]

            Observa que $W$ es la \alert{unión} de dos subespacios, los dos ejes.
            La unión de subespacios no es, en general, un subespacio.
        \end{column}
        \begin{column}{0.34\textwidth}
            \centering
            \begin{tikzpicture}[scale=0.85]
                \draw[very thick,alertcolor] (-1.6,0) -- (2,0);
                \draw[very thick,alertcolor] (0,-1.2) -- (0,2);
                \node[alertcolor,right] at (2,0) {$W$};
                \draw[dashed] (1,0) -- (1,1) -- (0,1);
                \fill[alertcolor] (1,0) circle (2.5pt)
                    node[below] {\footnotesize $(1,0)$};
                \fill[alertcolor] (0,1) circle (2.5pt)
                    node[left] {\footnotesize $(0,1)$};
                \fill[white] (1,1) circle (2.5pt);
                \draw[thick] (1,1) circle (2.5pt)
                    node[above right] {\footnotesize $(1,1)$};
            \end{tikzpicture}
        \end{column}
    \end{columns}
\end{frame}

\begin{frame}
    \frametitle{Contraejemplo: cerrado para la suma, no para escalares}

    \begin{example}
        El conjunto
        \[
            W=\{(v_1,v_2)\in\R^2:\ v_1\geq0,\ v_2\geq0\}
        \]
        (el primer cuadrante) \alert{no} es un subespacio de $\R^2$.
    \end{example}

    \begin{columns}[T]
        \begin{column}{0.62\textwidth}
            Contiene a $(0,0)$ y es cerrado para la suma, pues la suma de números
            no negativos es no negativa.

            \medskip

            Sin embargo, no es cerrado para el producto por escalares:
            \[
                (1,0)\in W,
                \qquad\text{pero}\qquad
                (-1)(1,0)=(-1,0)\notin W.
            \]

            Lo mismo ocurre con $\Z^2\subseteq\R^2$: es cerrado para la suma,
            pero $\tfrac12(1,0)\notin\Z^2$.
        \end{column}
        \begin{column}{0.34\textwidth}
            \centering
            \begin{tikzpicture}[scale=0.85]
                \fill[alertcolor!20] (0,0) rectangle (2,2);
                \draw[->] (-1.6,0) -- (2.2,0);
                \draw[->] (0,-1.2) -- (0,2.2);
                \draw[very thick,alertcolor] (2,0) -- (0,0) -- (0,2);
                \node[alertcolor] at (1.3,1.3) {$W$};
                \draw[dashed] (1,0) -- (-1,0);
                \fill[alertcolor] (1,0) circle (2.5pt)
                    node[below] {\footnotesize $(1,0)$};
                \fill[white] (-1,0) circle (2.5pt);
                \draw[thick] (-1,0) circle (2.5pt)
                    node[below] {\footnotesize $(-1,0)$};
            \end{tikzpicture}
        \end{column}
    \end{columns}
\end{frame}

\begin{frame}
    \frametitle{Contraejemplos en polinomios y matrices}

    \begin{example}
        Ninguno de los siguientes conjuntos es un subespacio:
        \begin{enumerate}
            \item los polinomios de grado \alert{exactamente} $n\geq1$, en $k[x]$;
            \item las matrices invertibles, en $M_n(k)$.
        \end{enumerate}
    \end{example}

    \medskip

    En ambos casos falta el vector cero: el polinomio nulo no tiene grado $n$
    y la matriz nula no es invertible.

    \medskip

    Tampoco son cerrados para la suma. Por ejemplo,
    \[
        x^n+(1-x^n)=1,
    \]
    que tiene grado $0$, y
    \[
        I+(-I)=0,
    \]
    que no es invertible.
\end{frame}

\begin{frame}[allowframebreaks, t]
    \frametitle{Soluciones de un sistema homogéneo}

    \begin{proposition}
        Sea $A\in M_{m\times n}(k)$. El conjunto de soluciones del sistema
        homogéneo $Av=0$,
        \[
            W=\{v\in k^n:\ Av=0\},
        \]
        es un subespacio de $k^n$ y se denota
        \[W\colon\ Ax=0.\]
    \end{proposition}

    Decimos que $Ax=0$ son unas \alert{ecuaciones implícitas} para el subespacio $W$.

    \framebreak

    \begin{proof}
        Es lo que vimos en el Tema 1. El vector nulo es solución, pues $A0=0$,
        luego $W\neq\varnothing$. Si $u,w\in W$ y $\lambda\in k$, entonces
        \[
            A(u+w)=Au+Aw=0+0=0,
            \qquad
            A(\lambda u)=\lambda Au=\lambda0=0,
        \]
        es decir, $u+w\in W$ y $\lambda u\in W$.
    \end{proof}


    Por ejemplo, el subespacio $\{v\in k^n: v_1+\cdots+v_n=0\}$ es el
    conjunto de soluciones de un sistema homogéneo con una sola ecuación.
\end{frame}

\begin{frame}
    \frametitle{Sistemas no homogéneos}

    \begin{example}
        Si $b\neq0$, el conjunto de soluciones de $Av=b$
        \alert{no} es un subespacio de $k^n$.
    \end{example}

    \medskip

    En efecto, $A0=0\neq b$, luego $0$ no es solución.

    \medskip

    \begin{columns}[T]
        \begin{column}{0.56\textwidth}
            La recta $\{(v_1,v_2)\in\R^2: v_2=v_1+1\}$ es de este tipo:
            es el conjunto de soluciones del sistema no homogéneo
            \[
                -v_1+v_2=1.
            \]
            En cambio, la recta paralela que pasa por el origen,
            $\{(v_1,v_2)\in\R^2: -v_1+v_2=0\}$, es el conjunto de soluciones
            del \alert{sistema homogéneo asociado}, y sí es un subespacio.
        \end{column}
        \begin{column}{0.40\textwidth}
            \centering
            \begin{tikzpicture}[scale=0.7]
                \draw[->] (-2.2,0) -- (2.4,0) node[below] {\footnotesize $v_1$};
                \draw[->] (0,-1.6) -- (0,3) node[left] {\footnotesize $v_2$};
                \draw[very thick,alertcolor] (-2.2,-1.2) -- (1.8,2.8)
                    node[above] {\footnotesize $-v_1+v_2=1$};
                \draw[very thick] (-1.6,-1.6) -- (2.2,2.2)
                    node[right] {\footnotesize $-v_1+v_2=0$};
                \fill (0,0) circle (2pt) node[below right] {\footnotesize $0$};
            \end{tikzpicture}
        \end{column}
    \end{columns}
\end{frame}


\begin{frame}[allowframebreaks]
    \frametitle{Ejemplo: un plano en $\R^3$}

    \begin{example}
        Consideremos el subespacio de $\R^3$
        \[
            W\colon\ v_1-2v_2+v_3=0.
        \]
        Es el conjunto de soluciones de un sistema homogéneo con una ecuación
        y tres incógnitas.
    \end{example}

    La matriz de coeficientes $\begin{pmatrix}1&-2&1\end{pmatrix}$ ya está
    escalonada, con pivote en la primera columna. Las incógnitas libres son
    $v_2$ y $v_3$. Tomando $v_2=\lambda$, $v_3=\mu$, se obtiene
    $v_1=2\lambda-\mu$, luego
    \[
        W=\{(2\lambda-\mu,\lambda,\mu):\ \lambda,\mu\in\R\}
        =\{\lambda(2,1,0)+\mu(-1,0,1):\ \lambda,\mu\in\R\}.
    \]

    \framebreak

    \begin{columns}[T]
        \begin{column}{0.52\textwidth}
            Así, los elementos de $W$ son exactamente las
            \alert{combinaciones lineales} de
            \[
                u_1=(2,1,0),
                \qquad
                u_2=(-1,0,1).
            \]
            Geométricamente, $W$ es el \alert{plano} de $\R^3$ que pasa
            por el origen y contiene a $u_1$ y $u_2$.
        \end{column}
        \begin{column}{0.44\textwidth}
            \centering
            \begin{tikzpicture}[x={(1cm,0cm)},y={(0.4cm,0.3cm)},
                    z={(0cm,1cm)},scale=0.66]
                \filldraw[alertcolor,fill opacity=0.15,draw opacity=0.6]
                    (-1,-1,-1) -- (3,1,-1) -- (1,1,1) -- (-3,-1,1) -- cycle;
                \draw[->] (0,0,0) -- (3.6,0,0) node[right] {\footnotesize $v_1$};
                \draw[->] (0,0,0) -- (0,3,0) node[right] {\footnotesize $v_2$};
                \draw[->] (0,0,0) -- (0,0,2) node[above] {\footnotesize $v_3$};
                \draw[->,very thick,alertcolor] (0,0,0) -- (2,1,0)
                    node[above] {\footnotesize $u_1$};
                \draw[->,very thick,alertcolor] (0,0,0) -- (-1,0,1)
                    node[above left] {\footnotesize $u_2$};
                \fill (0,0,0) circle (1.5pt);
                \node[alertcolor] at (-2.3,-0.8,0.3) {$W$};
            \end{tikzpicture}
        \end{column}
    \end{columns}

    \medskip

    Además, por lo visto en el Tema 1, cada vector de $W$ se escribe de
    manera \alert{única} como combinación lineal de $u_1$ y $u_2$, de modo
    que $u_1,u_2$ es una base de $W$ y $\dim W=2$.

    \medskip

    Esta descripción de un subespacio como el conjunto de combinaciones
    lineales de ciertos vectores es la idea del \alert{subespacio generado},
    que estudiamos a continuación.
\end{frame}


\begin{frame}
    \frametitle{Dimensión del espacio de soluciones}

    \begin{proposition}
        Sea $A\in M_{m\times n}(k)$ y sea $W\colon Ax=0$. Entonces
        \[
            \dim W=n-\rg(A).
        \]
    \end{proposition}

    \begin{proof}
        Al escalonar $A$ hay $\rg(A)$ pivotes, luego $s=n-\rg(A)$
        incógnitas libres. Sean $v_1,\ldots,v_s$ las soluciones obtenidas
        dando el valor $1$ a una incógnita libre y $0$ a las demás.
        Por el Tema~1, toda solución se escribe de manera \alert{única} como
        $\lambda_1v_1+\cdots+\lambda_sv_s$. Así, $v_1,\ldots,v_s$ generan $W$.
        Además son linealmente independientes: si
        $\lambda_1v_1+\cdots+\lambda_sv_s=0$, como también
        $0=0v_1+\cdots+0v_s$, la unicidad da $\lambda_1=\cdots=\lambda_s=0$.
        Por tanto, $v_1,\ldots,v_s$ es una base de $W$ y $\dim W=s$.
    \end{proof}
\end{frame}


\begin{frame}
    \frametitle{Ecuaciones independientes}

    \begin{definition}
        Sea $A\in M_{m\times n}(k)$. Las $m$ ecuaciones del sistema
        homogéneo $Av=0$ son \alert{independientes} si $\rg(A)=m$.
    \end{definition}

    \medskip

    Como $\rg(A)=\rg(A^t)$ y las columnas de $A^t$ son las filas de $A$,
    esto equivale a que las filas de $A$ sean vectores linealmente
    independientes de $k^n$: ninguna ecuación es combinación lineal
    de las demás.

    \medskip

    Por la proposición anterior, si $W$ viene dado por $m$ ecuaciones
    independientes, entonces
    \[
        \dim W=n-m,
        \qquad\text{es decir,}\qquad
        m=n-\dim W.
    \]
    Cualquier otro sistema que defina $W$ tiene al menos $n-\dim W$
    ecuaciones, pues su número de ecuaciones es $\geq$ que su rango.
\end{frame}


\begin{frame}
    \frametitle{Cómo obtener ecuaciones independientes}

    \begin{proposition}
        Todo subespacio $W\colon Ax=0$ de $k^n$ puede describirse mediante
        $\rg(A)$ ecuaciones independientes.
    \end{proposition}

    \begin{proof}
        Sea $E$ una forma escalonada de $A$. Por el Tema~1, las operaciones
        elementales por filas no cambian las soluciones, luego
        $W\colon Ex=0$. Las filas nulas de $E$ dan la ecuación
        $0=0$ y pueden suprimirse. Las $\rg(A)$ filas restantes forman una
        matriz escalonada sin filas nulas, cuyo rango es su número de
        filas.
    \end{proof}

    \medskip

    Es decir: para eliminar ecuaciones \alert{redundantes}, basta escalonar
    la matriz de coeficientes y quedarse con las filas no nulas.
\end{frame}


\begin{frame}
    \frametitle{Ejemplo: ecuaciones redundantes}

    \begin{example}
        Sea $W\subseteq\R^3$ el subespacio definido por
        \[
            \left\{
            \begin{aligned}
                v_1+v_2-v_3&=0,\\
                v_1-v_2+v_3&=0,\\
                2v_1&=0.
            \end{aligned}
            \right.
        \]
    \end{example}

    La tercera ecuación es la suma de las dos primeras, así que las
    ecuaciones no son independientes. Escalonando:
    \[
        \begin{pmatrix}1&1&-1\\1&-1&1\\2&0&0\end{pmatrix}
        \xrightarrow[F_3-2F_1]{F_2-F_1}
        \begin{pmatrix}1&1&-1\\0&-2&2\\0&-2&2\end{pmatrix}
        \xrightarrow{F_3-F_2}
        \begin{pmatrix}1&1&-1\\0&-2&2\\0&0&0\end{pmatrix}.
    \]
    El rango es $2$, luego $\dim W=3-2=1$, y $W$ viene dado por las dos
    ecuaciones independientes $v_1+v_2-v_3=0$ y $-2v_2+2v_3=0$.
    Resolviendo, $W=\langle(0,1,1)\rangle$.
\end{frame}


\begin{frame}
    \frametitle{Sistemas homogéneos en coordenadas}

    Sea $V$ un espacio vectorial de dimensión $n$ y sea $B$ una base de $V$.
    Recordemos que $u_B\in k^n$ denota el vector de coordenadas de $u\in V$
    en $B$.

    \begin{definition}
        Sea $A\in M_{m\times n}(k)$. El subespacio de $V$
        con \alert{ecuaciones implícitas $Ax=0$ respecto de $B$} es
        \[
            W=\{u\in V:\ Au_B=0\},
        \]
        es decir, el conjunto de vectores de $V$ cuyas coordenadas en $B$
        son soluciones de $Ax=0$. Lo denotamos
        \[
            W\colon\ Ax=0\quad\text{resp.\ }B.
        \]
    \end{definition}

    \medskip

    Las incógnitas $x_1,\ldots,x_n$ representan las \alert{coordenadas}
    en $B$ de un vector genérico de~$V$.
\end{frame}

\begin{frame}
    \frametitle{Sistemas homogéneos en coordenadas}

    \begin{proposition}
        Si $W\colon Ax=0$ resp.\ $B$, entonces $W$ es un subespacio de $V$.
    \end{proposition}

    \begin{proof}
        Como $0_B=0$ y $A0=0$, se tiene $0\in W$. Sean $u,w\in W$ y
        $\lambda\in k$. Por las propiedades de las coordenadas del Tema~2,
        \[
            A(u+w)_B=A(u_B+w_B)=Au_B+Aw_B=0+0=0,
        \]
        \[
            A(\lambda u)_B=A(\lambda u_B)=\lambda Au_B=\lambda0=0,
        \]
        luego $u+w\in W$ y $\lambda u\in W$.
    \end{proof}

    \medskip

    En $k^n$ con la base canónica $C$, se tiene $u_C=u$, así que
    $W\colon Ax=0$ resp.\ $C$ es simplemente el subespacio $W\colon Ax=0$
    de $k^n$.
\end{frame}


\begin{frame}
    \frametitle{Ecuaciones independientes respecto de una base}

    Como en $k^n$, si $A\in M_{m\times n}(k)$, decimos que las $m$
    ecuaciones de $W\colon Ax=0$ resp.\ $B$ son \alert{independientes}
    si $\rg(A)=m$.

    \begin{proposition}
        Si $W\colon Ax=0$ resp.\ $B$, entonces $\dim W=n-\rg(A)$.
        En particular, si las ecuaciones son independientes,
        $\dim W=n-m$.
    \end{proposition}

    \begin{proof}
        Sea $W'\colon Ax=0$ en $k^n$, de dimensión $n-\rg(A)$. Por
        definición, $u\in W$ si y solo si $u_B\in W'$. Todo vector de
        $k^n$ es el vector de coordenadas de algún vector de $V$, luego
        existen $u_1,\ldots,u_s\in V$ tales que $(u_1)_B,\ldots,(u_s)_B$ es
        una base de $W'$. Entonces $u_1,\ldots,u_s\in W$ son linealmente
        independientes y generan $W$, pues las coordenadas en $B$
        reflejan las combinaciones lineales (Tema~2). Luego
        $\dim W=\dim W'$.
    \end{proof}

    \medskip

    Igual que en $k^n$, escalonando $A$ y suprimiendo las filas nulas
    se obtienen ecuaciones independientes de $W$ respecto de $B$.
\end{frame}

\begin{frame}
    \frametitle{Ejemplo: un subespacio de polinomios}

    \begin{example}
        Sea $W=\{p(x)\in k_2[x]:\ p(1)=0\}$ el espacio de polinomios que tienen raíz $1$, y sea $B=\{1,x,x^2\}$ la base habitual de $k_2[x]$.
    \end{example}

    \medskip

    Si $p(x)=a_0+a_1x+a_2x^2$, entonces $p_B=(a_0,a_1,a_2)$ y
    \[
        p(1)=a_0+a_1+a_2.
    \]
    Por tanto, $p\in W$ si y solo si sus coordenadas en $B$ satisfacen
    $a_0+a_1+a_2=0$, es decir,
    \[
        W\colon\
        \begin{pmatrix}1&1&1\end{pmatrix}
        \begin{pmatrix}a_0\\a_1\\a_2\end{pmatrix}=0
        \quad\text{resp.\ }B.
    \]
\end{frame}

\begin{frame}
    \frametitle{Ejemplo: subespacios de matrices}

    \begin{example}
        En $M_2(k)$ consideramos la base
        \[
            B=\left\{
            \begin{pmatrix}1&0\\0&0\end{pmatrix},
            \begin{pmatrix}0&1\\0&0\end{pmatrix},
            \begin{pmatrix}0&0\\1&0\end{pmatrix},
            \begin{pmatrix}0&0\\0&1\end{pmatrix}
            \right\}.
        \]
    \end{example}

    Las coordenadas de $\begin{pmatrix}a&b\\c&d\end{pmatrix}$ en $B$ son
    $(a,b,c,d)$. Por tanto:
    \begin{itemize}
        \item las matrices \alert{simétricas} son las que cumplen $b=c$:
        \[
            \begin{pmatrix}0&1&-1&0\end{pmatrix}x=0\quad\text{resp.\ }B;
        \]
        \item las matrices \alert{triangulares superiores} son las que
        cumplen $c=0$:
        \[
            \begin{pmatrix}0&0&1&0\end{pmatrix}x=0\quad\text{resp.\ }B.
        \]
    \end{itemize}
\end{frame}

\begin{frame}
    \frametitle{Cambio de base en el sistema}

    \begin{proposition}
        Sean $B$ y $B'$ dos bases de $V$. Si
        \[
            W\colon\ Ax=0\quad\text{resp.\ }B,
        \]
        entonces
        \[
            W\colon\ A\,M(B',B)x=0\quad\text{resp.\ }B'.
        \]
    \end{proposition}

    \vspace{-10pt}

    \begin{proof}
        Por el Tema~2, para todo $u\in V$ se tiene $u_B=M(B',B)\,u_{B'}$.
        Por tanto,
        \[
            u\in W
            \iff Au_B=0
            \iff A\,M(B',B)\,u_{B'}=0.\qedhere
        \]
    \end{proof}

    Para cambiar de base, basta \alert{multiplicar por la derecha} la matriz
    de coeficientes por la matriz de cambio de base de $B'$ a $B$.
    El número de ecuaciones no cambia.
\end{frame}

\begin{frame}
    \frametitle{Ejemplo: cambio de base en polinomios}

    \begin{example}
        Sea $W=\{p(x)\in k_2[x]:\ p(1)=0\}$. Hemos visto que
        \[
            W\colon\ \begin{pmatrix}1&1&1\end{pmatrix}x=0
            \quad\text{resp.\ }B=\{1,x,x^2\}.
        \]
        Calculemos sus ecuaciones respecto de $B'=\{1,\ x-1,\ (x-1)^2\}$.
    \end{example}

    Como $(x-1)^2=1-2x+x^2$, las columnas de $M(B',B)$ son las coordenadas
    en $B$ de los vectores de $B'$:
    \[
        \begin{pmatrix}1&1&1\end{pmatrix}
        \begin{pmatrix}1&-1&1\\0&1&-2\\0&0&1\end{pmatrix}
        =\begin{pmatrix}1&0&0\end{pmatrix}.
    \]
    Por tanto,
    \[
        W\colon\ \begin{pmatrix}1&0&0\end{pmatrix}x=0
        \quad\text{resp.\ }B'.
    \]
    En efecto, si $p(x)=b_0+b_1(x-1)+b_2(x-1)^2$, entonces $p(1)=b_0$.
\end{frame}

\begin{frame}
    \frametitle{Ejemplo: una base adaptada a un plano}

    \begin{example}
        Sea $W\subseteq\R^3$ el plano $W\colon\ x_1-2x_2+x_3=0$ resp.\ la
        base canónica $C$, y sea
        \[
            B'=\{(2,1,0),\ (-1,0,1),\ (1,0,0)\},
        \]
        cuyos dos primeros vectores generan $W$.
    \end{example}

    Las columnas de $M(B',C)$ son los propios vectores de $B'$:
    \[
        \begin{pmatrix}1&-2&1\end{pmatrix}
        \begin{pmatrix}2&-1&1\\1&0&0\\0&1&0\end{pmatrix}
        =\begin{pmatrix}0&0&1\end{pmatrix}.
    \]
    ($B'$ es una base, pues esta matriz tiene determinante $1\neq0$.)
    Por tanto,
    \[
        W\colon\ x_3=0\quad\text{resp.\ }B'.
    \]
    En una base que contiene generadores de $W$, la ecuación de $W$ se
    vuelve \alert{mucho más sencilla}.
\end{frame}

\subsection{Subespacio generado}

\begin{frame}
    \frametitle{Subespacio generado}

    \begin{definition}
        Sean $v_1,\ldots,v_r$ vectores de un espacio vectorial $V$ sobre $k$.
        El \alert{subespacio generado} por $v_1,\ldots,v_r$ es el conjunto
        de todas sus combinaciones lineales:
        \[
            \langle v_1,\ldots,v_r\rangle
            =\{\lambda_1v_1+\cdots+\lambda_rv_r:\ \lambda_1,\ldots,\lambda_r\in k\}.
        \]
    \end{definition}

    \medskip

    El subespacio generado por el conjunto vacío es el trivial
    $\langle\varnothing\rangle=\{0\}$.

    \medskip

    El plano $W=\{(v_1,v_2,v_3)\in\R^3: v_1-2v_2+v_3=0\}$
    que acabamos de estudiar es
    \[
        W=\langle(2,1,0),(-1,0,1)\rangle.
    \]
\end{frame}

\begin{frame}
    \frametitle{El subespacio generado es un subespacio}

    \begin{proposition}
        $\langle v_1,\ldots,v_r\rangle$ es un subespacio de $V$ que contiene
        a $v_1,\ldots,v_r$.
    \end{proposition}

    \begin{proof}
        Contiene al vector cero, $0=0v_1+\cdots+0v_r$, luego no es vacío.
        Sean
        \[
            u=\lambda_1v_1+\cdots+\lambda_rv_r,
            \qquad
            w=\mu_1v_1+\cdots+\mu_rv_r
        \]
        dos de sus elementos y $\alpha\in k$. Agrupando términos,
        \[
            u+w=(\lambda_1+\mu_1)v_1+\cdots+(\lambda_r+\mu_r)v_r,
            \qquad
            \alpha u=(\alpha\lambda_1)v_1+\cdots+(\alpha\lambda_r)v_r,
        \]
        que vuelven a ser combinaciones lineales de $v_1,\ldots,v_r$.

        Finalmente, cada $v_i$ es la combinación lineal con coeficiente $1$
        en $v_i$ y $0$ en los demás.
    \end{proof}
\end{frame}

\begin{frame}
    \frametitle{El menor subespacio que contiene a unos vectores}

    \begin{proposition}
        Si $W\subseteq V$ es un subespacio y $v_1,\ldots,v_r\in W$, entonces
        \[
            \langle v_1,\ldots,v_r\rangle\subseteq W.
        \]
    \end{proposition}

    \begin{proof}
        Todo elemento de $\langle v_1,\ldots,v_r\rangle$ es una combinación
        lineal de $v_1,\ldots,v_r$, y los subespacios son cerrados para
        combinaciones lineales.
    \end{proof}

    \medskip

    Así, $\langle v_1,\ldots,v_r\rangle$ es el \alert{menor subespacio}
    de $V$ que contiene a $v_1,\ldots,v_r$: los contiene, y está contenido
    en cualquier otro subespacio que los contenga.
\end{frame}

\begin{frame}
    \frametitle{Generadores de un subespacio}

    Como los subespacios son espacios vectoriales, tiene sentido decir que
    unos vectores \alert{generan} un subespacio.

    \begin{proposition}
        \begin{enumerate}
            \item $\langle v_1,\ldots,v_r\rangle$ está generado por
            $v_1,\ldots,v_r$.
            \item Si un subespacio $W\subseteq V$ está generado por
            $v_1,\ldots,v_r$, entonces $W=\langle v_1,\ldots,v_r\rangle$.
        \end{enumerate}
    \end{proposition}

    \begin{proof}
        (1) Los vectores $v_1,\ldots,v_r$ están en
        $\langle v_1,\ldots,v_r\rangle$, y todo elemento de
        $\langle v_1,\ldots,v_r\rangle$ es, por definición, combinación
        lineal de ellos.

        (2) Que $v_1,\ldots,v_r$ generen $W$ significa que están en $W$ y que
        todo vector de $W$ es combinación lineal de ellos, es decir,
        $W\subseteq\langle v_1,\ldots,v_r\rangle$. La otra inclusión se debe
        a que $\langle v_1,\ldots,v_r\rangle$ es el menor subespacio que
        contiene a $v_1,\ldots,v_r$.
    \end{proof}
\end{frame}

\begin{frame}
    \frametitle{Rectas y planos generados}

    \begin{example}
        Sean $u,w\in\R^3$.
        \begin{itemize}
            \item Si $u\neq0$, $\langle u\rangle=\{\lambda u:\lambda\in\R\}$
            es la \alert{recta} que pasa por el origen con dirección $u$.
            \item Si $u$ y $w$ no son proporcionales, $\langle u,w\rangle$
            es el \alert{plano} que pasa por el origen y contiene a $u$ y $w$.
            \item Si $w=\alpha u$, entonces
            $\lambda u+\mu w=(\lambda+\mu\alpha)u$, luego
            \[
                \langle u,w\rangle=\langle u\rangle
            \]
            es una recta, aunque haya dos generadores.
        \end{itemize}
    \end{example}

    \medskip

    El subespacio generado depende de cómo son los vectores, no solo de
    cuántos hay.
\end{frame}

\begin{frame}
    \frametitle{Generadores de un subespacio de polinomios}

    \begin{example}
        Sea $W=\{p(x)\in k_2[x]:\ p(1)=0\}$. Entonces
        \[
            W=\langle x-1,\ x^2-1\rangle.
        \]
    \end{example}

    \vspace{-15pt}

    \begin{proof}[Comprobación]
        \alert{$\supseteq$.} Los polinomios $x-1$ y $x^2-1$ se anulan en $1$,
        luego están en el subespacio $W$, y por tanto
        $\langle x-1,x^2-1\rangle\subseteq W$.

        \alert{$\subseteq$.} Sea $p(x)=a_0+a_1x+a_2x^2\in W$. Como
        $p(1)=a_0+a_1+a_2=0$, se tiene $a_0=-a_1-a_2$, y
        \[
            p(x)=-a_1-a_2+a_1x+a_2x^2=a_1(x-1)+a_2(x^2-1).
        \]
    \end{proof}

    Además, $x-1$ y $x^2-1$ son linealmente independientes ya que no son proporcionales. Por tanto, $\{x-1,x^2-1\}$ es una
    \alert{base} de $W$.
\end{frame}

\begin{frame}
    \frametitle{Generadores de un subespacio de matrices}

    \begin{example}
        El subespacio de las matrices simétricas de $M_2(k)$ es
        \[
            \Bigl\langle
            \begin{pmatrix}1&0\\0&0\end{pmatrix},
            \begin{pmatrix}0&0\\0&1\end{pmatrix},
            \begin{pmatrix}0&1\\1&0\end{pmatrix}
            \Bigr\rangle.
        \]
    \end{example}

    \medskip

    Las tres matrices son simétricas. Recíprocamente, una matriz es simétrica
    si y solo si es de la forma
    \[
        \begin{pmatrix}a&b\\b&c\end{pmatrix}
        =a\begin{pmatrix}1&0\\0&0\end{pmatrix}
        +c\begin{pmatrix}0&0\\0&1\end{pmatrix}
        +b\begin{pmatrix}0&1\\1&0\end{pmatrix}.
    \]
    Los coeficientes de esta combinación lineal son las entradas $a$, $b$,
    $c$ de la matriz, luego son únicos. Por el teorema de representación
    única, las tres matrices forman una \alert{base} de este subespacio.
\end{frame}

\begin{frame}
    \frametitle{Generadores de las soluciones de un sistema homogéneo}

    \begin{proposition}
        Sea $A\in M_{m\times n}(k)$ y sean $v_1,\ldots,v_s$ las soluciones
        de $Av=0$ obtenidas en el Tema~1 dando el valor $1$ a una incógnita
        libre y $0$ a las demás, $s=n-\rg(A)$. Entonces
        \[
            W=\{v\in k^n:\ Av=0\}=\langle v_1,\ldots,v_s\rangle,
        \]
        y $\{v_1,\ldots,v_s\}$ es una base de este subespacio, en particular
        $\dim(W)=n-\rg(A)$.
    \end{proposition}

    \vspace{-8pt}

    \begin{proof}
        En el Tema~1 vimos que toda solución es combinación lineal de
        $v_1,\ldots,v_s$; esto da la inclusión $\subseteq$. La inclusión
        $\supseteq$ se debe a que $v_1,\ldots,v_s$ son soluciones y el
        conjunto de soluciones es un subespacio. También vimos en el Tema~1
        que los coeficientes de esa combinación lineal son \alert{únicos}
        (son los valores de las incógnitas libres), así que, por el teorema
        de representación única, $\{v_1,\ldots,v_s\}$ es una base.
    \end{proof}

    \vspace{-5pt}

    Por tanto, resolver un sistema homogéneo es encontrar una
    \alert{base} de su conjunto de soluciones.
\end{frame}

\begin{frame}
    \frametitle{Sistemas homogéneos en coordenadas}

    \begin{corollary}
        Sea $B$ una base de un espacio vectorial $V$ de dimensión $n$,
        sea $A\in M_{m\times n}(k)$ y sea
        \[
            W=\{v\in V:\ Av_B=0\}\colon Ax=0\text{ resp. }B.
        \]
        Sea $\{v_1,\ldots,v_s\}$ la base de $\{x\in k^n:\ Ax=0\}$ de la
        proposición anterior, $s=n-\rg(A)$, y sean $w_1,\ldots,w_s\in V$ los vectores
        con $(w_i)_B=v_i$. Entonces $\{w_1,\ldots,w_s\}$ es una base
        de $W$, luego $\dim(W)=n-\rg(A)$.
    \end{corollary}

    \vspace{-8pt}

    \begin{proof}
        Los $w_i$ existen porque $c_B$ es biyectiva. Como $c_B$
        preserva y refleja las combinaciones lineales,
        $v\in W\Leftrightarrow v_B\in\langle v_1,\ldots,v_s\rangle
        \Leftrightarrow v\in\langle w_1,\ldots,w_s\rangle$,
        así que $W=\langle w_1,\ldots,w_s\rangle$. Además,
        $w_1,\ldots,w_s$ son linealmente independientes porque sus
        coordenadas $v_1,\ldots,v_s$ lo son (Tema~2).
    \end{proof}

    \vspace{-5pt}

    Es decir: se resuelve el sistema en $k^n$ y se \alert{traducen} las
    soluciones básicas a vectores de $V$ mediante la base $B$.
\end{frame}

\begin{frame}
    \frametitle{Ejemplo: sistemas homogéneos en coordenadas}

    \begin{example}
        Volvamos a $W=\{p(x)\in k_2[x]:\ p(1)=0\}$, con la base
        $B=\{1,x,x^2\}$. Si $p(x)=a_0+a_1x+a_2x^2$, entonces
        \[
            p(x)\in W\iff a_0+a_1+a_2=0
            \iff
            \begin{pmatrix}1&1&1\end{pmatrix}p(x)_B=0.
        \]
        Las incógnitas libres son $a_1$ y $a_2$, y las soluciones básicas
        son $v_1=(-1,1,0)$ y $v_2=(-1,0,1)$.
        Los polinomios con esas coordenadas son $w_1=x-1$ y $w_2=x^2-1$,
        así que $\{x-1,x^2-1\}$ es una base de $W$.
    \end{example}

    \medskip

    Recuperamos el resultado obtenido antes, ahora sin necesidad de
    \alert{adivinar} los generadores.
\end{frame}

\begin{frame}
    \frametitle{¿Pertenece un vector al subespacio generado?}

    \begin{proposition}
        Sean $v_1,\ldots,v_r,u\in k^n$ y sea
        $A=(v_1\mid\cdots\mid v_r)\in M_{n\times r}(k)$. Entonces
        \[
            u\in\langle v_1,\ldots,v_r\rangle
            \iff
            \rg(A)=\rg(A\mid u).
        \]
    \end{proposition}

    \begin{proof}
        $u\in\langle v_1,\ldots,v_r\rangle$ significa que existen
        $\lambda_1,\ldots,\lambda_r\in k$ con
        $\lambda_1v_1+\cdots+\lambda_rv_r=u$, es decir, que el sistema
        \[
            A\begin{pmatrix}\lambda_1\\\vdots\\\lambda_r\end{pmatrix}=u
        \]
        es compatible. Por el teorema de Rouché--Frobenius, esto equivale
        a $\rg(A)=\rg(A\mid u)$.
    \end{proof}
\end{frame}

\begin{frame}[allowframebreaks, t]
    \frametitle{Ejemplo: pertenencia al subespacio generado}

    \begin{example}
        ¿Para qué valores de $a\in\R$ se tiene
        $u=(2,3,a,1)\in\langle v_1,v_2\rangle$, con
        \[
            v_1=(1,0,1,2),
            \qquad
            v_2=(0,1,1,-1)?
        \]
    \end{example}

    Escalonamos la matriz ampliada $(v_1\mid v_2\mid u)$ con
    $F_3-F_1-F_2$ y $F_4-2F_1+F_2$:
    \[
        \left(\begin{array}{cc|c}
            1&0&2\\
            0&1&3\\
            1&1&a\\
            2&-1&1
        \end{array}\right)
        \longrightarrow
        \left(\begin{array}{cc|c}
            1&0&2\\
            0&1&3\\
            0&0&a-5\\
            0&0&0
        \end{array}\right).
    \]

    \framebreak

    La matriz de coeficientes tiene rango $2$. La ampliada tiene rango $2$
    si $a=5$ y rango $3$ si $a\neq5$. Por tanto,
    \[
        u\in\langle v_1,v_2\rangle
        \iff
        a=5.
    \]

    \medskip

    En ese caso, la forma escalonada da directamente los coeficientes:
    $\lambda_1=2$, $\lambda_2=3$, es decir,
    \[
        (2,3,5,1)=2v_1+3v_2.
    \]
\end{frame}

\begin{frame}
    \frametitle{Comparar subespacios generados}

    \begin{proposition}
        Sean $v_1,\ldots,v_r,w_1,\ldots,w_s\in V$. Entonces
        \[
            \langle v_1,\ldots,v_r\rangle\subseteq\langle w_1,\ldots,w_s\rangle
            \iff
            v_1,\ldots,v_r\in\langle w_1,\ldots,w_s\rangle.
        \]
    \end{proposition}

    \begin{proof}
        \alert{$\Rightarrow$.} Cada $v_i$ está en
        $\langle v_1,\ldots,v_r\rangle$, luego en
        $\langle w_1,\ldots,w_s\rangle$.

        \alert{$\Leftarrow$.} $\langle w_1,\ldots,w_s\rangle$ es un subespacio
        que contiene a $v_1,\ldots,v_r$, luego contiene al menor subespacio
        que los contiene, que es $\langle v_1,\ldots,v_r\rangle$.
    \end{proof}

    \medskip

    Para probar que dos subespacios generados son \alert{iguales}, basta
    ver que cada generador de uno pertenece al otro.
\end{frame}

\begin{frame}
    \frametitle{Ejemplo: dos sistemas de generadores del mismo subespacio}

    \begin{example}
        En $\R^3$,
        \[
            \langle(1,1,0),(0,1,1)\rangle=\langle(1,2,1),(1,0,-1)\rangle.
        \]
    \end{example}

    \medskip

    Llamemos $v_1=(1,1,0)$, $v_2=(0,1,1)$, $w_1=(1,2,1)$, $w_2=(1,0,-1)$.

    \medskip

    \alert{$\supseteq$.} $w_1=v_1+v_2$ y $w_2=v_1-v_2$ están en
    $\langle v_1,v_2\rangle$.

    \medskip

    \alert{$\subseteq$.} Despejando,
    \[
        v_1=\tfrac12w_1+\tfrac12w_2,
        \qquad
        v_2=\tfrac12w_1-\tfrac12w_2
    \]
    están en $\langle w_1,w_2\rangle$.

    \medskip

    Un mismo subespacio tiene, en general, \alert{muchos} sistemas de
    generadores distintos.
\end{frame}

\begin{frame}
    \frametitle{Generadores redundantes}

    \begin{corollary}
        Sean $v_1,\ldots,v_r,u\in V$. Entonces
        \[
            \langle v_1,\ldots,v_r,u\rangle=\langle v_1,\ldots,v_r\rangle
            \iff
            u\in\langle v_1,\ldots,v_r\rangle.
        \]
    \end{corollary}

    \begin{proof}
        \alert{$\Rightarrow$.} El vector $u$ está en
        $\langle v_1,\ldots,v_r,u\rangle$, que por hipótesis es
        $\langle v_1,\ldots,v_r\rangle$.

        \alert{$\Leftarrow$.} La inclusión $\supseteq$ se cumple siempre,
        porque $v_1,\ldots,v_r\in\langle v_1,\ldots,v_r,u\rangle$. Para la
        inclusión $\subseteq$, los vectores $v_1,\ldots,v_r$ están en
        $\langle v_1,\ldots,v_r\rangle$, y $u$ también, por hipótesis.
        En ambos casos aplicamos la proposición anterior.
    \end{proof}

    \medskip

    Un generador que es combinación lineal de los demás se puede
    \alert{eliminar} sin cambiar el subespacio generado.
\end{frame}

\begin{frame}[allowframebreaks, t]
    \frametitle{De generadores a ecuaciones implícitas}

    \begin{proposition}
        Sean $v_1,\ldots,v_r\in k^n$ y $A=(v_1\mid\cdots\mid v_r)$.
        Las ecuaciones implícitas de $W=\langle v_1,\ldots,v_r\rangle$
        se obtienen escalonando la matriz $(A\mid u)$, donde
        $u=(u_1,\ldots,u_n)$ es simbólico: cada fila en que la parte
        de coeficientes se anule da una condición lineal
        $a_1u_1+\cdots+a_nu_n=0$, y $W$ es exactamente el conjunto de
        vectores que satisfacen todas esas condiciones.
    \end{proposition}

    \begin{proof}
        Un vector $u\in k^n$ pertenece a $W$ si y solo si el sistema
        $Ax=u$ es compatible. Al escalonar $(A\mid u)$ con $u$ simbólico,
        las filas en que la parte de coeficientes se anula expresan
        exactamente las condiciones de compatibilidad sobre $u_1,\ldots,u_n$.
    \end{proof}

    En particular, si $\dim W=s$, hay $n-s$ ecuaciones implícitas
    independientes.
\end{frame}


\begin{frame}[allowframebreaks, t]
    \frametitle{Ejemplo: ecuaciones implícitas del plano}

    \begin{example}
        Hallemos las ecuaciones implícitas de
        $W=\langle(2,1,0),(-1,0,1)\rangle\subseteq\R^3$.
    \end{example}

    Escalonamos la matriz ampliada $(A\mid u)$:
    \[
        \left(\begin{array}{cc|c}
            2&-1&u_1\\1&0&u_2\\0&1&u_3
        \end{array}\right)
        \xrightarrow{\;F_1\leftrightarrow F_2\;}
        \left(\begin{array}{cc|c}
            1&0&u_2\\2&-1&u_1\\0&1&u_3
        \end{array}\right)
        \xrightarrow{\;F_2-2F_1\;}
        \left(\begin{array}{cc|c}
            1&0&u_2\\0&-1&u_1-2u_2\\0&1&u_3
        \end{array}\right)
    \]
    \[
        \xrightarrow{\;F_3+F_2\;}
        \left(\begin{array}{cc|c}
            1&0&u_2\\0&-1&u_1-2u_2\\0&0&\alert{u_1-2u_2+u_3}
        \end{array}\right).
    \]
    La tercera fila da la condición $u_1-2u_2+u_3=0$. Por tanto,
    \[
        W\colon u_1-2u_2+u_3=0,
    \]
    que coincide con la ecuación del plano del ejemplo anterior.
\end{frame}

\subsection{Dimensión de un subespacio}

\begin{frame}[allowframebreaks]
    \frametitle{La dimensión de un subespacio}

    \begin{theorem}
        Sea $V$ un espacio vectorial de dimensión $n$ y sea $W\subseteq V$
        un subespacio. Entonces $W$ tiene dimensión finita y
        \[
            \dim(W)\leq n.
        \]
    \end{theorem}

    \medskip

    No es evidente: a priori no sabemos que $W$ tenga un sistema generador
    finito. Hay que \alert{construir} una base de $W$.

    \framebreak

    \begin{proof}
        Si $W=\{0\}$, entonces $\dim(W)=0$. Supongamos $W\neq\{0\}$.

        \medskip

        Unos vectores de $W$ son linealmente independientes en $W$ si y solo
        si lo son en $V$, pues la condición es la misma. Por el Tema~2, una
        familia independiente de $V$ tiene a lo sumo $n$ vectores.

        \medskip

        Como $W$ contiene algún vector no nulo, que por sí solo es
        independiente, podemos elegir $w_1,\ldots,w_m\in W$ linealmente
        independientes con $m\geq1$ \alert{lo mayor posible}; existe porque
        $m\leq n$.

        \medskip

        Veamos que $w_1,\ldots,w_m$ generan $W$. Si algún $u\in W$ no fuera
        combinación lineal de ellos, por el lema de añadir un vector a una
        familia independiente, $w_1,\ldots,w_m,u$ serían $m+1$ vectores
        independientes de $W$, en contra de la elección de $m$.

        \medskip

        Por tanto, $\{w_1,\ldots,w_m\}$ es una base de $W$ y
        $\dim(W)=m\leq n$.
    \end{proof}
\end{frame}

\begin{frame}
    \frametitle{Subespacios de la misma dimensión}

    \begin{corollary}
        Sea $V$ un espacio vectorial de dimensión finita y sean
        $W\subseteq W'$ subespacios de $V$. Entonces:
        \begin{enumerate}
            \item $\dim(W)\leq\dim(W')$;
            \item si $\dim(W)=\dim(W')$, entonces $W=W'$.
        \end{enumerate}
    \end{corollary}

    \begin{proof}
        Por el teorema, $W'$ es un espacio vectorial de dimensión finita, y
        $W$ es un subespacio suyo; aplicando de nuevo el teorema se obtiene
        (1).

        Para (2), sea $\{w_1,\ldots,w_m\}$ una base de $W$. Son
        $m=\dim(W')$ vectores linealmente independientes de $W'$, luego
        forman una base de $W'$ (Tema~2). Por tanto,
        $W'=\langle w_1,\ldots,w_m\rangle=W$.
    \end{proof}

    \medskip

    En particular, si $W\subsetneq W'$, entonces $\dim(W)<\dim(W')$; y el
    único subespacio de $V$ de dimensión $\dim(V)$ es el propio $V$.
\end{frame}

\begin{frame}
    \frametitle{Igualdad de subespacios mediante la dimensión}

    \begin{example}
        Volvamos a la igualdad
        \[
            \langle(1,1,0),(0,1,1)\rangle=\langle(1,2,1),(1,0,-1)\rangle
        \]
        en $\R^3$, con $v_1=(1,1,0)$, $v_2=(0,1,1)$, $w_1=(1,2,1)$,
        $w_2=(1,0,-1)$.
    \end{example}

    \medskip

    Antes comprobamos las dos inclusiones. Ahora basta \alert{una}:
    \begin{itemize}
        \item $w_1=v_1+v_2$ y $w_2=v_1-v_2$, luego
        $\langle w_1,w_2\rangle\subseteq\langle v_1,v_2\rangle$;
        \item ambos subespacios tienen dimensión $2$, pues cada uno está
        generado por dos vectores no proporcionales, que son por tanto
        una base.
    \end{itemize}
    Por el corolario, los dos subespacios coinciden.

    \medskip

    Para probar $W=W'$ basta ver $W\subseteq W'$ y $\dim(W)=\dim(W')$.
\end{frame}

\begin{frame}
    \frametitle{Los subespacios de $\R^3$}

    \begin{example}
        Los subespacios de $\R^3$ son exactamente: $\{0\}$, las rectas que
        pasan por el origen, los planos que pasan por el origen y $\R^3$.
    \end{example}

    \medskip

    Sea $W\subseteq\R^3$ un subespacio. Por el teorema,
    $\dim(W)\in\{0,1,2,3\}$:
    \begin{itemize}
        \item si $\dim(W)=0$, entonces $W=\{0\}$;
        \item si $\dim(W)=1$, entonces $W=\langle u\rangle$ con $u\neq0$, una
        \alert{recta} que pasa por el origen;
        \item si $\dim(W)=2$, entonces $W=\langle u,w\rangle$ con $u,w$
        independientes, es decir, no proporcionales: un \alert{plano} que
        pasa por el origen;
        \item si $\dim(W)=3$, entonces $W=\R^3$ por el corolario.
    \end{itemize}

    \medskip

    Recíprocamente, ya sabemos que todos estos conjuntos son subespacios.
\end{frame}

\begin{frame}
    \frametitle{En dimensión infinita}

    \begin{example}
        En $k[x]$, los subespacios
        \[
            k_0[x]\subsetneq k_1[x]\subsetneq k_2[x]\subsetneq\cdots
        \]
        forman una cadena estrictamente creciente que no termina.
    \end{example}

    \medskip

    En un espacio $V$ de dimensión $n$ esto es imposible: si
    $W_0\subsetneq W_1\subsetneq\cdots\subsetneq W_m$ son subespacios,
    el corolario da
    \[
        0\leq\dim(W_0)<\dim(W_1)<\cdots<\dim(W_m)\leq n,
    \]
    luego $m\leq n$.

    \medskip

    Además, $W=\{p(x)\in k[x]:\ p(0)=0\}$ es un subespacio
    \alert{propio} de $k[x]$, pues $1\notin W$, que tampoco tiene dimensión
    finita: si tuviera dimensión $d$, los $d+1$ polinomios
    $x,x^2,\ldots,x^{d+1}\in W$ serían independientes, en contra del Tema~2.
\end{frame}

\begin{frame}[allowframebreaks]
    \frametitle{Una base del subespacio generado}

    \begin{proposition}
        Sean $v_1,\ldots,v_r\in k^n$, sea
        $A=(v_1\mid\cdots\mid v_r)\in M_{n\times r}(k)$ y sea $E$ una forma
        escalonada de $A$. Los $v_j$ tales que la columna $j$ de $E$ tiene
        pivote forman una base de $\langle v_1,\ldots,v_r\rangle$. En
        particular,
        \[
            \dim\langle v_1,\ldots,v_r\rangle=\rg(A).
        \]
    \end{proposition}

    \medskip

    Es el método de existencia de bases en $k^n$ del Tema~2, sin suponer
    que los $v_j$ generan $k^n$.

    \medskip

    En particular, el rango de una matriz es la \alert{dimensión} del
    subespacio generado por sus columnas.

    \framebreak

    \begin{proof}
        Sean $j_1<\cdots<j_m$ las columnas de $E$ con pivote, $m=\rg(A)$.
        En el Tema~2 vimos que:
        \begin{itemize}
            \item $v_{j_1},\ldots,v_{j_m}$ son linealmente independientes,
            pues las mismas operaciones elementales transforman
            $(v_{j_1}\mid\cdots\mid v_{j_m})$ en una matriz escalonada con
            $m$ pivotes;
            \item si la columna $j$ de $E$ no tiene pivote, entonces $v_j$ es
            combinación lineal de $v_1,\ldots,v_{j-1}$.
        \end{itemize}

        Eliminamos uno a uno los $v_j$ cuya columna no tiene pivote,
        empezando por el de mayor índice. Cuando eliminamos $v_j$, los
        vectores $v_1,\ldots,v_{j-1}$ siguen presentes, así que, por el
        corolario de los generadores redundantes, el subespacio generado
        no cambia. Al final,
        \[
            \langle v_{j_1},\ldots,v_{j_m}\rangle
            =\langle v_1,\ldots,v_r\rangle,
        \]
        luego $\{v_{j_1},\ldots,v_{j_m}\}$ es una base y la dimensión es
        $m=\rg(A)$.
    \end{proof}
\end{frame}

\begin{frame}
    \frametitle{Una base del subespacio generado en coordenadas}

    \begin{corollary}
        Sea $B$ una base de un espacio vectorial $V$ de dimensión $n$, sean
        $u_1,\ldots,u_r\in V$ y sea
        $A=\bigl((u_1)_B\mid\cdots\mid(u_r)_B\bigr)$. Los $u_j$ tales que
        la columna $j$ de una forma escalonada de $A$ tiene pivote forman
        una base de $\langle u_1,\ldots,u_r\rangle$, y
        \[
            \dim\langle u_1,\ldots,u_r\rangle=\rg(A).
        \]
    \end{corollary}

    \begin{proof}
        Sean $j_1<\cdots<j_m$ las columnas con pivote. Como $c_B$ preserva
        y refleja las combinaciones lineales, y por la proposición anterior,
        \[
            u\in\langle u_1,\ldots,u_r\rangle
            \iff u_B\in\langle(u_{j_1})_B,\ldots,(u_{j_m})_B\rangle
            \iff u\in\langle u_{j_1},\ldots,u_{j_m}\rangle.
        \]
        Además, $u_{j_1},\ldots,u_{j_m}$ son linealmente independientes
        porque sus coordenadas lo son (Tema~2).
    \end{proof}
\end{frame}

\begin{frame}[allowframebreaks, t]
    \frametitle{Ejemplo: dimensión de un subespacio de $\R^4$}

    \begin{example}
        Calculemos la dimensión y una base de
        $W=\langle v_1,v_2,v_3,v_4,v_5\rangle\subseteq\R^4$, donde
        \[
            \begin{aligned}
                v_1&=(1,2,0,1), & v_2&=(2,4,1,3), & v_3&=(0,0,1,1),\\
                v_4&=(1,2,1,2), & v_5&=(0,1,0,1).
            \end{aligned}
        \]
    \end{example}

    Escalonamos la matriz que tiene por columnas $v_1,\ldots,v_5$:
    \[
        \begin{pmatrix}
            1&2&0&1&0\\
            2&4&0&2&1\\
            0&1&1&1&0\\
            1&3&1&2&1
        \end{pmatrix}
        \xrightarrow[\;F_4-F_1\;]{\;F_2-2F_1\;}
        \begin{pmatrix}
            1&2&0&1&0\\
            0&0&0&0&1\\
            0&1&1&1&0\\
            0&1&1&1&1
        \end{pmatrix}
        \xrightarrow[\;F_4-F_2,\ F_4-F_3\;]{\;F_2\leftrightarrow F_3\;}
        \begin{pmatrix}
            \alert{1}&2&0&1&0\\
            0&\alert{1}&1&1&0\\
            0&0&0&0&\alert{1}\\
            0&0&0&0&0
        \end{pmatrix}.
    \]

    \framebreak

    Los pivotes están en las columnas $1$, $2$ y $5$. Por tanto,
    \[
        \{v_1,v_2,v_5\}\ \text{es una base de }W,
        \qquad
        \dim(W)=3.
    \]

    \medskip

    En particular, $W\neq\R^4$: los cinco vectores \alert{no} generan
    $\R^4$.

    \medskip

    Los vectores eliminados son combinación lineal de los anteriores, como
    asegura la proposición:
    \[
        v_3=-2v_1+v_2,
        \qquad
        v_4=-v_1+v_2.
    \]
\end{frame}

\begin{frame}[allowframebreaks, t]
    \frametitle{Ejemplo: dimensión de un subespacio de polinomios}

    \begin{example}
        Sea $W=\langle 1+x,\ x+x^2,\ 1-x^2,\ 1+2x+x^2\rangle\subseteq\R_2[x]$.
    \end{example}

    Escalonamos la matriz de coordenadas respecto de $\{1,x,x^2\}$:
    \[
        \begin{pmatrix}
            1&0&1&1\\
            1&1&0&2\\
            0&1&-1&1
        \end{pmatrix}
        \xrightarrow{\;F_2-F_1\;}
        \begin{pmatrix}
            1&0&1&1\\
            0&1&-1&1\\
            0&1&-1&1
        \end{pmatrix}
        \xrightarrow{\;F_3-F_2\;}
        \begin{pmatrix}
            \alert{1}&0&1&1\\
            0&\alert{1}&-1&1\\
            0&0&0&0
        \end{pmatrix}.
    \]
    Los pivotes están en las columnas $1$ y $2$: $\{1+x,\ x+x^2\}$ es una
    base de $W$ y $\dim(W)=2$.

    \framebreak

    Podemos describir $W$ de otra forma. Sea
    \[
        W'=\{p(x)\in\R_2[x]:\ p(-1)=0\}.
    \]
    \begin{itemize}
        \item $W\subseteq W'$, porque los cuatro generadores de $W$ se anulan
        en $-1$.
        \item Si $p(x)=a_0+a_1x+a_2x^2$, entonces $p(-1)=a_0-a_1+a_2$, luego
        $W'\colon\begin{pmatrix}1&-1&1\end{pmatrix}x=0$ resp.\
        $\{1,x,x^2\}$, y $\dim(W')=3-1=2$.
    \end{itemize}

    \medskip

    Como $W\subseteq W'$ y tienen la misma dimensión, $W=W'$: los
    polinomios de $W$ son exactamente los de grado $\leq2$ que se anulan
    en $-1$.
\end{frame}

\begin{frame}[allowframebreaks, t]
    \frametitle{Ejemplo: dimensión de un subespacio de matrices}

    \begin{example}
        ¿Cuál es la dimensión de $W=\langle A_1,A_2,A_3\rangle\subseteq M_2(\R)$,
        según el valor de $t\in\R$?
        \[
            A_1=\begin{pmatrix}1&1\\0&1\end{pmatrix},
            \qquad
            A_2=\begin{pmatrix}0&1\\1&0\end{pmatrix},
            \qquad
            A_3=\begin{pmatrix}2&1\\-1&t\end{pmatrix}.
        \]
    \end{example}

    Usamos la base de $M_2(\R)$ en la que las coordenadas de
    $\begin{pmatrix}a&b\\c&d\end{pmatrix}$ son $(a,b,c,d)$:
    \[
        \begin{pmatrix}
            1&0&2\\
            1&1&1\\
            0&1&-1\\
            1&0&t
        \end{pmatrix}
        \xrightarrow[\;F_4-F_1\;]{\;F_2-F_1\;}
        \begin{pmatrix}
            1&0&2\\
            0&1&-1\\
            0&1&-1\\
            0&0&t-2
        \end{pmatrix}
        \xrightarrow[\;F_3\leftrightarrow F_4\;]{\;F_3-F_2\;}
        \begin{pmatrix}
            \alert{1}&0&2\\
            0&\alert{1}&-1\\
            0&0&t-2\\
            0&0&0
        \end{pmatrix}.
    \]

    \framebreak

    \begin{itemize}
        \item Si $t\neq2$, hay pivote en las tres columnas: $\{A_1,A_2,A_3\}$
        es una base de $W$ y $\dim(W)=3$.
        \item Si $t=2$, solo hay pivote en las columnas $1$ y $2$:
        $\{A_1,A_2\}$ es una base de $W$ y $\dim(W)=2$. En efecto,
        \[
            A_3=\begin{pmatrix}2&1\\-1&2\end{pmatrix}=2A_1-A_2.
        \]
    \end{itemize}

    \medskip

    En ningún caso $W=M_2(\R)$, que tiene dimensión $4$: tres matrices
    nunca bastan para generarlo.
\end{frame}

\begin{frame}[allowframebreaks]
    \frametitle{Bases adaptadas a un subespacio}

    \begin{proposition}
        Sea $V$ un espacio vectorial de dimensión $n$ y sea $W\subseteq V$
        un subespacio de dimensión $m$. Toda base $\{w_1,\ldots,w_m\}$ de
        $W$ se amplía a una base $B=\{w_1,\ldots,w_n\}$ de $V$, y entonces
        \[
            W\colon\ x_{m+1}=\cdots=x_n=0\quad\text{resp.\ }B.
        \]
    \end{proposition}

    \medskip

    Diremos que $B$ es una base de $V$ \alert{adaptada} a $W$.

    \framebreak

    \begin{proof}
        Los vectores $w_1,\ldots,w_m$ son linealmente independientes en
        $V$, así que, por el Tema~2, se pueden ampliar a una base
        $B=\{w_1,\ldots,w_n\}$ de $V$.

        \medskip

        Sea $u\in V$ con $u_B=(\lambda_1,\ldots,\lambda_n)$, es decir,
        $u=\lambda_1w_1+\cdots+\lambda_nw_n$.

        \medskip

        Si $\lambda_{m+1}=\cdots=\lambda_n=0$, entonces $u$ es combinación
        lineal de $w_1,\ldots,w_m$, luego $u\in W$.

        \medskip

        Recíprocamente, si $u\in W$, existen $\mu_1,\ldots,\mu_m\in k$ con
        \[
            u=\mu_1w_1+\cdots+\mu_mw_m+0w_{m+1}+\cdots+0w_n.
        \]
        Por la unicidad de las coordenadas, $\lambda_i=\mu_i$ para
        $i\leq m$ y $\lambda_{m+1}=\cdots=\lambda_n=0$.
    \end{proof}

    \medskip

    Es lo que ocurría en el ejemplo del plano
    $x_1-2x_2+x_3=0$ de $\R^3$.
\end{frame}

\begin{frame}
    \frametitle{Ejemplo: una base adaptada en polinomios}

    \begin{example}
        Ampliemos la base $\{1+x,\ x+x^2\}$ de
        $W=\{p(x)\in\R_2[x]:\ p(-1)=0\}$ a una base de $\R_2[x]$.
    \end{example}

    Como en el Tema~2, añadimos la base $\{1,x,x^2\}$ y aplicamos el método
    de existencia de bases en coordenadas:
    \[
        \begin{pmatrix}
            1&0&1&0&0\\
            1&1&0&1&0\\
            0&1&0&0&1
        \end{pmatrix}
        \xrightarrow{\;F_2-F_1\;}
        \begin{pmatrix}
            1&0&1&0&0\\
            0&1&-1&1&0\\
            0&1&0&0&1
        \end{pmatrix}
        \xrightarrow{\;F_3-F_2\;}
        \begin{pmatrix}
            \alert{1}&0&1&0&0\\
            0&\alert{1}&-1&1&0\\
            0&0&\alert{1}&-1&1
        \end{pmatrix}.
    \]
    Los pivotes están en las columnas $1$, $2$ y $3$, así que
    $B=\{1+x,\ x+x^2,\ 1\}$ es una base de $\R_2[x]$ adaptada a $W$, y
    \[
        W\colon\ x_3=0\quad\text{resp.\ }B.
    \]
    En efecto, si $p(x)=a(1+x)+b(x+x^2)+c$, entonces $p(-1)=c$.
\end{frame}

\subsection{Operaciones con subespacios}

\begin{frame}[allowframebreaks, t]
    \frametitle{Intersección de subespacios}

    \begin{proposition}
        Si $W_1$ y $W_2$ son subespacios de $V$, entonces
        $W_1\cap W_2$ también lo es.
    \end{proposition}

    \begin{proof}
        Como $0\in W_1$ y $0\in W_2$, se tiene $0\in W_1\cap W_2$.
        Sean $u,v\in W_1\cap W_2$ y $\lambda\in k$.
        Como $W_1$ es un subespacio, $u+v\in W_1$ y $\lambda u\in W_1$;
        análogamente para $W_2$.
        Por tanto $u+v\in W_1\cap W_2$ y $\lambda u\in W_1\cap W_2$.
    \end{proof}
\end{frame}


\begin{frame}
    \frametitle{Ecuaciones de la intersección en $k^n$}

    \begin{proposition}
        Sean $A_1\in M_{m_1\times n}(k)$ y $A_2\in M_{m_2\times n}(k)$. Si
        $W_1\colon A_1x=0$ y $W_2\colon A_2x=0$ en $k^n$,
        entonces
        \[
            W_1\cap W_2\colon\
            \begin{pmatrix}A_1\\ \hline A_2\end{pmatrix}x=0.
        \]
        Es decir, la unión de ecuaciones implícitas de $W_1$ y $W_2$
        es un sistema de ecuaciones implícitas de $W_1\cap W_2$.
    \end{proposition}

    \begin{proof}
        Las primeras $m_1$ entradas de
        $\left(\begin{smallmatrix}A_1\\ \hline A_2\end{smallmatrix}\right)v$ son las de $A_1v$, y las
        $m_2$ restantes son las de $A_2v$. Por tanto, este vector es nulo
        si y solo si $A_1v=0$ y $A_2v=0$, es decir, si y solo si
        $v\in W_1$ y $v\in W_2$.
    \end{proof}
\end{frame}


\begin{frame}[allowframebreaks]
    \frametitle{Ejemplo: intersección de dos planos en $\R^3$}

    \begin{example}
        Sean $W_1\colon v_1-2v_2+v_3=0$ y $W_2\colon v_1+v_2=0$ dos
        planos de $\R^3$. Calculemos $W_1\cap W_2$.
    \end{example}

    Por la proposición anterior, $W_1\cap W_2$ está definido por las dos
    ecuaciones a la vez. Escalonamos la matriz de coeficientes:
    \[
        \begin{pmatrix}1&-2&1\\1&1&0\end{pmatrix}
        \xrightarrow{F_2-F_1}
        \begin{pmatrix}1&-2&1\\0&3&-1\end{pmatrix}.
    \]
    El rango es $2$, así que las ecuaciones son independientes y
    $\dim(W_1\cap W_2)=3-2=1$. La incógnita libre es $v_3$. Tomando
    $v_3=1$, la segunda fila da $3v_2-1=0$, es decir, $v_2=\frac13$, y la
    primera da $v_1=2v_2-v_3=-\frac13$. Por tanto,
    \[
        W_1\cap W_2=\left\langle\left(-\tfrac13,\tfrac13,1\right)\right\rangle
        =\langle(-1,1,3)\rangle,
    \]
    una recta que pasa por el origen.

    \framebreak

    \begin{columns}[c]
        \begin{column}{0.42\textwidth}
            Geométricamente, $W_1$ y $W_2$ son dos \alert{planos} distintos
            que pasan por el origen. Su intersección es la \alert{recta}
            \[
                W_1\cap W_2=\langle(-1,1,3)\rangle.
            \]
            El plano $W_2\colon v_1+v_2=0$ es vertical: contiene al eje
            $v_3$.
        \end{column}
        \begin{column}{0.54\textwidth}
            \centering
            \begin{tikzpicture}[x={(1cm,0cm)},y={(0.4cm,0.3cm)},
                    z={(0cm,1cm)},scale=0.8]
                % W_1 = < (1,0,-1), (-1,1,3) >
                \filldraw[gray,fill opacity=0.2,draw opacity=0.6]
                    (2,-0.6,-3.2) -- (0.8,0.6,0.4) -- (-2,0.6,3.2)
                    -- (-0.8,-0.6,-0.4) -- cycle;
                % W_2 = < (1,-1,0), (-1,1,3) >
                \filldraw[secondarycolor,fill opacity=0.35,draw opacity=0.8]
                    (2.2,-2.2,-1.8) -- (1,-1,1.8) -- (-2.2,2.2,1.8)
                    -- (-1,1,-1.8) -- cycle;
                \draw[->] (0,0,0) -- (2.6,0,0) node[right] {\footnotesize $v_1$};
                \draw[->] (0,0,0) -- (0,3,0) node[right] {\footnotesize $v_2$};
                \draw[->] (0,0,0) -- (0,0,2.6) node[above] {\footnotesize $v_3$};
                \draw[very thick,alertcolor] (0.6,-0.6,-1.8) -- (-0.6,0.6,1.8);
                \fill (0,0,0) circle (1.5pt);
                \node[gray] at (-2.4,0.6,3.4) {\footnotesize $W_1$};
                \node[secondarycolor!80!black] at (2.6,-2.2,-1.4)
                    {\footnotesize $W_2$};
                \node[alertcolor,left] at (-0.6,0.6,1.8)
                    {\footnotesize $W_1\cap W_2$};
            \end{tikzpicture}
        \end{column}
    \end{columns}

    \medskip

    En general, la unión de ecuaciones independientes de $W_1$ y de $W_2$
    \alert{no} tiene por qué ser independiente; si no lo es, se escalona
    para eliminar las redundantes.
\end{frame}


\begin{frame}
    \frametitle{Ecuaciones de la intersección respecto de una base}

    \begin{proposition}
        Sea $V$ un espacio vectorial de dimensión $n$ y sea $B$ una base
        de $V$. Si $W_1\colon A_1x=0$ resp.\ $B$ y
        $W_2\colon A_2x=0$ resp.\ $B$, entonces
        \[
            W_1\cap W_2\colon\
            \begin{pmatrix}A_1\\ \hline A_2\end{pmatrix}x=0\quad\text{resp.\ }B.
        \]
    \end{proposition}

    \begin{proof}
        Para $u\in V$, aplicamos la proposición anterior a $v=u_B$:
        \[
            u\in W_1\cap W_2
            \iff A_1u_B=0\ \text{y}\ A_2u_B=0
            \iff \begin{pmatrix}A_1\\ \hline A_2\end{pmatrix}u_B=0.\qedhere
        \]
    \end{proof}

    Las ecuaciones de $W_1$ y $W_2$ deben estar dadas respecto de la
    \alert{misma base}; si no, antes hay que cambiar de base una de ellas.
\end{frame}


\begin{frame}
    \frametitle{Ejemplo: intersección de subespacios de matrices}

    \begin{example}
        En $M_2(k)$, sean $W_1$ el subespacio de las matrices simétricas
        y $W_2$ el de las triangulares superiores. Calculemos
        $W_1\cap W_2$.
    \end{example}

    Respecto de la base
    \[
        B=\left\{
        \begin{pmatrix}1&0\\0&0\end{pmatrix},
        \begin{pmatrix}0&1\\0&0\end{pmatrix},
        \begin{pmatrix}0&0\\1&0\end{pmatrix},
        \begin{pmatrix}0&0\\0&1\end{pmatrix}
        \right\}
    \]
    ya vimos que $W_1\colon x_2-x_3=0$ y $W_2\colon x_3=0$. Por tanto,
    \[
        W_1\cap W_2\colon\
        \begin{pmatrix}0&1&-1&0\\0&0&1&0\end{pmatrix}x=0
        \quad\text{resp.\ }B,
    \]
    es decir, $x_2=x_3=0$: la intersección está formada por las matrices
    \alert{diagonales}
    $\begin{pmatrix}a&0\\0&d\end{pmatrix}$.
\end{frame}


\begin{frame}
    \frametitle{La unión de subespacios no es un subespacio}

    \begin{columns}[T]
        \begin{column}{0.62\textwidth}
            A diferencia de la intersección, la unión $W_1\cup W_2$
            no es en general un subespacio.

            \medskip

            Ya vimos que la unión de los dos ejes de $\R^2$ no lo es:
            $(1,0)$ y $(0,1)$ están en ella, pero
            \[
                (1,0)+(0,1)=(1,1)
            \]
            no está en ninguno de los dos ejes.
        \end{column}
        \begin{column}{0.34\textwidth}
            \centering
            \begin{tikzpicture}[scale=0.85]
                \draw[very thick,alertcolor] (-1.6,0) -- (2,0);
                \draw[very thick,alertcolor] (0,-1.2) -- (0,2);
                \node[alertcolor,right] at (2,0) {$W_1$};
                \node[alertcolor,above] at (0,2) {$W_2$};
                \draw[dashed] (1,0) -- (1,1) -- (0,1);
                \fill[alertcolor] (1,0) circle (2.5pt)
                    node[below] {\footnotesize $(1,0)$};
                \fill[alertcolor] (0,1) circle (2.5pt)
                    node[left] {\footnotesize $(0,1)$};
                \fill[white] (1,1) circle (2.5pt);
                \draw[thick] (1,1) circle (2.5pt)
                    node[above right] {\footnotesize $(1,1)$};
            \end{tikzpicture}
        \end{column}
    \end{columns}

    \medskip

    El menor subespacio que contiene a $W_1$ y a $W_2$ es la
    \alert{suma}, que definimos a continuación.
\end{frame}


\begin{frame}[allowframebreaks, t]
    \frametitle{Suma de subespacios}

    \begin{definition}
        La \alert{suma} de dos subespacios $W_1$ y $W_2$ de $V$ es
        \[
            W_1+W_2=\{u+w:\; u\in W_1,\; w\in W_2\}.
        \]
    \end{definition}

    \begin{proposition}
        $W_1+W_2$ es el menor subespacio de $V$ que contiene a $W_1$
        y a $W_2$.
    \end{proposition}

    \begin{proof}
        Es un subespacio: $0=0+0\in W_1+W_2$; si $u_1+w_1,u_2+w_2\in
        W_1+W_2$ y $\lambda\in k$, entonces
        \[
            (u_1+w_1)+(u_2+w_2)=(u_1+u_2)+(w_1+w_2)\in W_1+W_2,
        \]
        \[
            \lambda(u_1+w_1)=\lambda u_1+\lambda w_1\in W_1+W_2.
        \]
        Contiene a $W_1$ (tomando $w=0$) y a $W_2$ (tomando $u=0$).
        Si $W$ es cualquier subespacio que contiene a $W_1$ y $W_2$,
        entonces $u+w\in W$ para todo $u\in W_1$, $w\in W_2$,
        luego $W_1+W_2\subseteq W$.
    \end{proof}
\end{frame}


\begin{frame}[allowframebreaks, t]
    \frametitle{Generadores de la suma}

    \begin{proposition}
        Si $W_1=\langle v_1,\ldots,v_p\rangle$ y
        $W_2=\langle w_1,\ldots,w_q\rangle$, entonces
        \[
            W_1+W_2=\langle v_1,\ldots,v_p,w_1,\ldots,w_q\rangle.
        \]
        Es decir, la unión de sistemas de generadores de $W_1$ y $W_2$
        es un sistema de generadores de $W_1+W_2$.
    \end{proposition}

    \begin{proof}
        Probamos las dos inclusiones.

        ($\supseteq$) Los vectores $v_1,\ldots,v_p$ están en $W_1$ y
        $w_1,\ldots,w_q$ están en $W_2$, luego todos están en $W_1+W_2$.
        Como $W_1+W_2$ es un subespacio y
        $\langle v_1,\ldots,v_p,w_1,\ldots,w_q\rangle$ es el menor
        subespacio que contiene a estos vectores, se tiene
        $\langle v_1,\ldots,v_p,w_1,\ldots,w_q\rangle\subseteq W_1+W_2$.

        ($\subseteq$) Todo elemento de $W_1+W_2$ es de la forma $u+w$ con
        $u\in W_1$ y $w\in W_2$. Como $v_1,\ldots,v_p$ generan $W_1$ y
        $w_1,\ldots,w_q$ generan $W_2$, existen escalares
        $\lambda_1,\ldots,\lambda_p,\mu_1,\ldots,\mu_q\in k$ tales que
        \[
            u=\lambda_1v_1+\cdots+\lambda_pv_p,
            \qquad
            w=\mu_1w_1+\cdots+\mu_qw_q.
        \]
        Sumando,
        \[
            u+w=\lambda_1v_1+\cdots+\lambda_pv_p+\mu_1w_1+\cdots+\mu_qw_q,
        \]
        que es una combinación lineal de
        $v_1,\ldots,v_p,w_1,\ldots,w_q$.
        Por tanto $W_1+W_2\subseteq\langle v_1,\ldots,v_p,w_1,\ldots,w_q\rangle$.
    \end{proof}
\end{frame}


\begin{frame}
    \frametitle{La unión de bases no es una base de la suma}

    \begin{example}
        En $\R^3$, sean
        \[
            W_1=\langle(1,0,0),(0,1,0)\rangle,
            \qquad
            W_2=\langle(1,1,0),(0,0,1)\rangle.
        \]
        Los dos pares de generadores son linealmente independientes, luego
        son bases de $W_1$ y $W_2$. Su unión genera $W_1+W_2$, pero
        \alert{no} es una base de $W_1+W_2$.
    \end{example}

    Por la proposición anterior,
    \[
        W_1+W_2=\langle(1,0,0),(0,1,0),(1,1,0),(0,0,1)\rangle=\R^3,
    \]
    pues entre estos vectores están $(1,0,0)$, $(0,1,0)$ y $(0,0,1)$,
    que generan $\R^3$.
\end{frame}


\begin{frame}
    \frametitle{La unión de bases no es una base de la suma}

    Sin embargo, los cuatro vectores
    \[
        (1,0,0),\quad(0,1,0),\quad(1,1,0),\quad(0,0,1)
    \]
    son linealmente dependientes:
    \[
        1\cdot(1,0,0)+1\cdot(0,1,0)+(-1)\cdot(1,1,0)+0\cdot(0,0,1)=0,
    \]
    con coeficientes no todos nulos.

    \medskip

    El problema es que $(1,1,0)$ está en $W_1\cap W_2$: es a la vez
    combinación lineal de la base de $W_1$ y elemento de la base de $W_2$.
    La unión de bases genera la suma, pero en general
    \alert{no es independiente}. Para obtener una base de $W_1+W_2$ hay que
    eliminar vectores sobrantes; aquí basta quitar $(1,1,0)$.
\end{frame}


\begin{frame}
    \frametitle{Ejemplo: intersección y suma en $\R^3$}

    \begin{example}
        En $\R^3$, sean $W_1\colon u_3=0$ y $W_2\colon u_1=0$, es decir,
        $W_1=\langle e_1,e_2\rangle$ y $W_2=\langle e_2,e_3\rangle$. Entonces
        $W_1\cap W_2\colon u_1=u_3=0$, es decir,
        \[
            W_1\cap W_2=\langle e_2\rangle,
            \qquad
            W_1+W_2=\R^3.
        \]
    \end{example}

    Para $W_1+W_2$: todo vector $(u_1,u_2,u_3)\in\R^3$ se descompone como
    \[
        (u_1,u_2,u_3)=\underbrace{(u_1,u_2,0)}_{\in W_1}
        +\underbrace{(0,0,u_3)}_{\in W_2}.
    \]
\end{frame}


\begin{frame}
    \frametitle{Recordatorio: principio de inclusión-exclusión}

    Si $A$ y $B$ son conjuntos finitos, entonces
    \[
        |A\cup B|=|A|+|B|-|A\cap B|.
    \]
    Al sumar $|A|+|B|$, los elementos de $A\cap B$ se cuentan
    \alert{dos veces}; restando $|A\cap B|$ se cuentan una sola vez.

    \begin{example}
        \begin{columns}[c]
            \begin{column}{0.5\textwidth}
                \centering
                \begin{tikzpicture}[scale=0.9]
                    \fill[alertcolor!15] (0,0) circle (1.4);
                    \fill[secondarycolor!30] (1.6,0) circle (1.4);
                    \begin{scope}
                        \clip (0,0) circle (1.4);
                        \fill[alertcolor!40] (1.6,0) circle (1.4);
                    \end{scope}
                    \draw (0,0) circle (1.4);
                    \draw (1.6,0) circle (1.4);
                    \node at (-0.9,1.45) {$A$};
                    \node at (2.5,1.45) {$B$};
                    \node at (-0.55,0.45) {$1$};
                    \node at (-0.85,-0.2) {$2$};
                    \node at (-0.4,-0.7) {$3$};
                    \node at (0.8,0.35) {$4$};
                    \node at (0.8,-0.35) {$5$};
                    \node at (2.1,0.45) {$6$};
                    \node at (2.3,-0.4) {$7$};
                \end{tikzpicture}
            \end{column}
            \begin{column}{0.5\textwidth}
                $A=\{1,2,3,4,5\}$,\quad $B=\{4,5,6,7\}$,
                \[
                    A\cap B=\{4,5\}.
                \]
                \[
                    |A\cup B|=5+4-2=7.
                \]
            \end{column}
        \end{columns}
    \end{example}

    La fórmula de Grassmann es el análogo para subespacios:
    $\cup$ se sustituye por $+$ y el cardinal por la dimensión.
\end{frame}


\begin{frame}[allowframebreaks, t]
    \frametitle{Fórmula de Grassmann}

    \begin{theorem}
        Si $W_1$ y $W_2$ son subespacios de dimensión finita de $V$, entonces
        \[
            \dim(W_1+W_2)=\dim W_1+\dim W_2-\dim(W_1\cap W_2).
        \]
    \end{theorem}

    \begin{proof}
        Sea $\{u_1,\ldots,u_d\}$ una base de $W_1\cap W_2$, que
        extendemos a bases
        \[
            \{u_1,\ldots,u_d,v_1,\ldots,v_p\}\quad\text{de }W_1,
            \qquad
            \{u_1,\ldots,u_d,w_1,\ldots,w_q\}\quad\text{de }W_2.
        \]
        Afirmamos que $B=\{u_1,\ldots,u_d,v_1,\ldots,v_p,w_1,\ldots,w_q\}$
        es una base de $W_1+W_2$.

        \alert{Generación.}
        $B$ es la unión de una base de $W_1$ y una base de $W_2$,
        luego genera $W_1+W_2$ por la proposición sobre generadores
        de la suma.

        \alert{Independencia lineal.}
        Supongamos $\sum_{i=1}^{d} a_i u_i+\sum_{j=1}^{p} b_j v_j+\sum_{k=1}^{q} c_k w_k=0$.
        Entonces
        \[
            W_2\ni\sum_{k=1}^{q} c_k w_k=-\sum_{i=1}^{d} a_i u_i-\sum_{j=1}^{p} b_j v_j\in W_1,
        \]
        así que se tiene \[\sum_{k=1}^{q} c_k w_k\in W_1\cap W_2,\]
        luego \[\sum_{k=1}^{q} c_k w_k=\sum_{i=1}^{d}\lambda_i u_i\] para ciertos $\lambda_i$.

        \framebreak

        Reordenando, \[-\sum_{i=1}^{d}\lambda_i u_i+\sum_{k=1}^{q} c_k w_k=0.\]
        Como $\{u_1,\ldots,u_d,w_1,\ldots,w_q\}$ es base de $W_2$,
        concluimos que $\lambda_i, c_k=0$ para todo $i,k$.
        Entonces \[\sum_{i=1}^{d} a_i u_i+\sum_{j=1}^{p} b_j v_j=0,\] y como
        $\{u_1,\ldots,u_d,v_1,\ldots,v_p\}$ es base de $W_1$,
        también $a_i=b_j=0$.

        \medskip

        Por tanto $\dim(W_1+W_2)=d+p+q=\dim W_1+\dim W_2-d$.
    \end{proof}
\end{frame}


\begin{frame}
    \frametitle{Consecuencias de la fórmula de Grassmann}

    Para $W_1=\langle e_1,e_2\rangle$ y $W_2=\langle e_2,e_3\rangle$ en $\R^3$:
    \[
        \dim(W_1+W_2)=2+2-1=3,
    \]
    lo que confirma $W_1+W_2=\R^3$.

    \begin{corollary}
        Si $\dim V=n$ y $W_1,W_2\subseteq V$, entonces
        \[
            \dim(W_1\cap W_2)\geq\dim W_1+\dim W_2-n.
        \]
        En particular, si $\dim W_1+\dim W_2>n$, la intersección
        no puede ser trivial.
    \end{corollary}
\end{frame}


\begin{frame}[allowframebreaks, t]
    \frametitle{Suma directa}

    \begin{definition}
        Se dice que $V$ es la \alert{suma directa} de $W_1$ y $W_2$,
        y se escribe $V=W_1\oplus W_2$, si
        $V=W_1+W_2$ y $W_1\cap W_2=\{0\}$.
    \end{definition}

    \begin{proposition}
        $V=W_1\oplus W_2$ si y solo si todo $v\in V$ se escribe de
        manera \alert{única} como $v=u+w$ con $u\in W_1$ y $w\in W_2$.
    \end{proposition}

    \begin{proof}
        ($\Rightarrow$) $V=W_1+W_2$ garantiza la existencia. Si
        $u_1+w_1=u_2+w_2$ con $u_i\in W_1$ y $w_i\in W_2$, entonces
        $u_1-u_2=w_2-w_1\in W_1\cap W_2=\{0\}$, luego $u_1=u_2$ y $w_1=w_2$.

        ($\Leftarrow$) La existencia da $V=W_1+W_2$. Si $v\in W_1\cap W_2$,
        entonces $v=v+0=0+v$ son dos descomposiciones de $v$; por
        unicidad, $v=0$.
    \end{proof}

    \begin{corollary}
        Si $V=W_1\oplus W_2$, entonces $\dim V=\dim W_1+\dim W_2$.
    \end{corollary}
\end{frame}


\begin{frame}
    \frametitle{Ejemplos de suma directa}

    \begin{example}
        En $\R^3$:
        \begin{enumerate}
            \item $\R^3=\langle e_1,e_2\rangle\oplus\langle e_3\rangle$:
            todo $(u_1,u_2,u_3)=(u_1,u_2,0)+(0,0,u_3)$, de forma única.
            \item $\R^3=\langle e_1\rangle\oplus\langle e_2\rangle\oplus\langle e_3\rangle$.
            \item $\langle e_1,e_2\rangle$ y $\langle e_2,e_3\rangle$
            no dan suma directa, pues su intersección es
            $\langle e_2\rangle\neq\{0\}$.
        \end{enumerate}
    \end{example}

    \medskip

    En general, si $\{v_1,\ldots,v_n\}$ es una base de $V$, entonces
    \[
        V=\langle v_1\rangle\oplus\cdots\oplus\langle v_n\rangle.
    \]
\end{frame}

% % Espacio cociente

\begin{frame}
    \frametitle{Identificar vectores que difieren en un subespacio}

    \begin{columns}[c]
        \begin{column}{0.55\textwidth}
            \begin{example}
                En $\R^2$, sea $W=\langle(1,1)\rangle$. Decidimos
                \alert{no distinguir} entre $u$ y $u'$ si $u-u'\in W$.
            \end{example}

            Los vectores que no distinguimos de $u$ son
            \[
                u+W=\{u+\lambda(1,1):\ \lambda\in\R\},
            \]
            la recta que pasa por $u$ y es \alert{paralela} a $W$.

            \medskip

            Cada una de estas rectas se convierte en un único objeto. ¿Se
            pueden sumar y multiplicar por escalares?
        \end{column}
        \begin{column}{0.42\textwidth}
            \centering
            \begin{tikzpicture}[scale=0.7]
                \clip (-2.6,-2.6) rectangle (2.8,2.8);
                \draw[->] (-2.5,0) -- (2.5,0) node[below left] {\footnotesize $u_1$};
                \draw[->] (0,-2.5) -- (0,2.5) node[below left] {\footnotesize $u_2$};
                \foreach \c in {-2,-1,1,2}
                    \draw[gray!60] (-3,-3+\c) -- (3,3+\c);
                \draw[very thick,alertcolor] (-3,-3) -- (3,3);
                \node[alertcolor] at (2.1,1.5) {\footnotesize $W$};
                \draw[very thick,secondarycolor!80!black] (-3,-1.5) -- (3,4.5);
                \draw[->,thick] (0,0) -- (0.5,2);
                \fill (0.5,2) circle (1.5pt);
                \node[right] at (0.5,2) {\footnotesize $u$};
                \node[secondarycolor!80!black] at (-1.2,1.7) {\footnotesize $u+W$};
            \end{tikzpicture}
        \end{column}
    \end{columns}
\end{frame}


\begin{frame}
    \frametitle{Clases módulo un subespacio}

    \begin{definition}
        Sea $W$ un subespacio de $V$ y sea $u\in V$. La \alert{clase} de $u$
        módulo $W$ es
        \[
            u+W=\{u+w:\ w\in W\}\subseteq V.
        \]
        Se dice que $u$ es un \alert{representante} de la clase $u+W$.
    \end{definition}

    \begin{proposition}
        Para $u,u'\in V$, son equivalentes:
        \begin{enumerate}[(1)]
            \item $u+W=u'+W$;
            \item $u'\in u+W$;
            \item $u-u'\in W$.
        \end{enumerate}
    \end{proposition}

    Una misma clase puede tener, por tanto, muchos representantes: todos los
    vectores que difieren de $u$ en un elemento de $W$.
\end{frame}


\begin{frame}[t]
    \frametitle{Clases módulo un subespacio: demostración}

    \begin{proof}
        \alert{(1)$\Rightarrow$(2).} $u'=u'+0\in u'+W=u+W$.

        \alert{(2)$\Rightarrow$(3).} Si $u'=u+w$ con $w\in W$, entonces
        $u-u'=-w\in W$.

        \alert{(3)$\Rightarrow$(1).} Sea $w_0=u-u'\in W$. Todo elemento de
        $u+W$ es de la forma
        \[
            u+w=u'+(w_0+w)\in u'+W,
        \]
        pues $w_0+w\in W$. Luego $u+W\subseteq u'+W$. Como también
        $u'-u=-w_0\in W$, el mismo argumento da $u'+W\subseteq u+W$.
    \end{proof}
\end{frame}


\begin{frame}[allowframebreaks, t]
    \frametitle{Las clases son iguales o disjuntas}

    \begin{corollary}
        Sean $u,u'\in V$.
        \begin{enumerate}[(1)]
            \item Las clases $u+W$ y $u'+W$ son \alert{iguales o disjuntas}.
            \item $u+W=W$ si y solo si $u\in W$.
        \end{enumerate}
    \end{corollary}

    \begin{proof}
        \alert{(1)} Supongamos que $u+W$ y $u'+W$ no son disjuntas y sea
        $u''\in(u+W)\cap(u'+W)$. Aplicando la implicación
        (2)$\Rightarrow$(1) de la proposición anterior a $u''\in u+W$ y a
        $u''\in u'+W$, obtenemos
        \[
            u+W=u''+W
            \qquad\text{y}\qquad
            u'+W=u''+W.
        \]
        Por tanto, $u+W=u'+W$.

        \medskip

        \alert{(2)} Se tiene $0+W=\{0+w:\ w\in W\}=W$. Por la equivalencia
        (1)$\Leftrightarrow$(3) de la proposición anterior,
        \[
            u+W=0+W\iff u-0\in W\iff u\in W.\qedhere
        \]
    \end{proof}
\end{frame}


\begin{frame}
    \frametitle{Operaciones con clases}

    \begin{definition}
        El \alert{cociente} de $V$ por $W$ es el conjunto de las clases
        \[
            V/W=\{u+W:\ u\in V\},
        \]
        con las operaciones
        \[
            (u_1+W)+(u_2+W)=(u_1+u_2)+W,
            \qquad
            \lambda(u_1+W)=\lambda u_1+W.
        \]
    \end{definition}

    Las operaciones se definen eligiendo representantes. Hay que comprobar
    que el resultado \alert{no depende} de la elección.
\end{frame}


\begin{frame}
    \frametitle{Las operaciones no dependen de los representantes}

    \begin{lemma}
        Si $u_1+W=u_1'+W$ y $u_2+W=u_2'+W$, entonces, dado cualquier $\lambda\in k$,
        \[
            (u_1+u_2)+W=(u_1'+u_2')+W
            \quad\text{y}\quad
            \lambda u_1+W=\lambda u_1'+W.
        \]
    \end{lemma}

    \begin{proof}
        Por hipótesis, $u_1-u_1',\,u_2-u_2'\in W$. Como $W$ es un subespacio,
        $(u_1+u_2)-(u_1'+u_2')=(u_1-u_1')+(u_2-u_2')\in W$ y
        $\lambda u_1-\lambda u_1'=\lambda(u_1-u_1')\in W$.
    \end{proof}
\end{frame}


\begin{frame}
    \frametitle{El espacio vectorial cociente}

    \begin{theorem}
        Con estas operaciones, $V/W$ es un espacio vectorial, llamado
        \alert{espacio cociente}. Su vector nulo es $0+W=W$ y el opuesto de
        $u+W$ es $(-u)+W$.
    \end{theorem}

    \begin{proof}
        Cada propiedad se hereda de la correspondiente de $V$, operando con
        representantes. Por ejemplo, la asociatividad:
        \begin{align*}
            \bigl((u_1+W)+(u_2+W)\bigr)+(u_3+W)
            &=\bigl((u_1+u_2)+u_3\bigr)+W\\
            &=\bigl(u_1+(u_2+u_3)\bigr)+W\\
            &=(u_1+W)+\bigl((u_2+W)+(u_3+W)\bigr).
        \end{align*}
        Del mismo modo, $(u+W)+(0+W)=u+W$ y $(u+W)+((-u)+W)=0+W$. Las
        demás propiedades se comprueban igual.
    \end{proof}

    \medskip

    Por el apartado (2) del corolario sobre las clases, en $V/W$ se tiene $u+W=0$ si y solo si
    $u\in W$: \alert{los vectores de $W$ se vuelven nulos} en el cociente.
\end{frame}


\begin{frame}
    \frametitle{Ejemplos de espacios cociente}

    \begin{example}
        \begin{enumerate}
            \item Si $W=\langle(1,1)\rangle\subseteq\R^2$, entonces
            $(u_1,u_2)+W=(u_1',u_2')+W$
            si y solo si $(u_1-u_1',u_2-u_2')\in\langle(1,1)\rangle$, es
            decir, $u_1-u_2=u_1'-u_2'$. La clase de $(u_1,u_2)$ queda
            determinada por el número $u_1-u_2$.
            \item Si $W=\{0\}$, entonces $u+W=\{u\}$: el cociente $V/\{0\}$
            es, en esencia, el propio $V$. 
            \item Si $W=V$, todas las clases son
            iguales a $V$ y $V/V=\{0\}$.
            \item En $k[x]$, sea $W=\{p(x):\ p(0)=0\}$. Entonces
            $p+W=q+W$ si y solo si $p(0)=q(0)$, y cada clase contiene una
            única constante: $p+W=p(0)+W$.
        \end{enumerate}
    \end{example}

    En el último ejemplo, $k[x]$ tiene dimensión infinita, pero las clases
    quedan determinadas por un solo número. Veremos que
    $\dim k[x]/W=1$.
\end{frame}


\begin{frame}[allowframebreaks, t]
    \frametitle{Dimensión y base del cociente}

    \begin{theorem}
        Sea $V$ de dimensión $n$ y sea $W\subseteq V$ un subespacio de
        dimensión $m$. Si $B=\{w_1,\ldots,w_m,w_{m+1},\ldots,w_n\}$ es una base de $V$ adaptada
        a $W$, es decir, $\{w_1,\ldots,w_m\}$ es una base de $W$, entonces
        \[
            \overline{B}=\{w_{m+1}+W,\ldots,w_n+W\}
        \]
        es una base de $V/W$. En particular,
        \[
            \dim V/W=\dim V-\dim W.
        \]
    \end{theorem}

    \framebreak

    \begin{proof}
        \alert{$\overline{B}$ genera $V/W$.} Sea $u\in V$ y escribamos
        $u=\lambda_1w_1+\cdots+\lambda_nw_n$. Por la definición de las
        operaciones en $V/W$,
        \[
            u+W=\lambda_1(w_1+W)+\cdots+\lambda_n(w_n+W).
        \]
        Para $i\leq m$ se tiene $w_i\in W$, luego $w_i+W=0$. Así,
        \[
            u+W=\lambda_{m+1}(w_{m+1}+W)+\cdots+\lambda_n(w_n+W).
        \]

        \framebreak

        \alert{$\overline{B}$ es linealmente independiente.} Si
        $\lambda_{m+1}(w_{m+1}+W)+\cdots+\lambda_n(w_n+W)=0$, entonces
        $(\lambda_{m+1}w_{m+1}+\cdots+\lambda_nw_n)+W=0$, es decir,
        $\lambda_{m+1}w_{m+1}+\cdots+\lambda_nw_n\in W$. Existen entonces
        $\mu_1,\ldots,\mu_m\in k$ con
        \[
            \lambda_{m+1}w_{m+1}+\cdots+\lambda_nw_n=\mu_1w_1+\cdots+\mu_mw_m.
        \]
        Pasando todo a un miembro obtenemos una combinación lineal nula de
        los vectores de $B$, que es una base; 
        \[
            -\mu_1w_1-\cdots-\mu_mw_m+\lambda_{m+1}w_{m+1}+\cdots+\lambda_nw_n=0,
        \]
        luego todos los coeficientes
        son nulos y, en particular, $\lambda_{m+1}=\cdots=\lambda_n=0$.

        \medskip

        Por tanto, $\overline{B}$ es una base de $V/W$ con $n-m$ elementos.
    \end{proof}
\end{frame}


\begin{frame}
    \frametitle{Coordenadas en el cociente}

    \begin{corollary}
        En las condiciones del teorema, si
        $u_B=(\lambda_1,\ldots,\lambda_m,\lambda_{m+1},\ldots,\lambda_n)$, entonces
        \[
            (u+W)_{\overline{B}}=(\lambda_{m+1},\ldots,\lambda_n).
        \]
    \end{corollary}

    \begin{proof}
        Basta observar que en la demostración del teorema hemos visto que
        \[u+W=\lambda_{m+1}(w_{m+1}+W)+\cdots+\lambda_n(w_n+W).\]
    \end{proof}
\end{frame}


\begin{frame}
    \frametitle{Ejemplo: una base de un cociente de $\R^4$}

    \begin{example}
        Sea $W=\langle w_1,w_2\rangle\subseteq\R^4$ con
        $w_1=(1,0,1,0)$ y $w_2=(0,1,0,1)$. Calculemos una base de $\R^4/W$.
    \end{example}

    Como $w_1,w_2$ no son proporcionales, $\dim W=2$ y
    $\dim\R^4/W=4-2=2$. Ampliamos $\{w_1,w_2\}$ con la base canónica,
    escalonando la matriz de columnas $w_1,w_2,e_1,e_2,e_3,e_4$:
    \[
        \begin{pmatrix}
            1&0&1&0&0&0\\
            0&1&0&1&0&0\\
            1&0&0&0&1&0\\
            0&1&0&0&0&1
        \end{pmatrix}
        \xrightarrow[\;F_4-F_2\;]{\;F_3-F_1\;}
        \begin{pmatrix}
            \alert{1}&0&1&0&0&0\\
            0&\alert{1}&0&1&0&0\\
            0&0&\alert{-1}&0&1&0\\
            0&0&0&\alert{-1}&0&1
        \end{pmatrix}.
    \]
    Los pivotes están en las columnas $1$ a $4$, así que
    $B=\{w_1,w_2,e_1,e_2\}$ es una base de $\R^4$ adaptada a $W$ y
    \[
        \overline{B}=\{e_1+W,\ e_2+W\}
    \]
    es una base de $\R^4/W$.
\end{frame}


\begin{frame}
    \frametitle{Ejemplo: coordenadas en un cociente de $\R^4$}

    Calculemos las coordenadas de $u+W$ en $\overline{B}$ para
    $u=(3,1,2,5)$. Para ello, hallamos las coordenadas $(\lambda_1,\lambda_2,\lambda_3,\lambda_4)$ de $u$ en $B$, que son las soluciones del sistema con matriz ampliada:
    \[
        \left(\begin{array}{cccc|c}
            1&0&1&0&3\\
            0&1&0&1&1\\
            1&0&0&0&2\\
            0&1&0&0&5
        \end{array}\right)
        \xrightarrow[\;F_4-F_2\;]{\;F_3-F_1\;}
        \left(\begin{array}{cccc|c}
            1&0&1&0&3\\
            0&1&0&1&1\\
            0&0&-1&0&-1\\
            0&0&0&-1&4
        \end{array}\right).
    \]
    Solo necesitamos las dos últimas incógnitas: $\lambda_3=1$ y
    $\lambda_4=-4$. Por tanto,
    \[
        (u+W)_{\overline{B}}=(1,-4),
        \qquad\text{es decir,}\qquad
        u+W=(e_1+W)-4(e_2+W).
    \]
    En general, el mismo cálculo da
    $\bigl((u_1,u_2,u_3,u_4)+W\bigr)_{\overline{B}}=(u_1-u_3,\ u_2-u_4)$,
    y $W\colon u_1-u_3=u_2-u_4=0$.
\end{frame}


\begin{frame}
    \frametitle{Ejemplo: un cociente de polinomios}

    \begin{example}
        Sea $W=\{p(x)\in\R_2[x]:\ p(-1)=0\}$. Describamos $\R_2[x]/W$.
    \end{example}

    Ya vimos que $B=\{1+x,\ x+x^2,\ 1\}$ es una base de $\R_2[x]$ adaptada
    a $W$. Por el teorema,
    \[
        \dim\R_2[x]/W=3-2=1,
        \qquad
        \overline{B}=\{1+W\}.
    \]
    Supongamos que las coordenadas de un polinomio $p(x)\in \R_2[x]$ en $B$ son $(a,b,c)$. Es decir, $p(x)=a(1+x)+b(x+x^2)+c$. Entonces $(p+W)_{\overline{B}}=c$. Además observamos que $p(-1)=c$, así que
    \[
        (p+W)_{\overline{B}}=p(-1).
    \]

    \medskip

    Así, $p+W=q+W$ si y solo si $p(-1)=q(-1)$: en el cociente solo
    \alert{recordamos el valor en $-1$}, que es lo que $W$ no ve.
\end{frame}


\begin{frame}
    \frametitle{El cociente $V/W$ como conjunto cociente}

    \begin{proposition}
        Sea $W$ un subespacio de $V$. La relación en $V$ dada por

        \vspace{-10pt}

        \[
            u\sim u'\iff u-u'\in W
        \]

        \vspace{-2pt}

        es de equivalencia, y la clase de equivalencia de $u$ es $[u]=u+W$.
        Por tanto, $V/W$ es el conjunto cociente $V/{\sim}$.
    \end{proposition}

    \vspace{-7pt}

    \begin{proof}
        Cada propiedad se sigue de que $W$ es un subespacio:
        \begin{itemize}
            \item reflexiva: $u-u=0\in W$;
            \item simétrica: si $u-u'\in W$, entonces $u'-u=-(u-u')\in W$;
            \item transitiva: si $u-u'\in W$ y $u'-u''\in W$, entonces
            $u-u''=(u-u')+(u'-u'')\in W$.
        \end{itemize}
        Por último, $u'\in[u]$ si y solo si $u'-u\in W$, es decir,
        $u'=u+w$ con $w\in W$, o sea, $u'\in u+W$.
    \end{proof}
\end{frame}
  % pendiente
% \input{temas/03-subespacios/ejercicios/...}
\section{Homomorfismos de espacios vectoriales}

% \input{temas/04-homomorfismos/teoria/...}
% \input{temas/04-homomorfismos/ejercicios/...}

\section{Forma canónica}

% \input{temas/05-forma-canonica/teoria/...}
% \input{temas/05-forma-canonica/ejercicios/...}

\section{Espacio vectorial euclídeo}

% \input{temas/06-espacio-vectorial-euclideo/teoria/...}
% \input{temas/06-espacio-vectorial-euclideo/ejercicios/...}

\section{Teoremas espectrales}

% \input{temas/07-teoremas-espectrales/teoria/...}
% \input{temas/07-teoremas-espectrales/ejercicios/...}

\section{Espacios afines}

% \input{temas/08-espacios-afines/teoria/...}
% \input{temas/08-espacios-afines/ejercicios/...}

\section{Aplicaciones afines}

% \input{temas/09-aplicaciones-afines/teoria/...}
% \input{temas/09-aplicaciones-afines/ejercicios/...}

\section{Espacios afines euclídeos}

% \input{temas/10-espacios-afines-euclideos/teoria/...}
% \input{temas/10-espacios-afines-euclideos/ejercicios/...}

\section{Movimientos}

% \input{temas/11-movimientos/teoria/...}
% \input{temas/11-movimientos/ejercicios/...}


\end{document}
